% Generated mechanically from frozen input p07/ko/derham.tex.
% Do not edit this generated file; edit the bound locale source instead.

% OK, start here
%

\setcounter{chapter}{49}
\chapter{드람 코호몰로지}



\phantomsection
\label{derham-section-phantom}


\section{서론}
\label{derham-section-introduction}

\noindent
이 장에서는 스킴의 사상에 대한 드람 복합체를 논의하는 데서 시작하여,
기초체의 표수가 영일 때 드람 코호몰로지가 베유 코호몰로지 이론을
이룬다는 것을 증명하는 것으로 끝맺는다.




\section{드람 복합체}
\label{derham-section-de-rham-complex}

\noindent
스킴의 사상 $p : X \to S$가 주어졌다고 하자. 이때
$$
\Omega^\bullet_{X/S} =
\mathcal{O}_{X/S} \to \Omega^1_{X/S} \to \Omega^2_{X/S} \to \ldots
$$
라는 $p^{-1}\mathcal{O}_S$-가군의 복합체가 있다. 여기서
$\Omega^i_{X/S} = \wedge^i(\Omega_{X/S})$는 차수 $i$에 놓이며,
미분은 국소 절단에 대한 규칙
$\text{d}(g_0 \text{d}g_1 \wedge \ldots \wedge \text{d}g_p) =
\text{d}g_0 \wedge \text{d}g_1 \wedge \ldots \wedge \text{d}g_p$
으로 정해진다. 가군층의 절 \ref{modules-section-de-rham-complex}을 보라.

\medskip\noindent
스킴의 가환 도식
$$
\xymatrix{
X' \ar[r]_f \ar[d] & X \ar[d] \\
S' \ar[r] & S
}
$$
이 주어지면 복합체의 표준 사상
$f^{-1}\Omega_{X/S}^\bullet \to \Omega^\bullet_{X'/S'}$와
$\Omega_{X/S}^\bullet \to f_*\Omega^\bullet_{X'/S'}$가 있다.
가군층의 절 \ref{modules-section-de-rham-complex}을 보라.
이를 선형화하면 모든 $p$에 대하여 선형 사상
$f^*\Omega^p_{X/S} \to \Omega^p_{X'/S'}$
을 얻는다.

\medskip\noindent
특히 $f : Y \to X$가 기초 스킴 $S$ 위 스킴의 사상이면 복합체의 사상
$$
\Omega^\bullet_{X/S} \longrightarrow f_*\Omega^\bullet_{Y/S}
$$
이 있다. 이를 선형화하면 모든 $p \geq 0$에 대하여 표준 사상
$$
\Omega^p_{X/S} \otimes_{\mathcal{O}_X} f_*\mathcal{O}_Y
\longrightarrow
f_*\Omega^p_{Y/S}
$$
을 얻는다.

\begin{lemma}
\label{derham-lemma-base-change-de-rham}
스킴의 데카르트 도식
$$
\xymatrix{
X' \ar[r]_f \ar[d] & X \ar[d] \\
S' \ar[r] & S
}
$$
이 주어졌다고 하자. 그러면 위에서 논한 사상들은 동형사상
$f^*\Omega^p_{X/S} \to \Omega^p_{X'/S'}$
을 유도한다.
\end{lemma}

\begin{proof}
스킴의 사상의 보조정리 \ref{morphisms-lemma-base-change-differentials}와
외대수의 각 차수 성분을 취하는 연산이 기저변환과 가환한다는 사실을
함께 적용한다.
\end{proof}

\begin{lemma}
\label{derham-lemma-etale}
스킴의 가환 도식
$$
\xymatrix{
X' \ar[r]_f \ar[d] & X \ar[d] \\
S' \ar[r] & S
}
$$
을 생각하자. $X' \to X$와 $S' \to S$가 에탈이면 위에서 논한
사상들은 동형사상
$f^*\Omega^p_{X/S} \to \Omega^p_{X'/S'}$
을 유도한다.
\end{lemma}

\begin{proof}
예컨대 스킴의 사상의 보조정리 \ref{morphisms-lemma-etale-at-point}에 의해
$\Omega_{S'/S} = 0$이고 $\Omega_{X'/X} = 0$이다. 따라서 스킴의 사상의
보조정리
\ref{morphisms-lemma-triangle-differentials}와
\ref{morphisms-lemma-triangle-differentials-smooth}
에 나오는 짧은 완전열들에 의해
$\Omega_{X'/S'} = \Omega_{X'/S} = f^*\Omega_{X/S}$임을 알 수 있다.
외대수의 각 차수 성분을 취하면 결론을 얻는다.
\end{proof}






\section{드람 코호몰로지}
\label{derham-section-de-rham-cohomology}

\noindent
스킴의 사상 $p : X \to S$가 주어졌다고 하자. {\it $X$의
$S$ 위 드람 코호몰로지}를 코호몰로지 군
$$
H^i_{dR}(X/S) = H^i(R\Gamma(X, \Omega^\bullet_{X/S}))
$$
으로 정의한다. $\Omega^\bullet_{X/S}$가 $p^{-1}\mathcal{O}_S$-가군의
복합체이므로 이 코호몰로지 군들은 자연스럽게
$H^0(S, \mathcal{O}_S)$-가군이다.

\medskip\noindent
스킴의 가환 도식
$$
\xymatrix{
X' \ar[r]_f \ar[d] & X \ar[d] \\
S' \ar[r] & S
}
$$
이 주어지면 절 \ref{derham-section-de-rham-complex}의 표준 사상을 사용하여
당김 사상
$$
f^* :
R\Gamma(X, \Omega^\bullet_{X/S})
\longrightarrow
R\Gamma(X', \Omega^\bullet_{X'/S'})
$$
과
$$
f^* : H^i_{dR}(X/S) \longrightarrow H^i_{dR}(X'/S')
$$
을 얻는다. 이 당김들은 자명한 합성 법칙을 만족한다.
특히 기초 스킴 $S$를 고정하여 작업하면 드람 코호몰로지는
$S$ 위 스킴의 범주에서 정의된 반변함자이다.

\begin{lemma}
\label{derham-lemma-de-rham-affine}
아핀 스킴의 사상 $X \to S$가 환 사상 $R \to A$로 주어진다고 하자.
그러면 $R\Gamma(X, \Omega^\bullet_{X/S}) = \Omega^\bullet_{A/R}$는
$D(R)$에서 성립하고
$H^i_{dR}(X/S) = H^i(\Omega^\bullet_{A/R})$이다.
\end{lemma}

\begin{proof}
이는 스킴의 코호몰로지의 보조정리
\ref{coherent-lemma-quasi-coherent-affine-cohomology-zero}
와 Leray 비순환성 보조정리(유도 범주의 보조정리
\ref{derived-lemma-leray-acyclicity})에서 따른다.
\end{proof}

\begin{lemma}
\label{derham-lemma-quasi-coherence-relative}
스킴의 사상 $p : X \to S$가 주어졌다고 하자. $p$가 준콤팩트이고
준분리이면
$Rp_*\Omega^\bullet_{X/S}$는 $D_\QCoh(\mathcal{O}_S)$의 대상이다.
\end{lemma}

\begin{proof}
첫째 쪽이 $E_1^{a, b} = R^bp_*\Omega^a_{X/S}$이고
$Rp_*\Omega^\bullet_{X/S}$의 코호몰로지로 수렴하는 스펙트럼 수열이 있다
(유도 범주의 보조정리
\ref{derived-lemma-two-ss-complex-functor}를 보라).
따라서 호몰로지 대수의 보조정리 \ref{homology-lemma-first-quadrant-ss}에 의해
$R^bp_*\Omega^a_{X/S}$가 준연접임을 보이면 충분하다. 이는 스킴의 코호몰로지의 보조정리
\ref{coherent-lemma-quasi-coherence-higher-direct-images}에서 따른다.
\end{proof}

\begin{lemma}
\label{derham-lemma-coherence-relative}
$p : X \to S$를 스킴의 고유 사상이라 하고 $S$가 국소 뇌터라고 하자.
그러면 $Rp_*\Omega^\bullet_{X/S}$는
$D_{\textit{Coh}}(\mathcal{O}_S)$의 대상이다.
\end{lemma}

\begin{proof}
이 경우 스킴의 사상의 보조정리
\ref{morphisms-lemma-finite-type-differentials}에 의해 가군
$\Omega^i_{X/S}$는 연접이다. 따라서 보조정리
\ref{derham-lemma-quasi-coherence-relative}의 증명과 똑같은 논증을 적용하되,
스킴의 코호몰로지의 명제
\ref{coherent-proposition-proper-pushforward-coherent}를 사용할 수 있다.
\end{proof}

\begin{lemma}
\label{derham-lemma-finite-de-Rham}
$A$를 뇌터 환이라 하고 $X$를 $S = \Spec(A)$ 위의 고유 스킴이라 하자.
그러면 $H^i_{dR}(X/S)$는 유한 $A$-가군이며, 이는 모든 $i$에 대하여
성립한다.
\end{lemma}

\begin{proof}
이는 보조정리 \ref{derham-lemma-coherence-relative}의 특수한 경우이다.
\end{proof}

\begin{lemma}
\label{derham-lemma-proper-smooth-de-Rham}
$f : X \to S$를 스킴의 고유하고 매끄러운 사상이라 하자. 그러면
$Rf_*\Omega^p_{X/S}$ ($p \geq 0$)와
$Rf_*\Omega^\bullet_{X/S}$는 $D(\mathcal{O}_S)$의 완전 대상이며,
그 구성은 임의의 기저변환과 가환한다.
\end{lemma}

\begin{proof}
$f$가 매끄러우므로 $\Omega^p_{X/S}$는 유한 국소 자유
$\mathcal{O}_X$-가군이다. 스킴의 사상의 보조정리
\ref{morphisms-lemma-smooth-omega-finite-locally-free}를 보라.
보조정리 \ref{derham-lemma-base-change-de-rham}에 의해 그 구성은 임의의
기저변환과 가환한다. 따라서 $Rf_*\Omega^p_{X/S}$는
$D(\mathcal{O}_S)$의 완전 대상이고 그 구성은 임의의 기저변환과
가환한다. 스킴의 유도 범주의 보조정리
\ref{perfect-lemma-flat-proper-perfect-direct-image-general}를 보라.
이로써 보조정리의 첫 번째 주장을 증명했다.

\medskip\noindent
$Rf_*\Omega^\bullet_{X/S}$가 $S$ 위에서 완전임을 증명할 때에는 $S$
위에서 국소적으로 작업해도 된다. 따라서 $S$가 준콤팩트라고 가정할
수 있다. 그러면 $\Omega^n_{X/S}$가 충분히 큰 $n$에 대하여 영이라고
가정할 수 있다. 모든 $p \geq 0$에 대하여
$Rf_*\sigma_{\geq p}\Omega^\bullet_{X/S}$가 $D(\mathcal{O}_S)$의
완전 대상이고 그 구성이 임의의 기저변환과 가환한다고 주장한다.
위의 결과로부터 이는 $p \gg 0$일 때 참이다. 주장이 $p$에 대하여
성립한다고 가정하고 완전 삼각형
$$
\sigma_{\geq p}\Omega^\bullet_{X/S} \to
\sigma_{\geq p - 1}\Omega^\bullet_{X/S} \to
\Omega^{p - 1}_{X/S}[-(p - 1)] \to
(\sigma_{\geq p}\Omega^\bullet_{X/S})[1]
$$
을 $D(f^{-1}\mathcal{O}_S)$에서 생각하자. 완전함자 $Rf_*$를 적용하면
$D(\mathcal{O}_S)$에서 완전 삼각형을 얻는다. 완전성에 대해서는 셋 중
둘 성질이 성립하므로(층 코호몰로지의 보조정리
\ref{cohomology-lemma-two-out-of-three-perfect}),
$Rf_*\sigma_{\geq p - 1}\Omega^\bullet_{X/S}$가
$D(\mathcal{O}_S)$의 완전 대상임을 얻는다. 임의의 기저변환과의
가환성도 마찬가지이다.
\end{proof}





\section{컵곱}
\label{derham-section-cup-product}

\noindent
사상
$\Omega^p_{X/S} \times \Omega^q_{X/S} \to \Omega^{p + q}_{X/S}$
을 생각하자. 이는 $(\omega , \eta) \longmapsto \omega \wedge \eta$로
주어진다. 절 \ref{derham-section-de-rham-complex}에 제시된 $\text{d}$의 공식과
$\text{d} : \mathcal{O}_X \to \Omega_{X/S}$에 대한 라이프니츠 법칙을 사용하면
$\text{d}(\omega \wedge \eta) = \text{d}(\omega) \wedge \eta +
(-1)^{\deg(\omega)} \omega \wedge \text{d}(\eta)$임을 알 수 있다. 따라서
$\wedge$는 복합체의 사상
\begin{equation}
\label{derham-equation-wedge}
\wedge :
\text{Tot}(
\Omega^\bullet_{X/S} \otimes_{p^{-1}\mathcal{O}_S} \Omega^\bullet_{X/S})
\longrightarrow
\Omega^\bullet_{X/S}
\end{equation}
을 정의하며, 이 복합체들은 $p^{-1}\mathcal{O}_S$-가군의 복합체이다.

\medskip\noindent
층 코호몰로지의 절 \ref{cohomology-section-cup-product}의 컵곱과
(\ref{derham-equation-wedge})를 결합하면 $H^0(S, \mathcal{O}_S)$-쌍선형인 컵곱 사상
$$
\cup : H^i_{dR}(X/S) \times H^j_{dR}(X/S) \longrightarrow H^{i + j}_{dR}(X/S)
$$
을 얻는다. 예를 들어 $\omega \in \Gamma(X, \Omega^i_{X/S})$와
$\eta \in \Gamma(X, \Omega^j_{X/S})$가 닫혀 있으면,
$\omega$와 $\eta$의 드람 코호몰로지류의 컵곱은
$\omega \wedge \eta$의 드람 코호몰로지류이다. 층 코호몰로지의 절
\ref{cohomology-section-cup-product}의 논의를 보라.

\medskip\noindent
스킴의 가환 도식
$$
\xymatrix{
X' \ar[r]_f \ar[d] & X \ar[d] \\
S' \ar[r] & S
}
$$
이 주어지면 당김 사상
$f^* : R\Gamma(X, \Omega^\bullet_{X/S}) \to R\Gamma(X', \Omega^\bullet_{X'/S'})$
과
$f^* : H^i_{dR}(X/S) \longrightarrow H^i_{dR}(X'/S')$
은 위에서 정의한 컵곱과 양립한다.

\begin{lemma}
\label{derham-lemma-cup-product-graded-commutative}
$p : X \to S$를 스킴의 사상이라 하자.
$H^*_{dR}(X/S)$ 위의 컵곱은 결합적이고 등급 가환이다.
\end{lemma}

\begin{proof}
이는 층 코호몰로지의 보조정리
\ref{cohomology-lemma-cup-product-associative}와
\ref{cohomology-lemma-cup-product-commutative}, 그리고 $\wedge$가 결합적이고
등급 가환이라는 사실에서 따른다.
\end{proof}

\begin{remark}
\label{derham-remark-relative-cup-product}
$p : X \to S$를 스킴의 사상이라 하자. 그러면
$\Omega^\bullet_{X/S}$를 미분 등급 $p^{-1}\mathcal{O}_S$-대수의 층으로
볼 수 있다. 미분 등급 층의 정의
\ref{sdga-definition-dga}를 보라. 특히 미분 등급 층의 절
\ref{sdga-section-misc}의 논의를 적용할 수 있다. 예를 들어 이는 임의의
스킴의 가환 도식
$$
\xymatrix{
X \ar[d]_p \ar[r]_f & Y \ar[d]^q \\
S \ar[r]^h & T
}
$$
에 대하여 표준 상대 컵곱
$$
\mu :
Rf_*\Omega^\bullet_{X/S}
\otimes_{q^{-1}\mathcal{O}_T}^\mathbf{L}
Rf_*\Omega^\bullet_{X/S}
\longrightarrow
Rf_*\Omega^\bullet_{X/S}
$$
이 $D(Y, q^{-1}\mathcal{O}_T)$에 존재함을 뜻한다. 이 곱은 결합적이며,
코호몰로지 위에서는 위에서 논한 컵곱을 재현한다.
\end{remark}

\begin{remark}
\label{derham-remark-cup-product-as-a-map}
$f : X \to S$를 스킴의 사상이라 하자. $\xi \in H_{dR}^n(X/S)$라 하자.
미분 등급 층의 절 \ref{sdga-section-misc}의 논의에 따르면
표준 사상
$$
\xi' : \Omega^\bullet_{X/S} \to \Omega^\bullet_{X/S}[n]
$$
이 $D(f^{-1}\mathcal{O}_S)$에 존재하며, 다음 성질 목록의 (1), (2)에 의해
유일하게 특징지어진다.
\begin{enumerate}
\item $\xi'$는 오른쪽 미분 등급 $\Omega^\bullet_{X/S}$-가군의 유도범주에
있는 사상으로 올릴 수 있고,
\item $\xi'(1) = \xi$가
$H^0(X, \Omega^\bullet_{X/S}[n]) = H^n_{dR}(X/S)$,
에서 성립하며,
\item 사상 $\xi'$는 $\eta \in H^m_{dR}(X/S)$를
$\xi \cup \eta$로 보내며, 그 결과는 $H^{n + m}_{dR}(X/S)$에 속하고,
\item $\xi'$의 구성은 열린집합으로의 제한과 가환한다. 즉,
$U \subset X$가 열린집합이면 제한 $\xi'|_U$는
상 $\xi|_U \in H^n_{dR}(U/S)$에 대응하는 사상이며,
\item 주석 \ref{derham-remark-relative-cup-product}에서와 같은 임의의 도식에
대하여 가환 도식
$$
\xymatrix{
Rf_*\Omega^\bullet_{X/S}
\otimes_{q^{-1}\mathcal{O}_T}^\mathbf{L}
Rf_*\Omega^\bullet_{X/S} \ar[d]_{\xi' \otimes \text{id}}
\ar[r]_-\mu &
Rf_*\Omega^\bullet_{X/S} \ar[d]^{\xi'} \\
Rf_*\Omega^\bullet_{X/S}[n]
\otimes_{q^{-1}\mathcal{O}_T}^\mathbf{L}
Rf_*\Omega^\bullet_{X/S}
\ar[r]^-\mu &
Rf_*\Omega^\bullet_{X/S}[n]
}
$$
을 $D(Y, q^{-1}\mathcal{O}_T)$에서 얻는다.
\end{enumerate}
\end{remark}




\section{호지 코호몰로지}
\label{derham-section-hodge-cohomology}

\noindent
$p : X \to S$를 스킴의 사상이라 하자. {\it $X$의 $S$ 위 호지
코호몰로지}를 코호몰로지 군
$$
H^n_{Hodge}(X/S) = \bigoplus\nolimits_{n = p + q} H^q(X, \Omega^p_{X/S})
$$
으로 정의하고 이를 등급 $H^0(X, \mathcal{O}_X)$-가군으로 본다. 미분형식의
쐐기곱과 층 코호몰로지의 절 \ref{cohomology-section-cup-product}의 컵곱을
결합하면 $H^0(X, \mathcal{O}_X)$-쌍선형인 컵곱
$$
\cup :
H^i_{Hodge}(X/S) \times H^j_{Hodge}(X/S)
\longrightarrow
H^{i + j}_{Hodge}(X/S)
$$
을 정의한다. 물론 $\xi \in H^q(X, \Omega^p_{X/S})$이고
$\xi' \in H^{q'}(X, \Omega^{p'}_{X/S})$이면 $\xi \cup \xi' \in
H^{q + q'}(X, \Omega^{p + p'}_{X/S})$이다.

\begin{lemma}
\label{derham-lemma-cup-product-hodge-graded-commutative}
$p : X \to S$를 스킴의 사상이라 하자.
$H^*_{Hodge}(X/S)$ 위의 컵곱은 결합적이고 등급 가환이다.
\end{lemma}

\begin{proof}
증명은 보조정리 \ref{derham-lemma-cup-product-graded-commutative}의 증명과 같다.
\end{proof}

\noindent
스킴의 가환 도식
$$
\xymatrix{
X' \ar[r]_f \ar[d] & X \ar[d] \\
S' \ar[r] & S
}
$$
이 주어지면 당김 사상
$f^* : H^i_{Hodge}(X/S) \longrightarrow H^i_{Hodge}(X'/S')$
이 있으며, 이는 등급 및 위에서 정의한 컵곱과 양립한다.







\section{두 스펙트럼 수열}
\label{derham-section-hodge-to-de-rham}

\noindent
$p : X \to S$를 스킴의 사상이라 하자. $p^{-1}\mathcal{O}_S$-가군의
$X$ 위 범주에는 단사 대상이 충분히 있으므로
$\Omega^\bullet_{X/S}$의 카르탕--아일렌베르크 분해가 존재한다. 유도 범주의 보조정리 \ref{derived-lemma-cartan-eilenberg}를 보라. 따라서
유도 범주의 보조정리
\ref{derived-lemma-two-ss-complex-functor}를 적용하면 두 스펙트럼 수열을
얻으며, 둘 다 $X$의 $S$ 위 드람 코호몰로지로 수렴한다.

\medskip\noindent
첫 번째 것은 관례적으로 {\it 호지--드람 스펙트럼 수열}이라 부른다.
이 스펙트럼 수열의 첫째 쪽은
$$
E_1^{p, q} = H^q(X, \Omega^p_{X/S})
$$
이며, 이들은 $X/S$의 호지 코호몰로지 군이다(그래서 이 이름을 쓴다).
이 쪽의 미분 $d_1$은 사상
$d_1^{p, q} : H^q(X, \Omega^p_{X/S}) \to H^q(X, \Omega^{p + 1}_{X/S})$
로 주어지며, 이는 미분
$\text{d} : \Omega^p_{X/S} \to \Omega^{p + 1}_{X/S}$가 유도한다.
그림으로 나타내면 다음과 같다.
$$
\xymatrix{
H^2(X, \mathcal{O}_X) \ar[r] \ar@{-->}[rrd] \ar@{..>}[rrrdd] &
H^2(X, \Omega^1_{X/S}) \ar[r] \ar@{-->}[rrd] &
H^2(X, \Omega^2_{X/S}) \ar[r] &
H^2(X, \Omega^3_{X/S}) \\
H^1(X, \mathcal{O}_X) \ar[r] \ar@{-->}[rrd] &
H^1(X, \Omega^1_{X/S}) \ar[r] \ar@{-->}[rrd] &
H^1(X, \Omega^2_{X/S}) \ar[r] &
H^1(X, \Omega^3_{X/S}) \\
H^0(X, \mathcal{O}_X) \ar[r] &
H^0(X, \Omega^1_{X/S}) \ar[r] &
H^0(X, \Omega^2_{X/S}) \ar[r] &
H^0(X, \Omega^3_{X/S})
}
$$
여기서 파선 화살표는 $E_2$-쪽의 미분의 정의역과 공역을 나타내고,
점선 화살표는 $E_3$-쪽의 미분을 나타낸다. 차수 $0$을 살펴보면
$$
H^0_{dR}(X/S) =
\Ker(\text{d} : H^0(X, \mathcal{O}_X) \to H^0(X, \Omega^1_{X/S}))
$$
임을 얻는다. 물론 드람 복합체가 차수 $0$에서
$\mathcal{O}_X \to \Omega^1_{X/S}$로 시작한다는 사실에서도 이를 곧바로
알 수 있다.

\medskip\noindent
두 번째 스펙트럼 수열은 보통 {\it 켤레 스펙트럼 수열}이라 부른다.
이 스펙트럼 수열의 둘째 쪽은
$$
E_2^{p, q} = H^p(X, H^q(\Omega^\bullet_{X/S})) = H^p(X, \mathcal{H}^q)
$$
이다. 여기서 $\mathcal{H}^q = H^q(\Omega^\bullet_{X/S})$는 제$q$
코호몰로지 층이며, 여기서 쓰는 복합체는 $X/S$의 드람 복합체이다. 이
쪽의 미분은
$E_2^{p, q} \to E_2^{p + 2, q - 1}$로 주어진다. 그림으로 나타내면 다음과
같다.
$$
\xymatrix{
H^0(X, \mathcal{H}^2) \ar[rrd] \ar@{..>}[rrrdd] &
H^1(X, \mathcal{H}^2) \ar[rrd] &
H^2(X, \mathcal{H}^2) &
H^3(X, \mathcal{H}^2) \\
H^0(X, \mathcal{H}^1) \ar[rrd] &
H^1(X, \mathcal{H}^1) \ar[rrd] &
H^2(X, \mathcal{H}^1) &
H^3(X, \mathcal{H}^1) \\
H^0(X, \mathcal{H}^0) &
H^1(X, \mathcal{H}^0) &
H^2(X, \mathcal{H}^0) &
H^3(X, \mathcal{H}^0)
}
$$
차수 $0$을 살펴보면
$$
H^0_{dR}(X/S) = H^0(X, \mathcal{H}^0)
$$
임을 얻는데, 생각해 보면 자명하다. 차수 $1$에서는 완전열
$$
0 \to H^1(X, \mathcal{H}^0) \to H^1_{dR}(X/S) \to
H^0(X, \mathcal{H}^1) \to H^2(X, \mathcal{H}^0) \to H^2_{dR}(X/S)
$$
을 얻는다. $X \to S$가 매끄럽고 $S$의 표수가 $p$이면 층
$\mathcal{H}^q$를 어떤 미분형식 층들을 사용하여 계산할 수 있으며,
켤레 스펙트럼 수열은 유용한 도구가 된다(여기에 추후 참조를 삽입할 것).



\section{호지 여과}
\label{derham-section-hodge-filtration}

\noindent
$X \to S$를 스킴의 사상이라 하자. $H^n_{dR}(X/S)$ 위의 호지 여과는
호지--드람 스펙트럼 수열이 유도하는 여과이다(호몰로지 대수의 정의
\ref{homology-definition-filtration-cohomology-filtered-complex}). 오해를
피하기 위하여 이를 다음과 같이 명시적으로 정의한다.

\begin{definition}
\label{derham-definition-hodge-filtration}
$X \to S$를 스킴의 사상이라 하자. $H^n_{dR}(X/S)$ 위의 {\it 호지 여과}는
그 항들이
$$
F^pH^n_{dR}(X/S) = \Im\left(H^n(X, \sigma_{\geq p}\Omega^\bullet_{X/S})
\longrightarrow H^n_{dR}(X/S)\right)
$$
로 주어지는 여과이다. 여기서 $\sigma_{\geq p}\Omega^\bullet_{X/S}$는
호몰로지 대수의 절 \ref{homology-section-truncations}에서 정의한 것이다.
\end{definition}

\noindent
물론 $\sigma_{\geq p}\Omega^\bullet_{X/S}$는 상대 드람 복합체의
부분복합체이며, 상대 드람 복합체의 여과
$$
\Omega^\bullet_{X/S} = \sigma_{\geq 0}\Omega^\bullet_{X/S} \supset
\sigma_{\geq 1}\Omega^\bullet_{X/S} \supset
\sigma_{\geq 2}\Omega^\bullet_{X/S} \supset
\sigma_{\geq 3}\Omega^\bullet_{X/S} \supset \ldots
$$
를 얻고
$\text{gr}^p(\Omega^\bullet_{X/S}) = \Omega^p_{X/S}[-p]$가 성립한다.
$\Omega^\bullet_{X/S}$를 여과된 층 복합체로 보아 층 코호몰로지의 보조정리
\ref{cohomology-lemma-spectral-sequence-filtered-object}에서 구성한 스펙트럼
수열은 절 \ref{derham-section-hodge-to-de-rham}에서 구성한 호지--드람 스펙트럼
수열과 같다. 층 코호몰로지의 예
\ref{cohomology-example-spectral-sequence-bis}를 보라. 또한 쐐기곱
(\ref{derham-equation-wedge})는
$\text{Tot}(\sigma_{\geq i}\Omega^\bullet_{X/S} \otimes_{p^{-1}\mathcal{O}_S}
\sigma_{\geq j}\Omega^\bullet_{X/S})$를
$\sigma_{\geq i + j}\Omega^\bullet_{X/S}$로 보낸다. 따라서 가환 도식
$$
\xymatrix{
H^n(X, \sigma_{\geq i}\Omega^\bullet_{X/S}))
\times 
H^m(X, \sigma_{\geq j}\Omega^\bullet_{X/S}))
\ar[r] \ar[d] &
H^{n + m}(X, \sigma_{\geq i + j}\Omega^\bullet_{X/S})) \ar[d] \\
H^n_{dR}(X/S) \times
H^m_{dR}(X/S)
\ar[r]^\cup &
H^{n + m}_{dR}(X/S)
}
$$
을 얻는다. 특히
$$
F^iH^n_{dR}(X/S) \cup F^jH^m_{dR}(X/S) \subset F^{i + j}H^{n + m}_{dR}(X/S)
$$
임을 알 수 있다.






\section{퀴네트 공식}
\label{derham-section-kunneth}

\noindent
드람 코호몰로지의 중요한 특징 가운데 하나는 퀴네트 공식이 존재한다는
것이다.

\medskip\noindent
$a : X \to S$와 $b : Y \to S$를 공역이 같은 스킴의 사상이라 하자.
$p : X \times_S Y \to X$와 $q : X \times_S Y \to Y$를 사영 사상이라 하고
$f = a \circ p = b \circ q$로 놓자. 그림으로 나타내면 다음과 같다.
$$
\xymatrix{
& X \times_S Y \ar[ld]^p \ar[rd]_q \ar[dd]^f \\
X \ar[rd]_a & & Y \ar[ld]^b \\
& S
}
$$
이 절에서는 $\mathcal{O}_X$-가군 $\mathcal{F}$와
$\mathcal{O}_Y$-가군 $\mathcal{G}$가 주어졌을 때
$$
\mathcal{F} \boxtimes \mathcal{G} =
p^*\mathcal{F} \otimes_{\mathcal{O}_{X \times_S Y}} q^*\mathcal{G}
$$
로 놓는다. 준연접 가군 위의 쌍함자
$(\mathcal{F}, \mathcal{G}) \mapsto \mathcal{F} \boxtimes \mathcal{G}$
는 준연접 가군과 $S$ 위의 유한 차수 미분 작용소 위의 쌍함자로 확장된다.
스킴의 사상의 주석
\ref{morphisms-remark-base-change-differential-operators}을 보라. 드람
복합체 $\Omega^\bullet_{X/S}$와 $\Omega^\bullet_{Y/S}$의 미분은 가군층의
보조정리 \ref{modules-lemma-differentials-relative-de-rham-complex-order-1}에
의해 차수 $1$인 $S$ 위의 미분 작용소이다. 따라서 복합체
$$
\text{Tot}(\Omega^\bullet_{X/S} \boxtimes \Omega^\bullet_{Y/S})
$$
를 생각하는 것은 의미가 있다. 스킴의 유도 범주의 절
\ref{perfect-section-kunneth-complexes}의 논의를 보라.

\begin{lemma}
\label{derham-lemma-de-rham-complex-product}
위 상황에서 복합체의 표준 동형사상
$$
\text{Tot}(\Omega^\bullet_{X/S} \boxtimes \Omega^\bullet_{Y/S})
\longrightarrow
\Omega^\bullet_{X \times_S Y/S}
$$
이 존재하며, 이 복합체들은 $f^{-1}\mathcal{O}_S$-가군의 복합체이다.
\end{lemma}

\begin{proof}
스킴의 사상의 보조정리 \ref{morphisms-lemma-differential-product}에 의해
$
\Omega_{X \times_S Y/S} = p^*\Omega_{X/S} \oplus q^*\Omega_{Y/S}
$
임을 알고 있다. 외대수의 각 차수 성분을 취하면
$$
\Omega^n_{X \times_S Y/S} =
\bigoplus\nolimits_{i + j = n}
p^*\Omega^i_{X/S} \otimes_{\mathcal{O}_{X \times_S Y}} q^*\Omega^j_{Y/S} =
\bigoplus\nolimits_{i + j = n}
\Omega^i_{X/S} \boxtimes \Omega^j_{Y/S}
$$
를 얻는다. 이는 외대수의 기본 성질에서 따른다. 이 식별들은 보조정리의
화살표 왼쪽과 오른쪽에 있는 복합체의 항들 사이의 동형사상들을 정한다.
이 사상들이 미분과 양립한다는 확인은 생략한다.
\end{proof}

\noindent
$A = \Gamma(S, \mathcal{O}_S)$로 놓자. 보조정리
\ref{derham-lemma-de-rham-complex-product}의 결과와 스킴의 유도 범주의 식 (\ref{perfect-equation-de-rham-kunneth})의 사상을 결합하면
컵곱
$$
R\Gamma(X, \Omega^\bullet_{X/S})
\otimes_A^\mathbf{L}
R\Gamma(Y, \Omega^\bullet_{Y/S})
\longrightarrow
R\Gamma(X \times_S Y, \Omega^\bullet_{X \times_S Y/S})
$$
을 얻는다. 코호몰로지 수준에서는 대수학 심화의 절
\ref{more-algebra-section-products-tor}의 논의를 사용하여 표준 사상
$$
H^i_{dR}(X/S) \otimes_A H^j_{dR}(Y/S)
\longrightarrow
H^{i + j}_{dR}(X \times_S Y/S),\quad
(\xi, \zeta) \longmapsto p^*\xi \cup q^*\zeta
$$
을 얻는다. 위 구성은 실제로 먼저 당긴 다음 컵곱을 취하는 순서로
이루어진다는 점에 유의하자.

\begin{lemma}
\label{derham-lemma-kunneth-de-rham}
$X$와 $Y$가 $S = \Spec(A)$ 위에서 매끄럽고 준콤팩트이며 대각 사상이
아핀이라고 가정하자. 그러면 사상
$$
R\Gamma(X, \Omega^\bullet_{X/S})
\otimes_A^\mathbf{L}
R\Gamma(Y, \Omega^\bullet_{Y/S})
\longrightarrow
R\Gamma(X \times_S Y, \Omega^\bullet_{X \times_S Y/S})
$$
은 $D(A)$에서 동형사상이다.
\end{lemma}

\begin{proof}
스킴의 사상의 보조정리
\ref{morphisms-lemma-smooth-omega-finite-locally-free}에 의해 층
$\Omega^n_{X/S}$와 $\Omega^m_{Y/S}$는 각각 유한 국소 자유
$\mathcal{O}_X$-가군과 $\mathcal{O}_Y$-가군이다. 한편 $X$와 $Y$는
$S$ 위에서 평탄하다(스킴의 사상의 보조정리
\ref{morphisms-lemma-smooth-flat}). 따라서 $\Omega^n_{X/S}$와
$\Omega^m_{Y/S}$는 $S$ 위에서 평탄하다. 또한
$\Omega^\bullet_{X/S}$는 국소 유계임에 유의하자. 이제 보조정리
\ref{derham-lemma-de-rham-complex-product}와 스킴의 유도 범주의
보조정리 \ref{perfect-lemma-kunneth-special}을 적용하면 결론을 얻는다.
\end{proof}

\noindent
컵곱에는 상대적 판본, 즉 사상
$$
Ra_*\Omega^\bullet_{X/S}
\otimes_{\mathcal{O}_S}^\mathbf{L}
Rb_*\Omega^\bullet_{Y/S}
\longrightarrow
Rf_*\Omega^\bullet_{X \times_S Y/S}
$$
이 $D(\mathcal{O}_S)$에 존재한다. 이 구성은 보조정리
\ref{derham-lemma-de-rham-complex-product}와 스킴의 유도 범주의 식
(\ref{perfect-equation-relative-de-rham-kunneth})의 사상을 결합한다.
이 구성에 따르면 이 사상은 다음 도식으로 주어진다.
$$
\xymatrix{
Ra_*\Omega^\bullet_{X/S}
\otimes_{\mathcal{O}_S}^\mathbf{L}
Rb_*\Omega^\bullet_{Y/S}
\ar[d]^{\text{수반의 단위원}}  \\
Rf_*(p^{-1}\Omega^\bullet_{X/S})
\otimes_{\mathcal{O}_S}^\mathbf{L}
Rf_*(q^{-1}\Omega^\bullet_{Y/S}) \ar[r] \ar[d]^{\text{상대 컵곱}} &
Rf_*(\Omega^\bullet_{X \times_S Y/S})
\otimes_{\mathcal{O}_S}^\mathbf{L}
Rf_*(\Omega^\bullet_{X \times_S Y/S}) \ar[d]^{\text{상대 컵곱}} \\
Rf_*(p^{-1}\Omega^\bullet_{X/S}
\otimes_{f^{-1}\mathcal{O}_S}^\mathbf{L}
q^{-1}\Omega^\bullet_{Y/S})
\ar[d]^{\text{유도 텐서곱에서 보통 텐서곱으로}} \ar[r] &
Rf_*(\Omega^\bullet_{X \times_S Y/S}
\otimes_{f^{-1}\mathcal{O}_S}^\mathbf{L}
\Omega^\bullet_{X \times_S Y/S})
\ar[d]^{\text{유도 텐서곱에서 보통 텐서곱으로}} \\
Rf_*\text{Tot}(p^{-1}\Omega^\bullet_{X/S}
\otimes_{f^{-1}\mathcal{O}_S}
q^{-1}\Omega^\bullet_{Y/S}) \ar[r] \ar[d]^{\text{표준 사상}} &
Rf_*\text{Tot}(\Omega^\bullet_{X \times_S Y/S}
\otimes_{f^{-1}\mathcal{O}_S}
\Omega^\bullet_{X \times_S Y/S})
\ar[d]^{\eta \otimes \omega \mapsto \eta \wedge \omega} \\
Rf_*\text{Tot}(\Omega^\bullet_{X/S} \boxtimes \Omega^\bullet_{Y/S})
\ar@{=}[r]
&
Rf_*\Omega^\bullet_{X \times_S Y/S}
}
$$
여기서 첫 번째 화살표는 수반의 단위원
$\text{id} \to Rp_* p^{-1}$와 $\text{id} \to Rq_* q^{-1}$, 그리고 식별
$Rf_* p^{-1} = Ra_* Rp_* p^{-1}$와
$Rf_* q^{-1} = Rb_* Rq_* q^{-1}$를 사용한다. 두 번째 화살표는 층 코호몰로지의
주석 \ref{cohomology-remark-cup-product}의 상대 컵곱이다. 세 번째 화살표는
복합체의 유도 텐서곱을 복합체 텐서곱의 전체화로 보내는 사상이다. 마지막
등식은 보조정리 \ref{derham-lemma-de-rham-complex-product}이다. 이 구성은 전역
절단 위에서 앞서 제시한 구성을 되찾는다.

\begin{lemma}
\label{derham-lemma-kunneth-de-rham-relative}
$X \to S$와 $Y \to S$가 매끄럽고 준콤팩트이며, 사상
$X \to X \times_S X$와 $Y \to Y \times_S Y$가 아핀이라고 가정하자.
그러면 상대 컵곱
$$
Ra_*\Omega^\bullet_{X/S}
\otimes_{\mathcal{O}_S}^\mathbf{L}
Rb_*\Omega^\bullet_{Y/S}
\longrightarrow
Rf_*\Omega^\bullet_{X \times_S Y/S}
$$
은 $D(\mathcal{O}_S)$에서 동형사상이다.
\end{lemma}

\begin{proof}
보조정리 \ref{derham-lemma-kunneth-de-rham}의 즉각적인 귀결이다.
\end{proof}







\section{드람 코호몰로지의 제1 천 특성류}
\label{derham-section-first-chern-class}

\noindent
$X \to S$를 스킴의 사상이라 하자. 복합체의 사상
$$
\text{d}\log : \mathcal{O}_X^*[-1] \longrightarrow \Omega^\bullet_{X/S}
$$
이 존재하며, 이는 절단 $g \in \mathcal{O}_X^*(U)$를 절단
$\text{d}\log(g) = g^{-1}\text{d}g$, 즉
$\Omega^1_{X/S}(U)$의 절단으로 보낸다.
따라서 사상
$$
\Pic(X) = H^1(X, \mathcal{O}_X^*) =
H^2(X, \mathcal{O}_X^*[-1]) \longrightarrow H^2_{dR}(X/S)
$$
을 생각할 수 있다. 여기서 첫 번째 등식은 층 코호몰로지의 보조정리
\ref{cohomology-lemma-h1-invertible}이다.
가역 가군 $\mathcal{L}$의 동형류의 상을
$c^{dR}_1(\mathcal{L}) \in H^2_{dR}(X/S)$로 표기한다.

\medskip\noindent
또한 사상 $\text{d}\log : \mathcal{O}_X^* \to \Omega^1_{X/S}$를
사용하여 호지 코호몰로지의 천 특성류
$$
c_1^{Hodge} : \Pic(X) \longrightarrow H^1(X, \Omega^1_{X/S})
\subset H^2_{Hodge}(X/S)
$$
를 정의할 수 있다. 이 구성들은 당김과 양립한다.

\begin{lemma}
\label{derham-lemma-pullback-c1}
스킴의 가환도식
$$
\xymatrix{
X' \ar[r]_f \ar[d] & X \ar[d] \\
S' \ar[r] & S
}
$$
이 주어지면 다음 도식들은 가환한다.
$$
\xymatrix{
\Pic(X') \ar[d]_{c_1^{dR}} &
\Pic(X) \ar[d]^{c_1^{dR}} \ar[l]^{f^*} \\
H^2_{dR}(X'/S') &
H^2_{dR}(X/S) \ar[l]_{f^*}
}
\quad
\xymatrix{
\Pic(X') \ar[d]_{c_1^{Hodge}} &
\Pic(X) \ar[d]^{c_1^{Hodge}} \ar[l]^{f^*} \\
H^1(X', \Omega^1_{X'/S'}) &
H^1(X, \Omega^1_{X/S}) \ar[l]_{f^*}
}
$$
\end{lemma}

\begin{proof}
생략한다.
\end{proof}

\noindent
Čech 코호몰로지에서 원소 $c^{dR}_1(\mathcal{L})$를 ``계산''해 보자
(Čech 미분의 부호 규칙은 층 코호몰로지의 절
\ref{cohomology-section-cech-cohomology-of-complexes}에서와 같다).
구체적으로, 열린 덮개
$\mathcal{U} : X = \bigcup_{i \in I} U_i$를 택하되, $s_i$가
$\mathcal{L}|_{U_i}$의 자명화 절단이 되게 한다. 이는 모든 $i$에
대해 이루어진다.
겹침 $U_{i_0i_1} = U_{i_0} \cap U_{i_1}$ 위에는
가역 함수 $f_{i_0i_1}$가 있어서
$f_{i_0i_1} = s_{i_1}|_{U_{i_0i_1}} s_{i_0}|_{U_{i_0i_1}}^{-1}$이다\footnote{$0$-사이클
$\{a_{i_0}\}$의 Čech 미분은 $a_{i_1} - a_{i_0}$이며, 이는
$U_{i_0i_1}$ 위의 식이다.}.
물론
$$
f_{i_1i_2}|_{U_{i_0i_1i_2}}
f_{i_0i_2}^{-1}|_{U_{i_0i_1i_2}}
f_{i_0i_1}|_{U_{i_0i_1i_2}} = 1
$$
이다. $\mathcal{L}$의 $H^1(X, \mathcal{O}_X^*)$에서의 코호몰로지류는
코사이클 $\{f_{i_0i_1}\}$의
$\check{\mathcal{C}}^\bullet(\mathcal{U}, \mathcal{O}_X^*)$에서의
Čech 코호몰로지류의 상이다.
따라서 $c_1^{dR}(\mathcal{L})$은
Čech 코사이클 $\{\alpha_{i_0 \ldots i_p}\}$에 대응하는
코호몰로지류의 상이다. 이 코사이클은
$\text{Tot}(\check{\mathcal{C}}^\bullet(\mathcal{U}, \Omega_{X/S}^\bullet))$
안에서 차수 $2$이고 다음과 같이 주어진다.
\begin{enumerate}
\item $\alpha_{i_0} = 0$이며, 이는 $\Omega^2_{X/S}(U_{i_0})$에서의
등식이다.
\item $\alpha_{i_0i_1} = f_{i_0i_1}^{-1}\text{d}f_{i_0i_1}$이며, 이는
$\Omega^1_{X/S}(U_{i_0i_1})$에서의 등식이다.
\item $\alpha_{i_0i_1i_2} = 0$이며, 이는
$\mathcal{O}_{X/S}(U_{i_0i_1i_2})$에서의 등식이다.
\end{enumerate}
가역 가군 $\mathcal{L}_k$, $k = 1, \ldots, a$가 주어져 있고,
각각이 $U_i$ 위에서 자명화되어 있다고 하자. 이는 모든 $i \in I$에
대해 성립하며, 위와 같은
코사이클 $f_{k, i_0i_1}$와
$\alpha_k = \{\alpha_{k, i_0 \ldots i_p}\}$를 준다고 하자.
층 코호몰로지의 절 \ref{cohomology-section-cech-cohomology-of-complexes}의
규칙을 사용하면
$$
\beta = \alpha_1 \cup \alpha_2 \cup \ldots \cup \alpha_a
$$
는 다음과 같은 코사이클 $\beta = \{\beta_{i_0 \ldots i_p}\}$로
주어짐을 계산할 수 있다.
\begin{enumerate}
\item $\beta_{i_0 \ldots i_p} = 0$이며, 이는
$\Omega^{2a - p}_{X/S}(U_{i_0 \ldots i_p})$에서의 등식이다. 단,
$p = a$인 경우는 제외한다.
\item $\beta_{i_0 \ldots i_a} = (-1)^{a(a - 1)/2}
\alpha_{1, i_0i_1} \wedge \alpha_{2, i_1 i_2} \wedge \ldots \wedge
\alpha_{a, i_{a - 1}i_a}$이며, 이는
$\Omega^a_{X/S}(U_{i_0 \ldots i_a})$에서의 등식이다.
\end{enumerate}
따라서 이는
$c_1^{dR}(\mathcal{L}_1) \cup \ldots \cup c_1^{dR}(\mathcal{L}_a)$를
대표하는 코사이클이다. 물론 같은 계산으로 코사이클
$\{\beta_{i_0 \ldots i_a}\}$가
$\check{\mathcal{C}}^a(\mathcal{U}, \Omega_{X/S}^a))$ 안에서 코호몰로지류
$c_1^{Hodge}(\mathcal{L}_1) \cup \ldots \cup c_1^{Hodge}(\mathcal{L}_a)$를
대표함을 알 수 있다.

\begin{remark}
\label{derham-remark-truncations}
위의 계산을 더 추상적인 용어로 다시 서술해 보자.
$p : X \to S$를 스킴의 사상이라 하고, $\mathcal{L}$을 가역
$\mathcal{O}_X$-가군이라 하자. $\text{d}\log$를 사상
$$
\mathcal{O}_X^*[-1] \to \sigma_{\geq 1}\Omega^\bullet_{X/S}
$$
으로 보면, 위와 같이 $\Pic(X) = H^1(X, \mathcal{O}_X^*)$를 사용하여
코호몰로지류
$$
\gamma_1(\mathcal{L}) \in H^2(X, \sigma_{\geq 1}\Omega^\bullet_{X/S})
$$
를 얻는다. $\gamma_1(\mathcal{L})$의 사상
$\sigma_{\geq 1}\Omega^\bullet_{X/S} \to \Omega^\bullet_{X/S}$에 의한
상은 $c_1^{dR}(\mathcal{L})$이다. 특히
$c_1^{dR}(\mathcal{L}) \in F^1H^2_{dR}(X/S)$임을 알 수 있다.
절 \ref{derham-section-hodge-filtration}을 보라. $\gamma_1(\mathcal{L})$의 사상
$\sigma_{\geq 1}\Omega^\bullet_{X/S} \to \Omega^1_{X/S}[-1]$에 의한
상은 $c_1^{Hodge}(\mathcal{L})$이다.
컵곱을 취하면(절 \ref{derham-section-hodge-filtration}을 보라)
$$
\xi = \gamma_1(\mathcal{L}_1) \cup \ldots \cup \gamma_1(\mathcal{L}_a) \in
H^{2a}(X, \sigma_{\geq a}\Omega^\bullet_{X/S})
$$
를 얻는다. 절 \ref{derham-section-hodge-filtration}의 가환도식들에 따르면,
$\xi$는
$c_1^{dR}(\mathcal{L}_1) \cup \ldots \cup c_1^{dR}(\mathcal{L}_a)$
의 $H^{2a}_{dR}(X/S)$ 안의 원소로 사상
$\sigma_{\geq a}\Omega^\bullet_{X/S} \to \Omega^\bullet_{X/S}$에 의해
보내진다. 또한 이로부터
$c_1^{dR}(\mathcal{L}_1) \cup \ldots \cup c_1^{dR}(\mathcal{L}_a)$
가 $F^a H^{2a}_{dR}(X/S)$에 속함이 따른다. 마찬가지로 사상
$\sigma_{\geq a}\Omega^\bullet_{X/S} \to \Omega^a_{X/S}[-a]$
은 $\xi$를
$c_1^{Hodge}(\mathcal{L}_1) \cup \ldots \cup c_1^{Hodge}(\mathcal{L}_a)$
의 $H^a(X, \Omega^a_{X/S})$ 안의 원소로 보낸다.
\end{remark}

\begin{remark}
\label{derham-remark-log-forms}
$p : X \to S$를 스킴의 사상이라 하자. $i > 0$일 때,
다음 꼴의 국소 절단들로 생성되는 아벨 부분층
$\Omega^i_{X/S, log} \subset \Omega^i_{X/S}$를 생각하자.
$$
\text{d}\log(u_1) \wedge \ldots \wedge \text{d}\log(u_i)
$$
여기서 $u_1, \ldots, u_n$은 $\mathcal{O}_X$의 가역 국소 절단들이다.
$i = 0$일 때 부분층 $\Omega^0_{X/S, log} \subset \mathcal{O}_X$는
$\mathbf{Z} \to \mathcal{O}_X$의 상이다. 모든 $i \geq 0$에 대해
복합체의 사상
$$
\Omega^i_{X/S, log}[-i] \longrightarrow \Omega^\bullet_{X/S}
$$
이 존재한다. 로그 형식의 미분이 영이기 때문이다. 또한 로그 형식들의
쐐기곱은 다시 로그 형식이므로 쌍선형 사상
$$
\wedge :  \Omega^i_{X/S, log} \times
\Omega^j_{X/S, log} \longrightarrow \Omega^{i + j}_{X/S, log}
$$
을 얻으며, 이는 (\ref{derham-equation-wedge}) 및 위의 사상들과 양립한다.
$\mathcal{L}$을 가역 $\mathcal{O}_X$-가군이라 하자. 아벨 층의 사상
$\text{d}\log : \mathcal{O}_X^* \to \Omega^1_{X/S, log}$
과 식별 $\Pic(X) = H^1(X, \mathcal{O}_X^*)$을 사용하면 표준 코호몰로지류
$$
\tilde \gamma_1(\mathcal{L}) \in H^1(X, \Omega^1_{X/S, log})
$$
를 얻는다. 이 코호몰로지류들은 다음 성질들을 갖는다.
\begin{enumerate}
\item $\tilde \gamma_1(\mathcal{L})$의 표준 사상
$\Omega^1_{X/S, log}[-1] \to \sigma_{\geq 1}\Omega^\bullet_{X/S}$에
의한 상은 $\tilde \gamma_1(\mathcal{L})$을 류
$\gamma_1(\mathcal{L}) \in 
H^2(X, \sigma_{\geq 1}\Omega^\bullet_{X/S})$
로 보낸다. 여기서 이 류는 주석 \ref{derham-remark-truncations}의 류이다.
\item $\tilde \gamma_1(\mathcal{L})$의 표준 사상
$\Omega^1_{X/S, log}[-1] \to \Omega^\bullet_{X/S}$에 의한 상은
$\tilde \gamma_1(\mathcal{L})$을 $c_1^{dR}(\mathcal{L})$의
$H^2_{dR}(X/S)$ 안의 원소로 보낸다.
\item $\tilde \gamma_1(\mathcal{L})$의 표준 사상
$\Omega^1_{X/S, log} \to \Omega^1_{X/S}$에 의한 상은
$\tilde \gamma_1(\mathcal{L})$을 $c_1^{Hodge}(\mathcal{L})$의
$H^1(X, \Omega^1_{X/S})$ 안의 원소로 보낸다.
\item 이 류들의 구성은 당김과 양립한다.
\item 여기에 더 추가한다.
\end{enumerate}
\end{remark}





\section{선다발의 드람 코호몰로지}
\label{derham-section-line-bundle}

\noindent
선다발은 벡터 다발의 특수한 경우이고, 벡터 다발은 다시 어떤 추가 구조가
주어진 원뿔이다. 이들의 드람 복합체를 제대로 논하려면 먼저 등급환의 드람
복합체를 논하는 것이 자연스럽다.

\begin{remark}[등급환의 드람 복합체]
\label{derham-remark-de-rham-complex-graded}
$G$를 중성원 $0$을 갖고 덧셈으로 표기한 아벨 모노이드라 하자.
$R \to A$를 환 사상이라 하고, $A$에 등급
$A = \bigoplus_{g \in G} A_g$가 주어졌다고 하자. 이는 $R$-가군들에 의한
등급이고, $R$은 $A_0$으로 가며
$A_g \cdot A_{g'} \subset A_{g + g'}$라고 하자. 그러면 미분 가군에는 등급
$$
\Omega_{A/R} = \bigoplus\nolimits_{g \in G} \Omega_{A/R, g}
$$
이 주어진다. 여기서 $\Omega_{A/R, g}$는 $R$-부분가군이고
$\Omega_{A/R}$에 들어가며, $a_0 \text{d}a_1$들로 생성된다. 여기서
$a_i \in A_{g_i}$이고 $g = g_0 + g_1$이다. 마찬가지로
$$
\Omega^p_{A/R} = \bigoplus\nolimits_{g \in G} \Omega^p_{A/R, g}
$$
를 얻는다. 여기서 $\Omega^p_{A/R, g}$는 $R$-부분가군이고
$\Omega^p_{A/R}$에 들어가며,
$a_0 \text{d}a_1 \wedge \ldots \wedge \text{d}a_p$들로 생성된다.
여기서 $a_i \in A_{g_i}$이고 $g = g_0 + g_1 + \ldots + g_p$이다.
물론 미분은 등급을 보존하고, 쐐기곱은 자명한 방식으로 등급들과 양립한다.
\end{remark}

\noindent
$f : X \to S$를 스킴의 사상이라 하자. $\pi : C \to X$를 원뿔이라 하자.
스킴의 구성의 정의 \ref{constructions-definition-abstract-cone}를 보라.
이는 $\pi$가 아핀이고 등급
$\pi_*\mathcal{O}_C = \bigoplus_{n \geq 0} \mathcal{A}_n$이 있으며
$\mathcal{A}_0 = \mathcal{O}_X$라는 뜻임을 상기하자.
아핀 열린집합 위에서 주석 \ref{derham-remark-de-rham-complex-graded}의 논의를
사용하면\footnote{이 표기법에 대해서는 양해를 구한다!}
$$
\pi_*(\Omega^\bullet_{C/S}) =
\bigoplus\nolimits_{n \geq 0} \Omega^\bullet_{C/S, n}
$$
가 부분복합체들의 표준 직합임을 알 수 있다. 더욱이 분해
$$
\Omega^\bullet_{X/S} \to \Omega^\bullet_{C/S, 0} \to
\pi_*(\Omega^\bullet_{C/S})
$$
가 있고, $\omega \wedge \eta \in \Omega^{p + q}_{C/S, n + m}$임을 안다.
여기서 $\omega \in \Omega^p_{C/S, n}$이고
$\eta \in \Omega^q_{C/S, m}$이다.

\medskip\noindent
$f : X \to S$를 스킴의 사상이라 하자. $\pi : L \to X$를 가역
$\mathcal{O}_X$-가군 $\mathcal{L}$에 결부한 선다발이라 하자.
이는 $\pi$가 다음 등식을 만족하는 유일한 아핀 사상이라는 뜻이다.
$$
\pi_*\mathcal{O}_L = 
\bigoplus\nolimits_{n \geq 0} \mathcal{L}^{\otimes n}
$$
이는 $\mathcal{O}_X$-대수의 등식이다. 따라서 $L$은 $X$ 위의 원뿔이다.
위의 논의로부터 표준 직합 분해
$$
\pi_*(\Omega^\bullet_{L/S}) =
\bigoplus\nolimits_{n \geq 0} \Omega^\bullet_{L/S, n}
$$
를 얻는다. 이는 쐐기곱 및 위의 $\pi_*\mathcal{O}_L$의 분해와 양립하고,
$\Omega_{X/S}$는 부분 $\Omega_{L/S, 0}$으로 가며 이 부분의 차수는 $0$이다.

\medskip\noindent
우리에게 유용할 또 다른 경우가 있다. 즉, 다음 여집합을
생각하자\footnote{스킴 $L^\star$는 $\mathbf{G}_m$-토서로서 $X$ 위에 있고
$L$에 결부한다. 이것이 아래에서 얻는 등급이 $\mathbf{Z}$-등급인 이유이다.
준군 스킴의 예
\ref{groupoids-example-Gm-on-affine} 및 보조정리
\ref{groupoids-lemma-complete-reducibility-Gm}와
\ref{groupoids-lemma-Gm-equivariant-module}를 비교하라.}
$L^\star \subset L$, 즉 영절단 $o : X \to L$의 선다발 $L$ 안에서의
여집합이다. 국소 계산으로 표준 동형
$$
(L^\star \to X)_*\mathcal{O}_{L^\star} =
\bigoplus\nolimits_{n \in \mathbf{Z}} \mathcal{L}^{\otimes n}
$$
을 얻으며, 이는 $\mathcal{O}_X$-대수의 동형이다. 우변은
$\mathbf{Z}$-등급 준연접 $\mathcal{O}_X$-대수이다.
아핀 열린집합 위에서 주석 \ref{derham-remark-de-rham-complex-graded}의 논의를
사용하면
$$
(L^\star \to X)_*(\Omega^\bullet_{L^\star/S}) =
\bigoplus\nolimits_{n \in \mathbf{Z}} \Omega^\bullet_{L^\star/S, n}
$$
를 얻는다. 이는 쐐기곱 및 위의 $(L^\star \to X)_*\mathcal{O}_{L^\star}$의
분해와 양립하며, $\Omega_{X/S}$는 부분 $\Omega_{L^\star/S, 0}$으로 가고
이 부분의 차수는 $0$이다. 우리는 특히 복합체
$\Omega^\bullet_{L^\star/S, 0}$에 관심을 둘 것이다.

\begin{lemma}
\label{derham-lemma-the-complex-for-L-star}
위의 표기 아래 복합체의 짧은 완전열
$$
0 \to \Omega^\bullet_{X/S} \to
\Omega^\bullet_{L^\star/S, 0} \to
\Omega^\bullet_{X/S}[-1] \to 0
$$
\end{lemma}

\begin{proof}
위에서 사상
$\Omega^\bullet_{X/S} \to \Omega^\bullet_{L^\star/S, 0}$을 구성하였다.

\medskip\noindent
$\text{Res} : \Omega^\bullet_{L^\star/S, 0} \to \Omega^\bullet_{X/S}[-1]$의
구성. $U \subset X$를 열린집합이라 하고, 절단
$s \in \mathcal{L}(U)$와 $s' \in \mathcal{L}^{\otimes -1}(U)$가
$s' s = 1$을 만족한다고 하자. 그러면 $s$는 대수층
$(L^\star \to X)_*\mathcal{O}_{L^\star}$의 $U$ 위 가역 절단을 주며 그 역은
$s' = s^{-1}$이다. 따라서 $1$-형식
$\text{d}\log(s) = s' \text{d}(s)$를 생각할 수 있다. 이는
$\Omega^1_{L^\star/S, 0}(U)$의 원소이며, 그 사실은
$\Omega^1_{L^\star/S}$에 둔 등급의 구성에 따른다. 아래의 아핀 계산으로부터
$1$과 $\text{d}\log(s)$가 $\Omega^\bullet_{L^\star/S, 0}|_U$를 자유롭게
생성함을 알 수 있다. 여기서 이는 $\Omega^\bullet_{X/S}|_U$ 위의 오른쪽
가군이다. 따라서 $\text{Res}$를 $U$ 위에서 다음 규칙으로 정의할 수 있다.
$$
\text{Res}(\omega' + \text{d}\log(s) \wedge \omega) = \omega
$$
이는 모든 $\omega', \omega \in \Omega^\bullet_{X/S}(U)$에 대해 적용한다.
이 사상은 국소 생성원 $s$의 선택과 무관하므로 붙여서 전역 사상을 준다.
실제로 $s$의 다른 선택은 $gs$ 꼴이고, 여기서
$g \in \mathcal{O}_X(U)$는 가역원이다. 이때
$\text{d}\log(gs) = g^{-1}\text{d}(g) + \text{d}\log(s)$를 얻으므로
독립성이 쉽게 따른다. 마지막으로 $\text{Res}$에 대한 이 규칙은 미분과
양립한다. 실제로
$\text{d}(\omega' + \text{d}\log(s) \wedge \omega) =
\text{d}(\omega') - \text{d}\log(s) \wedge \text{d}(\omega)$이고,
호몰로지 대수의 정의 \ref{homology-definition-shift-cochain}에서 택한 부호 규약에
따라 $\Omega^\bullet_{X/S}[-1]$의 미분은 $\omega'$를
$-\text{d}(\omega')$로 보낸다.

\medskip\noindent
국소 계산. $X$를 아핀 열린집합 $U \subset X$들로 덮되
$\mathcal{L}|_U \cong \mathcal{O}_U$이고, 각 열린집합이 어떤 아핀 열린집합
$V \subset S$로 가게 할 수 있다. $U = \Spec(A)$, $V = \Spec(R)$로 쓰고
$s$를 생성원으로 택하되, 이는 $\mathcal{L}$의 생성원이다. 그러면
$$
L^\star \times_X U = \Spec(A[s, s^{-1}])
$$
이다. 미분을 계산하면
$$
\Omega^1_{A[s, s^{-1}]/R} =
A[s, s^{-1}] \otimes_A \Omega^1_{A/R} \oplus A[s, s^{-1}] \text{d}\log(s)
$$
이고, 따라서 외대수의 각 차수 성분을 취하여
$$
\Omega^p_{A[s, s^{-1}]/R} =
A[s, s^{-1}] \otimes_A \Omega^p_{A/R}
\oplus
A[s, s^{-1}] \text{d}\log(s) \otimes_A \Omega^{p - 1}_{A/R}
$$
을 얻는다. 차수 $0$인 부분을 취하면
$$
\Omega^p_{A[s, s^{-1}]/R, 0} =
\Omega^p_{A/R} \oplus \text{d}\log(s) \otimes_A \Omega^{p - 1}_{A/R}
$$
을 얻으며, 보조정리의 증명이 끝난다.
\end{proof}

\begin{lemma}
\label{derham-lemma-the-complex-for-L-star-gives-chern-class}
보조정리 \ref{derham-lemma-the-complex-for-L-star}의 짧은 완전열에서 나오는
``경계'' 사상
$\delta : \Omega^\bullet_{X/S} \to \Omega^\bullet_{X/S}[2]$
은 $D(X, f^{-1}\mathcal{O}_S)$ 안에서
$\xi = c_1^{dR}(\mathcal{L})$에 대한 주석
\ref{derham-remark-cup-product-as-a-map}의 사상이다.
\end{lemma}

\begin{proof}
정확히 말하면, 보조정리 \ref{derham-lemma-the-complex-for-L-star}의 짧은 완전열을
밀어서 얻는
$$
0 \to \Omega^\bullet_{X/S}[1] \to
\Omega^\bullet_{L^\star/S, 0}[1] \to
\Omega^\bullet_{X/S} \to 0
$$
을 생각한다. 등급이 주어진
$(L^\star \to X)_*\Omega^\bullet_{L^\star/S}$
의 차수 영인 부분이므로, $\Omega^\bullet_{L^\star/S, 0}$은 미분 등급
$\mathcal{O}_X$-대수이고 사상
$\Omega^\bullet_{X/S} \to \Omega^\bullet_{L^\star/S, 0}$
은 미분 등급 $\mathcal{O}_X$-대수의 준동형임을 알 수 있다. 따라서
$\Omega^\bullet_{X/S}[1] \to
\Omega^\bullet_{L^\star/S, 0}[1]$을 오른쪽 미분 등급
$\Omega^\bullet_{X/S}$-가군들의 $X$ 위 사상으로 볼 수 있다. 사상
$\text{Res} : \Omega^\bullet_{L^\star/S, 0}[1] \to \Omega^\bullet_{X/S}$
은 오른쪽 미분 등급 $\Omega^\bullet_{X/S}$-가군들의 사상이다. 실제로 이는
국소적으로 규칙
$\text{Res}(\omega' + \text{d}\log(s) \wedge \omega) = \omega$로 정의된다.
보조정리 \ref{derham-lemma-the-complex-for-L-star}의 증명을 보라.
그러므로 미분 등급 층의 절 \ref{sdga-section-misc}의 논의에
의해 $\delta$는 사상
$\delta' : \Omega^\bullet_{X/S} \to \Omega^\bullet_{X/S}[2]$에서 온다.
이는 드람 복합체 위의 오른쪽 미분 등급 가군의 유도범주
$D(\Omega^\bullet_{X/S}, \text{d})$ 안의 사상이다.
주석 \ref{derham-remark-cup-product-as-a-map}에서 단언한 유일성에 의해
$\delta(1) = c_1^{dR}(\mathcal{L})$임을 보이면 충분하다.

\medskip\noindent
다음 가환도식이 존재한다고 주장한다.
$$
\xymatrix{
0 \ar[r] &
\mathcal{O}_X^* \ar[r] \ar[d]_{\text{d}\log} &
E \ar[r] \ar[d] &
\underline{\mathbf{Z}} \ar[d] \ar[r] &
0 \\
0 \ar[r] &
\Omega^\bullet_{X/S}[1] \ar[r] &
\Omega^\bullet_{L^\star/S, 0}[1] \ar[r] &
\Omega^\bullet_{X/S} \ar[r] &
0
}
$$
여기서 윗행은 아벨 층들의 짧은 완전열이고, 그 경계 사상은 $1$을
$\mathcal{L}$의 $H^1(X, \mathcal{O}_X^*)$ 안의 류로 보낸다. 경계 사상이
짧은 완전열 사이의 사상과 양립하므로 이 주장을 증명하면 충분하다.
$E$를 다음 규칙의 층화로 정의한다.
$$
U \longmapsto \{(s, n) \mid
n \in \mathbf{Z},\ s \in \mathcal{L}^{\otimes n}(U)\text{ 생성원}\}
$$
군 구조는 $(s, n) \cdot (t, m) = (s \otimes t, n + m)$로 주어진다.
가운데 세로 사상은 $(s, n)$을 $\text{d}\log(s)$로 보낸다. 이는 짧은
완전열들의 사상을 준다. 실제로 사상
$Res : \Omega^1_{L^\star/S, 0} \to \mathcal{O}_X$을 보조정리
\ref{derham-lemma-the-complex-for-L-star}의 증명에서 구성하였다. 이 사상은
$\text{d}\log(s)$를 $1$로 보낸다. 여기서 $s$는 $\mathcal{L}$의 국소
생성원이다. 윗행에서 $1$의 경계를 계산하기 위해 국소 자명화 $s_i$를
택하자. 이는 $\mathcal{L}$의 것이고 열린집합 $U_i$ 위에 있다. 절
\ref{derham-section-first-chern-class}에서와 같이 택한다. 겹침
$U_{i_0i_1} = U_{i_0} \cap U_{i_1}$
위에는 다음을 만족하는 가역 함수 $f_{i_0i_1}$가 있다.
$f_{i_0i_1} = s_{i_1}|_{U_{i_0i_1}} s_{i_0}|_{U_{i_0i_1}}^{-1}$
그리고 $\mathcal{L}$의 코호몰로지류는 Čech 코사이클
$\{f_{i_0i_1}\}$로 주어진다. 그러면 물론
$$
(f_{i_0i_1}, 0) = (s_{i_1}, 1)|_{U_{i_0i_1}} \cdot
(s_{i_0}, 1)|_{U_{i_0i_1}}^{-1}
$$
가 $E$의 절단으로서 성립하며, 이로써 증명이 끝난다.
\end{proof}

\begin{lemma}
\label{derham-lemma-push-omega-a}
위의 표기 아래 다음이 성립한다.
\begin{enumerate}
\item $\Omega^p_{L^\star/S, n} =
\Omega^p_{L^\star/S, 0} \otimes_{\mathcal{O}_X} \mathcal{L}^{\otimes n}$
이다. 이는 모든 $n \in \mathbf{Z}$에 대해 준연접 $\mathcal{O}_X$-가군의
등식이다.
\item $\Omega^\bullet_{X/S} = \Omega^\bullet_{L/X, 0}$
가 복합체의 등식으로서 성립한다.
\item $n > 0$ 및 $p \geq 0$이면
$\Omega^p_{L/X, n} = \Omega^p_{L^\star/S, n}$.
\end{enumerate}
\end{lemma}

\begin{proof}
각 경우에 전역적으로 정의된 표준 사상이 있고, 생략하는 국소 계산에 의해
그 사상은 동형이다.
\end{proof}

\begin{lemma}
\label{derham-lemma-line-bundle-characteristic-zero}
위의 상황에서 사상 $S \to \Spec(\mathbf{Q})$가 존재한다고 가정하자.
그러면 $\Omega^\bullet_{X/S} \to \pi_*\Omega^\bullet_{L/S}$는
준동형이고 $H_{dR}^*(X/S) = H_{dR}^*(L/S)$이다.
\end{lemma}

\begin{proof}
$R$을 $\mathbf{Q}$-대수라 하고, $A$를 $R$-대수라 하자. 아핀 국소 명제는 사상
$$
\Omega^\bullet_{A/R} \longrightarrow \Omega^\bullet_{A[t]/R}
$$
이 $R$-가군의 복합체들의 준동형이라는 것이다. 사실 이는 호모토피 동치이고,
그 호모토피 역은 $g \omega + g' \text{d}t \wedge \omega'$를
$g(0)\omega$로 보내는 사상으로 주어진다. 여기서
$g, g' \in A[t]$이고 $\omega, \omega' \in \Omega^\bullet_{A/R}$이다.
호모토피는
$g \omega + g' \text{d}t \wedge \omega'$를 $(\int g') \omega'$로 보낸다.
여기서 $\int g' \in A[t]$는 상수항이 영이고 $t$에 관한 미분이 $g'$인
다항식이다. 물론 여기서는 $R$이 $\mathbf{Q}$를 포함한다는 사실을
사용한다. 실제로 $\int t^n = (1/n)t^{n + 1}$이다.
\end{proof}

\begin{example}
\label{derham-example-affine-line}
양의 표수에서는 보조정리 \ref{derham-lemma-line-bundle-characteristic-zero}가
거짓이다. $\mathbf{A}^1_k = \Spec(k[x])$의 체 $k$ 위 드람 복합체는 직합
$$
k \oplus
\bigoplus\nolimits_{n \geq 1}
(k \cdot t^n \xrightarrow{n}
k \cdot t^{n - 1} \text{d}t)
$$
처럼 보인다. 따라서 $k$의 표수가 $p > 0$이면
$H^0_{dR}(\mathbf{A}^1_k/k)$와
$H^1_{dR}(\mathbf{A}^1_k/k)$
는 둘 다 $k$ 위에서 무한 차원이다.
\end{example}








\section{사영공간의 드람 코호몰로지}
\label{derham-section-projective-space}

\noindent
$A$를 환이라 하고 $n \geq 1$이라 하자. 구조 사상
$\mathbf{P}^n_A \to \Spec(A)$은 고유하고 매끄러우며 상대차원이
$n$이다. 이 사상이 상대차원 $n$으로 매끄럽고 유한형인 것은
$\mathbf{P}^n_A$가 각각 $\mathbf{A}^n_A$와 동형인 스킴들로 이루어진
유한 아핀 열린 덮개를 갖기 때문이다. 스킴의 구성의 보조정리
\ref{constructions-lemma-standard-covering-projective-space}를 보라.
또한 이 사상은 분리이고 보편적으로 닫혀 있으므로 고유하다
(스킴의 구성의 보조정리 \ref{constructions-lemma-projective-space-separated}).
$\mathcal{O}$와 $\mathcal{O}(d)$로 각각 구조층
$\mathcal{O}_{\mathbf{P}^n_A}$와 세르 꼬임
$\mathcal{O}_{\mathbf{P}^n_A}(d)$를 나타내자.
상대 미분층을 $\Omega = \Omega_{\mathbf{P}^n_A/A}$로 나타내고,
그 외대수의 p차 성분을 $\Omega^p$로 나타내자.

\begin{lemma}
\label{derham-lemma-euler-sequence}
다음 짧은 완전열이 존재한다.
$$
0 \to \Omega \to \mathcal{O}(-1)^{\oplus n + 1} \to \mathcal{O} \to 0
$$
\end{lemma}

\begin{proof}
이를 설명하기 위해
$\mathbf{P}^n_A = \text{Proj}(A[T_0, \ldots, T_n])$,
임을 상기하고, 기호상 다음과 같이 쓰자.
$$
\mathcal{O}(-1)^{\oplus n + 1} =
\bigoplus\nolimits_{j = 0, \ldots, n} \mathcal{O}(-1) \text{d}T_j
$$
위 짧은 완전열의 첫 번째 화살표
$$
\Omega \to
\bigoplus\nolimits_{j = 0, \ldots, n} \mathcal{O}(-1) \text{d}T_j
$$
는 위에서 언급한 각 표준 열린집합
$D_+(T_i) = \Spec(A[T_0/T_i, \ldots, T_n/T_i])$
위에서 다음 규칙으로 주어진다.
$$
\sum\nolimits_{j \not = i} g_j \text{d}(T_j/T_i)
\longmapsto
\sum\nolimits_{j \not = i} g_j/T_i \text{d}T_j
- (\sum\nolimits_{j \not = i} g_jT_j/T_i^2) \text{d}T_i
$$
$1/T_i$가 $\mathcal{O}(-1)$의 $D_+(T_i)$ 위 단면이므로 이는
의미가 있다. 사상
$$
\bigoplus\nolimits_{j = 0, \ldots, n} \mathcal{O}(-1) \text{d}T_j
\to
\mathcal{O}
$$
는 $\text{d}T_j$를 $T_j$로 보내어 주어진다. 더 정확히 말해
$D_+(T_i)$ 위에서 단면 $\sum g_j \text{d}T_j$를
$\sum T_jg_j$로 보낸다. 이것이 짧은 완전열을 준다는 확인은 생략한다.
\end{proof}

\noindent
정수 $k \in \mathbf{Z}$와 준연접
$\mathcal{O}_{\mathbf{P}^n_A}$-가군 $\mathcal{F}$가 주어졌을 때,
통상대로 $\mathcal{F}(k)$로 나타내자. 이는 $k$번째 세르 꼬임을
$\mathcal{F}$에 취한 것이다. 스킴의 구성의 정의
\ref{constructions-definition-twist}를 보라.

\begin{lemma}
\label{derham-lemma-twisted-hodge-cohomology-projective-space}
위 상황에서 코호몰로지 군들은 다음 성질을 갖는다.
\begin{enumerate}
\item $H^q(\mathbf{P}^n_A, \Omega^p) = 0$
이며, 유일한 예외는 $0 \leq p = q \leq n$인 경우이다.
\item $0 \leq p \leq n$이면 $A$-가군
$H^p(\mathbf{P}^n_A, \Omega^p)$는 계수 $1$인 자유 가군이다.
\item $q > 0$, $k > 0$이고 $p$가 임의이면
$H^q(\mathbf{P}^n_A, \Omega^p(k)) = 0$이다.
\item 여기에 더 추가한다.
\end{enumerate}
\end{lemma}

\begin{proof}
이하에서는 스킴의 코호몰로지의 보조정리
\ref{coherent-lemma-cohomology-projective-space-over-ring}
의 결과를 더 언급하지 않고 사용한다. 특히 이 명제들은
$H^q(\mathbf{P}^n_A, \mathcal{O}(k))$에 대해 참이다.

\medskip\noindent
$p = 1$인 경우를 증명하자. 짧은 완전열
$$
0 \to \Omega \to \mathcal{O}(-1)^{\oplus n + 1} \to \mathcal{O} \to 0
$$
을 생각하자. 이는 보조정리 \ref{derham-lemma-euler-sequence}의 완전열이다.
$\mathcal{O}(-1)$의 코호몰로지가 모든 차수에서 소멸하므로,
$H^q(\mathbf{P}^n_A, \Omega)$는 차수 $1$을 제외하면 영이고,
그 차수에서는 $1$의 경계에 의해 자유롭게 생성된다. 이 원소는
$H^0(\mathbf{P}^n_A, \mathcal{O})$에서 취한다.

\medskip\noindent
$p > 1$이라고 가정하자. 위 짧은 완전열을 $2$단계 여과를
$\mathcal{O}(-1)^{\oplus n + 1}$ 위에 정의하는 것으로
생각하자. $\wedge^p\mathcal{O}(-1)^{\oplus n + 1}$ 위에 유도되는 여과는
다음과 같다.
$$
0 \to \Omega^p \to \wedge^p\left(\mathcal{O}(-1)^{\oplus n + 1}\right)
\to \Omega^{p - 1} \to 0
$$
$\wedge^p\mathcal{O}(-1)^{\oplus n + 1}$은 $n + 1$개에서 $p$개를
고르는 수만큼의 $\mathcal{O}(-p)$들의 직합과 동형이고, 따라서 모든
차수에서 코호몰로지가 소멸한다. 귀납 가정에 의해
$H^q(\mathbf{P}^n_A, \Omega^p)$는 $q = p$가 아니면 영이고,
$H^p(\mathbf{P}^n_A, \Omega^p)$는 계수 $1$인 자유 가군이며 그 생성원은
$H^{p - 1}(\mathbf{P}^n_A, \Omega^{p - 1})$의 생성원의 경계이다.

\medskip\noindent
$k > 0$이라 하자. 예를 들어 $\Omega^n = \mathcal{O}(-n - 1)$이고,
이는 위 짧은 완전열에서 $p = n + 1$로 놓으면 알 수 있다. 따라서 $\Omega^n(k)$의
코호몰로지는 양의 차수에서 소멸한다. 짧은 완전열
$$
0 \to \Omega^p(k) \to \wedge^p\left(\mathcal{O}(-1)^{\oplus n + 1}\right)(k)
\to \Omega^{p - 1}(k) \to 0
$$
과 $p$에 관한 {\it 내림} 귀납법을 사용하면 $\Omega^p(k)$의
코호몰로지가 양의 차수에서 소멸함을 얻으며, 이는 모든 $p$에 대해 성립한다.
\end{proof}

\begin{lemma}
\label{derham-lemma-hodge-cohomology-projective-space}
$H^q(\mathbf{P}^n_A, \Omega^p) = 0$이며, 유일한 예외는
$0 \leq p = q \leq n$인 경우이다. $0 \leq p \leq n$이면 $A$-가군
$H^p(\mathbf{P}^n_A, \Omega^p)$는 계수 $1$인 자유 가군이고, 기저원은
$c_1^{Hodge}(\mathcal{O}(1))^p$이다.
\end{lemma}

\begin{proof}
소멸성과 자유성은 보조정리
\ref{derham-lemma-twisted-hodge-cohomology-projective-space}에서 따른다.
$p = 0$이면 $1 \in H^0(\mathbf{P}^n_A, \mathcal{O})$가 생성원임은
분명하다.

\medskip\noindent
$p = 1$인 경우를 증명하자. 짧은 완전열
$$
0 \to \Omega \to \mathcal{O}(-1)^{\oplus n + 1} \to \mathcal{O} \to 0
$$
을 생각하자. 이는 보조정리 \ref{derham-lemma-euler-sequence}의 완전열이다. 보조정리
\ref{derham-lemma-twisted-hodge-cohomology-projective-space}의 증명에서
$H^1(\mathbf{P}^n_A, \Omega)$의 생성원이 경계 $\xi$이고, 이것은
$1 \in H^0(\mathbf{P}^n_A, \mathcal{O})$에서 나온 것임을 보았다.
보조정리 \ref{derham-lemma-euler-sequence}의 증명에서와 같이
$\mathcal{O}(-1)^{\oplus n + 1}$을
$\bigoplus_{j = 0, \ldots, n} \mathcal{O}(-1)\text{d}T_j$와 동일시하자.
열린 덮개
$$
\mathcal{U} : 
\mathbf{P}^n_A =
\bigcup\nolimits_{i = 0, \ldots, n} D_{+}(T_i)
$$
를 생각하자. 전역 단면 $1$은 $\mathcal{O}$의 단면이다. 이를
$U_i = D_+(T_i)$에 제한한 것은 단면 $T_i^{-1} \text{d}T_i$로
올릴 수 있다. 이 단면은 $\bigoplus \mathcal{O}(-1)\text{d}T_j$의
$U_i$ 위 단면이다. 따라서 $\xi$를 나타내는
코사이클은 다음과 같다.
$$
T_{i_1}^{-1} \text{d}T_{i_1} - T_{i_0}^{-1} \text{d}T_{i_0} =
\text{d}\log(T_{i_1}/T_{i_0}) \in \Omega(U_{i_0i_1})
$$
한편 각 $i$에 대해 단면 $T_i$는 $\mathcal{O}(1)$의
$U_i$ 위 자명화 단면이다. 따라서
$f_{i_0i_1} = T_{i_1}/T_{i_0} \in \mathcal{O}^*(U_{i_0i_1})$는
$\mathcal{O}(1)$을 $\Pic(\mathbf{P}^n_A)$에서 나타내는 코사이클이다.
절 \ref{derham-section-first-chern-class}을 보라. 그러므로
$c_1^{Hodge}(\mathcal{O}(1))$은 코사이클
$\text{d}\log(T_{i_1}/T_{i_0})$로 주어지며, 이는 위에서 $\xi$에 대해
얻은 것과 일치한다.

\medskip\noindent
일반적인 $p$에 대해서는 귀납법으로 증명한다. 기초 단계 $p = 0, 1$은
위에서 다루었다. $p > 1$이라 가정하자. 보조정리
\ref{derham-lemma-twisted-hodge-cohomology-projective-space}의 증명에서
$H^p(\mathbf{P}^n_A, \Omega^p)$의 생성원은 다음 짧은 완전열에 수반되는
코호몰로지 긴 완전열에서 $c_1^{Hodge}(\mathcal{O}(1))^{p - 1}$의
경계임을 보았다.
$$
0 \to \Omega^p \to \wedge^p\left(\mathcal{O}(-1)^{\oplus n + 1}\right)
\to \Omega^{p - 1} \to 0
$$
절 \ref{derham-section-first-chern-class}의 계산에 의해 코호몰로지 류
$c_1^{Hodge}(\mathcal{O}(1))^{p - 1}$은 부호를 제외하면 다음 항들을
갖는 코사이클로 표현된다.
$$
\beta_{i_0 \ldots i_{p - 1}} =
\text{d}\log(T_{i_1}/T_{i_0}) \wedge
\text{d}\log(T_{i_2}/T_{i_1}) \wedge \ldots \wedge
\text{d}\log(T_{i_{p - 1}}/T_{i_{p - 2}})
$$
이는 $\Omega^{p - 1}(U_{i_0 \ldots i_{p - 1}})$의 원소이다. 이들
$\beta_{i_0 \ldots i_{p - 1}}$는 단면
$\tilde \beta_{i_0 \ldots i_{p -1}} =
T_{i_0}^{-1}\text{d}T_{i_0} \wedge \beta_{i_0 \ldots i_{p - 1}}$로
올릴 수 있다. 이는 $\wedge^p(\bigoplus \mathcal{O}(-1) \text{d}T_j)$의
$U_{i_0 \ldots i_{p - 1}}$ 위 단면이다. 따라서
$H^p(\mathbf{P}^n_A, \Omega^p)$의 생성원은 성분이 다음과 같은
코사이클로 주어진다.
\begin{align*}
\sum\nolimits_{a = 0}^p (-1)^a
\tilde \beta_{i_0 \ldots \hat{i_a} \ldots i_p}
& =
T_{i_1}^{-1}\text{d}T_{i_1} \wedge \beta_{i_1 \ldots i_p}
+ \sum\nolimits_{a = 1}^p (-1)^a
T_{i_0}^{-1}\text{d}T_{i_0} \wedge
\beta_{i_0 \ldots \hat{i_a} \ldots i_p} \\
& =
(T_{i_1}^{-1}\text{d}T_{i_1} - T_{i_0}^{-1}\text{d}T_{i_0}) \wedge
\beta_{i_1 \ldots i_p} +
T_{i_0}^{-1}\text{d}T_{i_0} \wedge \text{d}(\beta)_{i_0 \ldots i_p} \\
& =
\text{d}\log(T_{i_1}/T_{i_0}) \wedge \beta_{i_1 \ldots i_p}
\end{align*}
이를 $\Omega^p$의 $U_{i_0 \ldots i_p}$ 위 단면으로 본다.
이는 부호를 제외하면 $c_1^{Hodge}(\mathcal{O}(1))^p$를 나타내는
코사이클과 같으므로 증명이 끝난다.
\end{proof}

\begin{lemma}
\label{derham-lemma-de-rham-cohomology-projective-space}
$0 \leq i \leq n$이면 드람 코호몰로지
$H^{2i}_{dR}(\mathbf{P}^n_A/A)$는 자유 $A$-가군이고 계수가 $1$이며,
기저원은 $c_1^{dR}(\mathcal{O}(1))^i$이다. 그 밖의 모든 차수에서
$\mathbf{P}^n_A$의 $A$ 위 드람 코호몰로지는 영이다.
\end{lemma}

\begin{proof}
절 \ref{derham-section-hodge-to-de-rham}의 호지--드람 스펙트럼 수열을
생각하자. 보조정리 \ref{derham-lemma-hodge-cohomology-projective-space}에서 계산한
$\mathbf{P}^n_A$의 $A$ 위 호지 코호몰로지에 의해, 이 스펙트럼 수열은
$E_1$ 쪽에서 퇴화한다. 따라서 $H^{2i}_{dR}(\mathbf{P}^n_A/A)$는 자유
$A$-가군이고 계수가 $1$이다. 이는 $0 \leq i \leq n$일 때 성립하고,
그 밖에는 영이다.
$c_1^{dR}(\mathcal{O}(1))^i \in H^{2i}_{dR}(\mathbf{P}^n_A/A)$이며
$i = 0, \ldots, n$임을 주목하자. $i = n$일 때 이 원소는 복합체의 사상
아래에서 $c_1^{Hodge}(\mathcal{L})^n$의 상이다.
$$
\Omega^n_{\mathbf{P}^n_A/A}[-n]
\longrightarrow
\Omega^\bullet_{\mathbf{P}^n_A/A}
$$
이는 예를 들어 비고 \ref{derham-remark-truncations}의 논의 또는 절
\ref{derham-section-first-chern-class}에서 이 류들을 나타내는 코사이클의
명시적 기술로부터 따른다. 스펙트럼 수열에 의하면 유도된 사상
$$
H^n(\mathbf{P}^n_A, \Omega^n_{\mathbf{P}^n_A/A}) \longrightarrow
H^{2n}_{dR}(\mathbf{P}^n_A/A)
$$
은 동형이다. 또한 $c_1^{Hodge}(\mathcal{L})^n$이 그 정의역의 생성원이므로
(보조정리 \ref{derham-lemma-hodge-cohomology-projective-space}),
$c_1^{dR}(\mathcal{L})^n$이 공역의 생성원임을 알 수 있다. 컵곱의
$A$-쌍선형성에 의해 $c_1^{dR}(\mathcal{L})^i$도
$H^{2i}_{dR}(\mathbf{P}^n_A/A)$의 생성원이며, 이는
$0 \leq i \leq n$에 대해 성립한다.
\end{proof}









\section{매끄러운 사상에 대한 스펙트럼 수열}
\label{derham-section-relative-spectral-sequence}

\noindent
스킴들의 가환 도식
$$
\xymatrix{
X \ar[rr]_f \ar[rd]_p & & Y \ar[ld]^q \\
& S
}
$$
을 생각하자. 여기서 $f$는 매끄러운 사상이다. 그러면 국소적으로 분할되는
짧은 완전열
$$
0 \to f^*\Omega_{Y/S} \to \Omega_{X/S} \to \Omega_{X/Y} \to 0
$$
을 얻는다(스킴의 사상의 보조정리
\ref{morphisms-lemma-triangle-differentials-smooth}). 이를 내림 여과 $F$로
생각하자. 이 여과는 $\Omega_{X/S}$ 위에 놓이며, 여기서
$F^0\Omega_{X/S} = \Omega_{X/S}$, $F^1\Omega_{X/S} = f^*\Omega_{Y/S}$이고
$F^2\Omega_{X/S} = 0$이다. 함자 $\wedge^p$를 적용하면 각 $p$에 대해
유도된 여과
$$
\Omega^p_{X/S} = F^0\Omega^p_{X/S} \supset
F^1\Omega^p_{X/S} \supset
F^2\Omega^p_{X/S} \supset \ldots \supset F^{p + 1}\Omega^p_{X/S} = 0
$$
를 얻으며, 그 연속 몫은
$$
\text{gr}^k\Omega^p_{X/S} =
F^k\Omega^p_{X/S}/F^{k + 1}\Omega^p_{X/S} =
f^*\Omega^k_{Y/S} \otimes_{\mathcal{O}_X} \Omega^{p - k}_{X/Y} =
f^{-1}\Omega^k_{Y/S} \otimes_{f^{-1}\mathcal{O}_Y} \Omega^{p - k}_{X/Y}
$$
$k = 0, \ldots, p$에 대해 위와 같다. 실제로 라이프니츠 법칙을 사용하면
$F^k\Omega^\bullet_{X/S}$가 $\Omega^\bullet_{X/S}$의 부분복합체임을
확인할 수 있다. 이로써 $\Omega^\bullet_{X/S}$는 여과 복합체의 구조를
갖는다. 이는 다음 식을 관찰하여도 알 수 있다.
$$
F^k\Omega^\bullet_{X/S} =
\Im\left(\wedge :
\text{Tot}(
f^{-1}\sigma_{\geq k}\Omega^\bullet_{Y/S} \otimes_{p^{-1}\mathcal{O}_S}
\Omega^\bullet_{X/S})
\longrightarrow
\Omega^\bullet_{X/S}\right)
$$
이것은 $X$ 위 복합체 사상의 상이다. 여과 복합체
$$
\Omega^\bullet_{X/S} = F^0\Omega^\bullet_{X/S} \supset
F^1\Omega^\bullet_{X/S} \supset F^2\Omega^\bullet_{X/S} \supset \ldots
$$
의 수반 등급 성분은 다음과 같다.
$$
\text{gr}^k\Omega^\bullet_{X/S} = 
f^{-1}\Omega^k_{Y/S}[-k] \otimes_{f^{-1}\mathcal{O}_Y} \Omega^\bullet_{X/Y}
$$
이는 위의 논의에서 따른다.

\begin{lemma}
\label{derham-lemma-spectral-sequence-smooth}
$f : X \to Y$를 바탕 스킴 $S$ 위에서 준콤팩트하고 준분리이며 매끄러운
스킴 사상이라 하자. 첫째 쪽이
$$
E_1^{p, q} =
H^q(\Omega^p_{Y/S} \otimes_{\mathcal{O}_Y}^\mathbf{L} Rf_*\Omega^\bullet_{X/Y})
$$
이고 $R^{p + q}f_*\Omega^\bullet_{X/S}$로 수렴하는 유계 스펙트럼 수열이
존재한다.
\end{lemma}

\begin{proof}
위에서 도입한 여과를 갖춘 $\Omega^\bullet_{X/S}$를 여과 복합체로 생각하자.
이 스펙트럼 수열은 층 코호몰로지의 보조정리
\ref{cohomology-lemma-relative-spectral-sequence-filtered-object}에서 주어진 것이다.
스킴의 유도 범주의 보조정리
\ref{perfect-lemma-cohomology-de-rham-base-change}에 의해
$$
Rf_*\text{gr}^k\Omega^\bullet_{X/S} =
\Omega^k_{Y/S}[-k] \otimes_{\mathcal{O}_Y}^\mathbf{L} Rf_*\Omega^\bullet_{X/Y}
$$
이고, 이로써 결론을 얻는다.
\end{proof}

\begin{remark}
\label{derham-remark-gauss-manin}
보조정리 \ref{derham-lemma-spectral-sequence-smooth}에서 코호몰로지 층
$$
\mathcal{H}^q_{dR}(X/Y) = H^q(Rf_*\Omega^\bullet_{X/Y})
$$
을 생각하자. $f$가 매끄러울 뿐 아니라 고유하고 $S$가
$\mathbf{Q}$ 위의 스킴이면 $\mathcal{H}^q_{dR}(X/Y)$는 유한 국소 자유이다
(향후 참고문헌을 여기에 삽입한다). $\mathcal{H}^q_{dR}(X/Y)$가 평탄한
$\mathcal{O}_Y$-가군이라고만 가정해도 다음을 얻는다(짧은 논증은 생략한다).
$$
E_1^{p, q} =
\Omega^p_{Y/S} \otimes_{\mathcal{O}_Y} \mathcal{H}^q_{dR}(X/Y)
$$
그리고 스펙트럼 수열의 미분들은 사상
$$
d_1^{p, q} :
\Omega^p_{Y/S} \otimes_{\mathcal{O}_Y} \mathcal{H}^q_{dR}(X/Y)
\longrightarrow
\Omega^{p + 1}_{Y/S} \otimes_{\mathcal{O}_Y} \mathcal{H}^q_{dR}(X/Y)
$$
이다. 특히 $p = 0$이면 사상
$d_1^{0, q} : \mathcal{H}^q_{dR}(X/Y) \to
\Omega^1_{Y/S} \otimes_{\mathcal{O}_Y} \mathcal{H}^q_{dR}(X/Y)$
을 얻는데, 이는 적분가능 접속 $\nabla$가 된다
(향후 참고문헌을 여기에 삽입한다). 그리고 복합체
$$
\mathcal{H}^q_{dR}(X/Y) \to
\Omega^1_{Y/S} \otimes_{\mathcal{O}_Y} \mathcal{H}^q_{dR}(X/Y) \to
\Omega^2_{Y/S} \otimes_{\mathcal{O}_Y} \mathcal{H}^q_{dR}(X/Y) \to \ldots
$$
에서 미분이 $d_1^{\bullet, q}$로 주어질 때, 이는 $\nabla$의 드람 복합체이다.
접속 $\nabla$를 {\it 가우스--마닌 접속}이라 한다.
\end{remark}






\section{르레--히르슈형 정리}
\label{derham-section-leray-hirsch}

\noindent
이 절에서는 매끄럽고 고유한 사상에 대해 위쪽의 드람 코호몰로지를
아래쪽의 드람 코호몰로지로 나타낼 수 있는 경우가 있음을 증명한다.

\begin{lemma}
\label{derham-lemma-relative-global-generation-on-fibres}
$f : X \to Y$를 매끄럽고 고유한 스킴 사상이라 하자.
$N$과 $n_1, \ldots, n_N \geq 0$을 정수라 하고
$\xi_i \in H^{n_i}_{dR}(X/Y)$, $1 \leq i \leq N$이라 하자.
모든 점 $y \in Y$에 대해 $\xi_1, \ldots, \xi_N$의 상들이
$H^*_{dR}(X_y/y)$의 $\kappa(y)$ 위 기저를 이룬다고 가정하자. 그러면 사상
$$
\bigoplus\nolimits_{i = 1}^N \mathcal{O}_Y[-n_i]
\longrightarrow
Rf_*\Omega^\bullet_{X/Y}
$$
은 $\xi_1, \ldots, \xi_N$에 결부된 사상이며 동형이다.
\end{lemma}

\begin{proof}
보조정리 \ref{derham-lemma-proper-smooth-de-Rham}에 의해
$Rf_*\Omega^\bullet_{X/Y}$는 $D(\mathcal{O}_Y)$의 완전 대상이며, 그 구성은
임의의 기저변환과 가환한다. 따라서 보조정리의 사상 $a : K \to L$은
$D(\mathcal{O}_Y)$의 완전 대상들 사이의 사상이고,
올에 대한 가정에 의해 임의의 점으로의 도출 제한은 동형이다. 그러면 원뿔
$C$, 즉 $a$의 원뿔은 $D(\mathcal{O}_Y)$의 완전 대상이며(층 코호몰로지의 보조정리
\ref{cohomology-lemma-two-out-of-three-perfect}), 임의의 점으로의 도출 제한은
영이다. 따라서 대수학 심화의 보조정리
\ref{more-algebra-lemma-lift-perfect-from-residue-field}에 의해 $C$는 영이고
$a$는 동형이다. (이를 대수로 옮기는 데에는 스킴의 유도 범주의
보조정리들 \ref{perfect-lemma-affine-compare-bounded}와
\ref{perfect-lemma-perfect-affine}도 사용한다.)
\end{proof}

\noindent
먼저 이 절의 주결과를 다음 특수한 경우에 증명한다.

\begin{lemma}
\label{derham-lemma-global-generation-on-fibres}
$f : X \to Y$를 바탕 $S$ 위에서 매끄럽고 고유한 스킴 사상이라 하자.
다음을 가정하자.
\begin{enumerate}
\item $Y$와 $S$는 아핀이다.
\item 정수 $N$과 $n_1, \ldots, n_N \geq 0$ 및
$\xi_i \in H^{n_i}_{dR}(X/S)$, $1 \leq i \leq N$이 존재하여,
모든 점 $y \in Y$에 대해 $\xi_1, \ldots, \xi_N$의 상들이
$H^*_{dR}(X_y/y)$의 $\kappa(y)$ 위 기저를 이룬다.
\end{enumerate}
그러면 사상
$$
\bigoplus\nolimits_{i = 1}^N H^*_{dR}(Y/S) \longrightarrow
H^*_{dR}(X/S), \quad
(a_1, \ldots, a_N) \longmapsto  \sum \xi_i \cup f^*a_i
$$
은 동형이다.
\end{lemma}

\begin{proof}
$Y = \Spec(A)$이고 $S = \Spec(R)$라 하자. 이 경우
$\Omega^\bullet_{A/R}$은 보조정리 \ref{derham-lemma-de-rham-affine}에 의해
$R\Gamma(Y, \Omega^\bullet_{Y/S})$를 계산한다. 유한 아핀 열린 덮개
$\mathcal{U} : X = \bigcup_{i \in I} U_i$를 택하자. 복합체
$$
K^\bullet =
\text{Tot}(\check{\mathcal{C}}^\bullet(\mathcal{U}, \Omega_{X/S}^\bullet))
$$
를 생각하자. 이는 층 코호몰로지의 절
\ref{cohomology-section-cech-cohomology-of-complexes}에 나오는 복합체이다.
이 복합체에 관한 몇 가지 사실을 모아 적자. 그 대부분은 방금 인용한 곳에서
찾을 수 있다.
\begin{enumerate}
\item $K^\bullet$은 $R$-가군의 복합체이고 그 항들은 $A$-가군이다.
\item $K^\bullet$은 $R\Gamma(X, \Omega^\bullet_{X/S})$를 $D(R)$에서 나타낸다
(스킴의 코호몰로지의 보조정리
\ref{coherent-lemma-quasi-coherent-affine-cohomology-zero} 및
층 코호몰로지의 보조정리 \ref{cohomology-lemma-cech-complex-complex-computes}).
\item 자연스러운 사상 $\Omega^\bullet_{A/R} \to K^\bullet$이 존재한다.
이는 $R$-가군 복합체의 사상이고 각 항에서
$A$-선형이며 코호몰로지 위에서 당김 사상
$H^*_{dR}(Y/S) \to H^*_{dR}(X/S)$을 유도한다.
\item $K^\bullet$에는 $\wedge$로 나타내는 곱셈이 있고, 이 곱셈은 이를
미분 등급 $R$-대수로 만든다.
\item $K^\bullet$ 위의 곱셈은 $H^*_{dR}(X/S)$ 위의 컵곱을 유도한다
(층 코호몰로지의 절 \ref{cohomology-section-cup-product}).
\item 여과 $F$, 즉 $\Omega^*_{X/S}$ 위의 여과는 다음 여과를 유도한다.
$$
K^\bullet =
F^0K^\bullet \supset F^1K^\bullet \supset F^2K^\bullet \supset \ldots
$$
이는 $K^\bullet$의 부분복합체들로 이루어지고 다음 성질을 갖는다.
\begin{enumerate}
\item $F^kK^n \subset K^n$은 $A$-부분가군이다.
\item $F^kK^\bullet \wedge F^lK^\bullet \subset F^{k + l}K^\bullet$이다.
\item $\text{gr}^kK^\bullet$은 $A$-가군의 복합체이다.
\item $\text{gr}^0K^\bullet =
\text{Tot}(\check{\mathcal{C}}^\bullet(\mathcal{U}, \Omega_{X/Y}^\bullet))$
이고, 이는 $R\Gamma(X, \Omega^\bullet_{X/Y})$를 $D(A)$에서 나타낸다.
\item 곱셈은 동형
$\Omega^k_{A/R}[-k] \otimes_A \text{gr}^0K^\bullet \to \text{gr}^kK^\bullet$
을 유도한다.
\end{enumerate}
\end{enumerate}
이 명제들의 상세한 증명은 생략한다. 보조정리
\ref{derham-lemma-spectral-sequence-smooth}에서 스펙트럼 수열을 구성하기까지의
논의를 보라.

\medskip\noindent
각 $i = 1, \ldots, N$에 대해 코사이클 $x_i \in K^{n_i}$를 택하자.
이 코사이클은 $\xi_i$를 나타낸다. 다음으로 복합체의 사상
$$
\tilde x :
M^\bullet = \bigoplus\nolimits_{i = 1, \ldots, N}
\Omega^\bullet_{A/R}[-n_i]
\longrightarrow
K^\bullet
$$
을 생각하자. 이 사상은 $\omega$를 $i$번째 직합 성분에서
$x_i \wedge \omega$로 보낸다. 이제 이 사상이 준동형임을 보이면 된다.
$M^\bullet$에 다음 규칙으로 여과 복합체의 구조를 준다.
$$
F^kM^\bullet =
\bigoplus\nolimits_{i = 1, \ldots, N}
(\sigma_{\geq k}\Omega^\bullet_{A/R})[-n_i]
$$
이 선택 아래 사상 $\tilde x$는 여과 복합체의 사상이다.
$\text{gr}^0M^\bullet = \bigoplus A[-n_i]$임을 주목하자. 또한 곱셈은 동형
$\Omega^k_{A/R}[-k] \otimes_A \text{gr}^0M^\bullet \to \text{gr}^kM^\bullet$.
을 유도한다. 구성과 보조정리
\ref{derham-lemma-relative-global-generation-on-fibres}에 의해
$$
\text{gr}^0\tilde x :
\text{gr}^0M^\bullet \longrightarrow
\text{gr}^0K^\bullet
$$
는 $D(A)$에서 동형이다. 따라서 모든 $k \geq 0$에 대해 동형
$$
\text{gr}^k \tilde x :
\text{gr}^kM^\bullet = \Omega^k_{A/R}[-k] \otimes_A \text{gr}^0M^\bullet
\longrightarrow
\Omega^k_{A/R}[-k] \otimes_A \text{gr}^0K^\bullet =
\text{gr}^kK^\bullet
$$
을 $D(A)$에서 얻는다. 실제로 복합체
$\text{gr}^0K^\bullet =
\text{Tot}(\check{\mathcal{C}}^\bullet(\mathcal{U}, \Omega_{X/Y}^\bullet))$
는 스킴의 유도 범주의 보조정리
\ref{perfect-lemma-K-flat}에 의해 $A$-가군의 복합체로서 K-평탄하다.
따라서 우변의 텐서곱은 도출 텐서곱이며, 좌변도 직접 확인하면 마찬가지이다.
마지막으로 도출 텐서곱을 취하는 연산
$\Omega^k_{A/R}[-k] \otimes_A^\mathbf{L} -$는 $D(A)$ 위의 함자이므로
동형을 동형으로 보낸다. $k$에 관한 귀납법으로
$$
\tilde x : M^\bullet/F^kM^\bullet \to K^\bullet/F^kK^\bullet
$$
가 $D(R)$에서 동형임을 얻는다. 실제로 짧은 완전열
$$
0 \to F^kM^\bullet/F^{k + 1}M^\bullet \to
M^\bullet/F^{k + 1}M^\bullet \to
\text{gr}^kM^\bullet \to 0
$$
을 가지며 $K^\bullet$에 대해서도 마찬가지이다. 각 주어진 차수에서 여과가
유한하므로, 이로부터 $\tilde x$가 준동형임이 따른다.
\end{proof}

\begin{proposition}
\label{derham-proposition-global-generation-on-fibres}
$f : X \to Y$를 바탕 $S$ 위에서 매끄럽고 고유한 스킴 사상이라 하자.
$N$과 $n_1, \ldots, n_N \geq 0$을 정수라 하고
$\xi_i \in H^{n_i}_{dR}(X/S)$, $1 \leq i \leq N$이라 하자.
모든 점 $y \in Y$에 대해 $\xi_1, \ldots, \xi_N$의 상들이
$H^*_{dR}(X_y/y)$의 $\kappa(y)$ 위 기저를 이룬다고 가정하자. 사상
$$
\tilde \xi = \bigoplus \tilde \xi_i[-n_i] :
\bigoplus \Omega^\bullet_{Y/S}[-n_i]
\longrightarrow
Rf_*\Omega^\bullet_{X/S}
$$
은 $D(Y, (Y \to S)^{-1}\mathcal{O}_S)$에서 동형이며(증명을 보라), 이에
대응하여 사상
$$
\bigoplus\nolimits_{i = 1}^N H^*_{dR}(Y/S) \longrightarrow
H^*_{dR}(X/S), \quad
(a_1, \ldots, a_N) \longmapsto  \sum \xi_i \cup f^*a_i
$$
도 동형이다.
\end{proposition}

\begin{proof}
구조 사상을 $p : X \to S$와 $q : Y \to S$로 나타내자.
$\xi'_i : \Omega^\bullet_{X/S} \to \Omega^\bullet_{X/S}[n_i]$를 비고
\ref{derham-remark-cup-product-as-a-map}에서 $\xi_i$에 대응하는 사상이라 하자.
다음 사상을
$$
\tilde \xi_i :
\Omega^\bullet_{Y/S} \to Rf_*\Omega^\bullet_{X/S}[n_i]
$$
$\xi'_i$와 표준 사상
$\Omega^\bullet_{Y/S} \to Rf_*\Omega^\bullet_{X/S}$의 합성으로 나타내자.
등식
$$
R\Gamma(Y, Rf_*\Omega^\bullet_{X/S}) = R\Gamma(X, \Omega^\bullet_{X/S})
$$
을 사용하면 코호몰로지 위에서 $\tilde \xi_i$는 사상
$\eta \mapsto \xi_i \cup f^*\eta$이며, 이는 $H^m_{dR}(Y/S)$에서
$H^{m + n}_{dR}(X/S)$로 간다. 또한 $\xi'_i$의 구성이 열린집합으로의
제한과 가환하므로, $\tilde \xi_i$의 구성도 열린집합으로의 제한과 가환한다.

\medskip\noindent
따라서 사상
$$
\tilde \xi = \bigoplus \tilde \xi_i[-n_i] :
\bigoplus \Omega^\bullet_{Y/S}[-n_i]
\longrightarrow
Rf_*\Omega^\bullet_{X/S}
$$
을 생각할 수 있다. 보조정리를 증명하려면 이것이
$D(Y, q^{-1}\mathcal{O}_S)$에서 동형임을 보이면 충분하다. $\tilde \xi$가
적절한 여과를 갖춘 여과 복합체 사상에서 온다는 것을 보일 수 있다면,
보조정리 \ref{derham-lemma-spectral-sequence-smooth}의 스펙트럼 수열을 사용하여
증명을 끝낼 수 있다. 그러나 이는 필요한 것보다 더 많은 작업이므로, 대신
아핀 경우로 환원한다(그 경우의 증명은 어떤 의미에서는 스펙트럼 수열을 쓴다).

\medskip\noindent
실제로 $Y' \subset Y$가 임의의 열린집합이고 그 역상이
$X' \subset X$이면, $\tilde \xi|_{X'}$는 사상
$$
\bigoplus\nolimits_{i = 1}^N H^*_{dR}(Y'/S) \longrightarrow
H^*_{dR}(X'/S), \quad
(a_1, \ldots, a_N) \longmapsto  \sum \xi_i|_{X'} \cup f^*a_i
$$
을 $Y'$ 위의 코호몰로지에 유도한다. 위의 논의를 보라. 따라서 $Y$의
위상에 대한 기저를 찾아, 그 기저의 각 원소에 대해 명제가 성립함을 보이면
충분하다(특히 이렇게 할 때 사상 $\tilde \xi$는 잊어도 된다). 이는 $Y$와
$S$가 아핀인 경우로 환원되며, 그 경우는 보조정리
\ref{derham-lemma-global-generation-on-fibres}에서 다루었으므로 증명이 끝난다.
\end{proof}





\section{사영공간 다발 공식}
\label{derham-section-projective-space-bundle-formula}

\noindent
제목이 내용을 모두 말해 준다.

\begin{proposition}
\label{derham-proposition-projective-space-bundle-formula}
$X \to S$를 스킴의 사상이라 하자. $\mathcal{E}$를 국소 자유
$\mathcal{O}_X$-가군으로서 일정한 계수 $r$을 갖는다고 하자. 사상
$p : P = \mathbf{P}(\mathcal{E}) \to X$를 생각하자. 그러면 다음 규칙으로
주어지는 사상
$$
\bigoplus\nolimits_{i = 0, \ldots, r - 1} H^*_{dR}(X/S)
\longrightarrow
H^*_{dR}(P/S)
$$
$$
(a_0, \ldots, a_{r - 1}) \longmapsto
\sum\nolimits_{i = 0, \ldots, r - 1} c_1^{dR}(\mathcal{O}_P(1))^i \cup p^*(a_i)
$$
은 동형이다.
\end{proposition}

\begin{proof}
$\Spec(A) \subset X$를 아핀 열린집합으로 택하여, $\mathcal{E}$의 제한이
자명한 국소 자유 가군 $\mathcal{O}_{\Spec(A)}^{\oplus r}$이 되게 하자.
그러면 $P \times_X \Spec(A) = \mathbf{P}^{r - 1}_A$이다. 따라서 $p$는
고유하고 매끄럽다. 절 \ref{derham-section-projective-space}을 보라. 또한 류들
$c_1^{dR}(\mathcal{O}_P(1))^i$, $i = 0, 1, \ldots, r - 1$을 올
$X_y = \mathbf{P}^{r - 1}_y$에 제한하면 드람 코호몰로지
$H^*_{dR}(X_y/y)$를 $\kappa(y)$ 위에서 자유롭게 생성한다. 보조정리
\ref{derham-lemma-de-rham-cohomology-projective-space}를 보라. 이로써 명제
\ref{derham-proposition-global-generation-on-fibres}의 조건을 확인했으므로 결론을 얻는다.
\end{proof}

\begin{remark}
\label{derham-remark-projective-space-bundle-formula}
명제 \ref{derham-proposition-projective-space-bundle-formula}의 상황에서는 사상
$$
\tilde \xi :
\bigoplus\nolimits_{t = 0, \ldots, r - 1}
\Omega^\bullet_{X/S}[-2t]
\longrightarrow
Rp_*\Omega^\bullet_{P/S}
$$
이 $D(X, (X \to S)^{-1}\mathcal{O}_X)$에서 동형임도 얻는다. 이는 명제
\ref{derham-proposition-global-generation-on-fibres}를 적용하면 바로 따른다.
$t = 0$에 대응하는 화살표는 표준 사상
$c_{P/X} : \Omega^\bullet_{X/S} \to Rp_*\Omega^\bullet_{P/S}$
일 뿐이다. 절 \ref{derham-section-de-rham-complex}을 보라. 사실 이 특수한 경우에는
이 사상을 더 구체적으로 정할 수 있다. 즉 표준 사상
$$
\xi' : \Omega^\bullet_{P/S} \to \Omega^\bullet_{P/S}[2]
$$
을 생각하자. 이는 비고 \ref{derham-remark-cup-product-as-a-map}에서
$c_1^{dR}(\mathcal{O}_P(1))$에 대응하는 사상이다. 그러면
$$
\xi'[2(t - 1)] \circ \ldots \circ \xi'[2] \circ \xi' : 
\Omega^\bullet_{P/S} \to \Omega^\bullet_{P/S}[2t]
$$
은 비고 \ref{derham-remark-cup-product-as-a-map}에서
$c_1^{dR}(\mathcal{O}_P(1))^t$에 대응하는 사상이다. 명제
\ref{derham-proposition-global-generation-on-fibres}의 증명에서 한 선택들을 따라가면
다음 값을 얻는다.
$$
\tilde \xi|_{\Omega^\bullet_{X/S}[-2t]} =
Rp_*\xi'[-2] \circ \ldots \circ Rp_*\xi'[-2(t - 1)] \circ
Rp_*\xi'[-2t] \circ c_{P/X}[-2t]
$$
이는 우리 동형을 직합 성분 $\Omega^\bullet_{X/S}[-2t]$에 제한한 것이다.
여기서 아래에 사용할 다음 간단한 결과를 얻는다. 다음과 같이 놓자.
$$
M = \bigoplus\nolimits_{t = 1, \ldots, r - 1} \Omega^\bullet_{X/S}[-2t]
\quad\text{및}\quad
K = \bigoplus\nolimits_{t = 0, \ldots, r - 2} \Omega^\bullet_{X/S}[-2t]
$$
이들을 화살표 $\tilde \xi$의 정의역의 부분복합체로 본다. 위의 논의에서
형식적으로
$$
c_{P/X} \oplus
\tilde \xi|_M :
\Omega^\bullet_{X/S} \oplus M \longrightarrow
Rp_*\Omega^\bullet_{P/S}
$$
가 동형이고, 도식
$$
\xymatrix{
K \ar[d]_{\tilde \xi|_K} \ar[r]_{\text{id}} &
M[2] \ar[d]^{(\tilde \xi|_M)[2]} \\
Rp_*\Omega^\bullet_{P/S} \ar[r]^{Rp_*\xi'} &
Rp_*\Omega^\bullet_{P/S}[2]
}
$$
이 가환함이 따른다. 여기서 $\text{id} : K \to M[2]$는 $t$에 대응하는
$K$의 분해의 직합 성분을 $t + 1$에 대응하는 $M$의 분해의 직합
성분과 동일시한다.
\end{remark}









\section{인자를 따른 로그 극}
\label{derham-section-divisor}

\noindent
스킴의 사상 $X \to S$와 유효 카르티에 인자 $Y \subset X$가 주어졌다고
하자. $X$가 $Y$를 따라 에탈 국소적으로 $Y \times \mathbf{A}^1$과 같은
모양이면, 복합체들의 표준적인 짧은 완전열
$$
0 \to \Omega^\bullet_{X/S} \to
\Omega^\bullet_{X/S}(\log Y) \to
\Omega^\bullet_{Y/S}[-1] \to 0
$$
이 존재하며, 이 절에서는 그 여러 좋은 성질을 논한다. 정규 교차 인자에서
출발하는 이 구성의 변형도 있다
(스킴의 에탈 사상의 정의 \ref{etale-definition-strict-normal-crossings}).
이는 다른 곳에서 논할 것이다(향후 참조를 여기에 넣는다).

\begin{definition}
\label{derham-definition-local-product}
스킴의 사상 $X \to S$와 유효 카르티에 인자 $Y \subset X$가 주어졌다고
하자. {\it $Y \subset X$에 대한 로그 극 드람 복합체가 $S$ 위에서
정의된다}고 하는 것은 다음을 뜻한다. 모든 $y \in Y$와
$f \in \mathcal{O}_{X, y}$가 $Y$의 국소 방정식일 때 다음이 성립한다.
\begin{enumerate}
\item $\mathcal{O}_{X, y} \to \Omega_{X/S, y}$, $g \mapsto g \text{d}f$
는 분할 단사이고,
\item $\Omega^p_{X/S, y}$에는 $f$-꼬임이 없으며, 이는 모든 $p$에 대해
성립한다.
\end{enumerate}
\end{definition}

\noindent
간단한 국소 계산에 따르면 각 $y \in Y$마다 조건 (1), (2)를 만족하는
국소 방정식 $f$를 하나 찾으면 충분하다.

\begin{lemma}
\label{derham-lemma-log-complex}
스킴의 사상 $X \to S$와 유효 카르티에 인자 $Y \subset X$가 주어졌다고
하자. $Y \subset X$에 대한 로그 극 드람 복합체가 $S$ 위에서 정의된다고
가정하자. 그러면 복합체들의 표준적인 짧은 완전열
$$
0 \to \Omega^\bullet_{X/S} \to
\Omega^\bullet_{X/S}(\log Y) \to
\Omega^\bullet_{Y/S}[-1] \to 0
$$
이 존재한다.
\end{lemma}

\begin{proof}
가정에 의해, 모든 $y \in Y$와 국소 방정식
$f \in \mathcal{O}_{X, y}$가 $Y$를 정의할 때
$$
\Omega_{X/S, y} = \mathcal{O}_{X, y}\text{d}f \oplus M
\quad\text{및}\quad
\Omega^p_{X/S, y} = \wedge^{p - 1}(M)\text{d}f \oplus \wedge^p(M)
$$
어떤 가군 $M$에 대해 위 등식이 성립하며, 이때 $f$-꼬임이 없는
외대수 거듭제곱은 $\wedge^p(M)$이다. 따라서
$$
\Omega^p_{Y/S, y} = \wedge^p(M/fM) = \wedge^p(M)/f\wedge^p(M)
$$
이다. 이하에서는 이 사실들을 암묵적으로 사용한다. 특히 층
$\Omega^p_{X/S}$에는 지지가 $Y$에 들어가는 영이 아닌 국소 단면이 없고,
표준적인 포함
$$
\Omega^p_{X/S} \subset \Omega^p_{X/S}(Y)
$$
을 얻는다. 평탄성 심화의 절 \ref{flat-section-eta}를 보라.
아핀 열린 부분스킴 $U = \Spec(A)$를 잡자. 이때
$Y \cap U = V(f)$이고 $f \in A$는 영인자가 아닌 원소라고 하자.
$\mathcal{O}_U$-부분가군 중 $\Omega^p_{X/S}(Y)|_U$ 안에서
$\Omega^p_{X/S}|_U$와
$\text{d}\log(f) \wedge \Omega^{p - 1}_{X/S}$로 생성되는 것을 생각하자.
여기서 $\text{d}\log(f) = f^{-1}\text{d}(f)$이다. 이 부분가군은 $f$의
선택에 무관하다. 실제로 $Y$의 아이디얼을 $U$ 위에서 생성하는 다른
원소는 $uf$와 같다. 여기서 $u \in A$는 단원이며,
$$
\text{d}\log(uf) - \text{d}\log(f) = \text{d}\log(u) = u^{-1}\text{d}u
$$
를 얻는다. 이는 $\Omega_{X/S}$의 $U$ 위 단면이다. 이 국소 층들은
서로 붙어서 준연접 부분가군
$$
\Omega^p_{X/S} \subset \Omega^p_{X/S}(\log Y) \subset \Omega^p_{X/S}(Y)
$$
을 이룬다. $\Omega^p_{Y/S}$를 준연접 $\mathcal{O}_X$-가군으로 간주하기로
하자. 다음 규칙으로 정의되는 유일한 전사 $\mathcal{O}_X$-선형 사상
$$
\text{Res} : \Omega^p_{X/S}(\log Y) \to \Omega^{p - 1}_{Y/S}
$$
이 존재한다.
$$
\text{Res}(\eta' + \text{d}\log(f) \wedge \eta) = \eta|_{Y \cap U}
$$
이는 위와 같은 모든 열린집합 $U$와 모든
$\eta' \in \Omega^p_{X/S}(U)$ 및 $\eta \in \Omega^{p - 1}_{X/S}(U)$에
대한 규칙이다. 미분형식 $\eta$가 $U$ 위에서 $Y \cap U$에 제한될 때
영이면
$\eta = \text{d}f \wedge \eta' + f\eta''$이다. 여기서 $\eta'$, $\eta''$은
$U$ 위의 어떤 미분형식들이다. 따라서 짧은 완전열
$$
0 \to \Omega^p_{X/S} \to \Omega^p_{X/S}(\log Y) \to \Omega^{p - 1}_{Y/S} \to 0
$$
을 얻으며, 이는 모든 $p$에 대해 성립한다. 이제 미분
$\Omega^p_{X/S}(\log Y) \to \Omega^{p + 1}_{X/S}(\log Y)$를 정의해야 한다.
부분층 $\Omega^p_{X/S}$에서는 $X$의 $S$ 위 드람 복합체의 미분을
사용한다. 마지막으로
$\text{d}(\text{d}\log(f) \wedge \eta) = -\text{d}\log(f) \wedge \text{d}\eta$.
로 정의한다. 이 부호는 $\Omega^\bullet_{X/S}$ 위의 라이프니츠 법칙에 의해
강제되며 $\Omega^\bullet_{Y/S}[-1]$ 위의 미분과 양립한다. 이 후자의
미분은 호몰로지 대수의 정의 \ref{homology-definition-shift-cochain}에서 택한
부호 규약에 따라 결국 $-\text{d}_{Y/S}$이다. 이로써 보조정리에 적은
복합체들의 짧은 완전열을 얻는다.
\end{proof}

\begin{definition}
\label{derham-definition-log-complex}
스킴의 사상 $X \to S$와 유효 카르티에 인자 $Y \subset X$가 주어졌다고
하자. $Y \subset X$에 대한 로그 극 드람 복합체가 $S$ 위에서 정의된다고
가정하자. 그러면 보조정리 \ref{derham-lemma-log-complex}에서 구성한 복합체
$$
\Omega^\bullet_{X/S}(\log Y)
$$
를 {\it $Y \subset X$에 대한 $S$ 위의 로그 극 드람 복합체}라 한다.
\end{definition}

\noindent
이 복합체는 여러 좋은 성질을 갖는다.

\begin{lemma}
\label{derham-lemma-multiplication-log}
$p : X \to S$를 스킴의 사상이라 하고 $Y \subset X$를 유효 카르티에
인자라 하자. $Y \subset X$에 대한 로그 극 드람 복합체가 $S$ 위에서
정의된다고 가정하자.
\begin{enumerate}
\item 사상들
$\wedge : \Omega^p_{X/S} \times \Omega^q_{X/S} \to \Omega^{p + q}_{X/S}$
은 $\mathcal{O}_X$-쌍선형 사상들
$$
\wedge : \Omega^p_{X/S}(\log Y) \times \Omega^q_{X/S}(\log Y)
\to \Omega^{p + q}_{X/S}(\log Y)
$$
로 유일하게 연장되며, 라이프니츠 법칙
$
\text{d}(\omega \wedge \eta) = \text{d}(\omega) \wedge \eta +
(-1)^{\deg(\omega)} \omega \wedge \text{d}(\eta)$,
을 만족한다.
\item (1)의 곱셈에 관해 사상
$\Omega^\bullet_{X/S} \to \Omega^\bullet_{X/S}(\log(Y)$
은 미분 등급 $\mathcal{O}_S$-대수의 준동형이다.
\item (1)의 사상들을 통해 $\Omega^p_{X/S}(\log Y) =
\wedge^p(\Omega^1_{X/S}(\log Y))$를 얻는다.
\item 사상
$\text{Res} : \Omega^\bullet_{X/S}(\log Y) \to \Omega^\bullet_{Y/S}[-1]$
은
$$
\text{Res}(\omega \wedge \eta) = \text{Res}(\omega) \wedge \eta|_Y
$$
를 만족한다. 여기서 $\omega$는 $\Omega^p_{X/S}(\log Y)$의 국소 단면이고
$\eta$는 $\Omega^q_{X/S}$의 국소 단면이다.
\end{enumerate}
\end{lemma}

\begin{proof}
이는 보조정리 \ref{derham-lemma-log-complex}의 증명에서 복합체를 국소적으로
구성한 방식으로부터 직접 계산하면 따른다. 세부 사항은 생략한다.
\end{proof}

\noindent
스킴들의 가환 도식
$$
\xymatrix{
X' \ar[r]_f \ar[d] & X \ar[d] \\
S' \ar[r] & S
}
$$
을 생각하자. $Y \subset X$를 당김 $Y' = f^*Y$가 정의되는 유효 카르티에
인자라 하자(인자의 정의
\ref{divisors-definition-pullback-effective-Cartier-divisor}).
$Y \subset X$에 대한 로그 극 드람 복합체가 $S$ 위에서 정의되고,
$Y' \subset X'$에 대한 로그 극 드람 복합체도 $S'$ 위에서 정의된다고
가정하자. 그러면 복합체들의 짧은 완전열 사이의 사상
$$
\xymatrix{
0 \ar[r] &
f^{-1}\Omega^\bullet_{X/S} \ar[r] \ar[d] &
f^{-1}\Omega^\bullet_{X/S}(\log Y) \ar[r] \ar[d] &
f^{-1}\Omega^\bullet_{Y/S}[-1] \ar[r] \ar[d] &
0 \\
0 \ar[r] &
\Omega^\bullet_{X'/S'} \ar[r] &
\Omega^\bullet_{X'/S'}(\log Y') \ar[r] &
\Omega^\bullet_{Y'/S'}[-1] \ar[r] &
0
}
$$
을 얻는다. 선형화하면 모든 $p$에 대해 선형 사상
$f^*\Omega^p_{X/S}(\log Y) \to \Omega^p_{X'/S'}(\log Y')$을 얻는다.

\begin{lemma}
\label{derham-lemma-gysin-via-log-complex}
$f : X \to S$를 스킴의 사상이라 하고 $Y \subset X$를 유효 카르티에
인자라 하자. $Y \subset X$에 대한 로그 극 드람 복합체가 $S$ 위에서
정의된다고 가정하자.
$$
\delta : \Omega^\bullet_{Y/S} \to \Omega^\bullet_{X/S}[2]
$$
를 $D(X, f^{-1}\mathcal{O}_S)$ 안에서 보조정리
\ref{derham-lemma-log-complex}의 짧은 완전열로부터 나오는 ``경계'' 사상이라
하자. 또한
$$
\xi' : \Omega^\bullet_{X/S} \to \Omega^\bullet_{X/S}[2]
$$
를 $D(X, f^{-1}\mathcal{O}_S)$ 안에서 비고
\ref{derham-remark-cup-product-as-a-map}의 사상 중
$\xi = c_1^{dR}(\mathcal{O}_X(-Y))$에 대응하는 것이라 하자. 또한
$$
\zeta' : \Omega^\bullet_{Y/S} \to \Omega^\bullet_{Y/S}[2]
$$
를 $D(Y, f|_Y^{-1}\mathcal{O}_S)$ 안에서 비고
\ref{derham-remark-cup-product-as-a-map}의 사상 중
$\zeta = c_1^{dR}(\mathcal{O}_X(-Y)|_Y)$에 대응하는 것이라 하자. 그러면 도식
$$
\xymatrix{
\Omega^\bullet_{X/S} \ar[d]_{\xi'} \ar[r] &
\Omega^\bullet_{Y/S} \ar[d]^{\zeta'} \ar[ld]_\delta \\
\Omega^\bullet_{X/S}[2] \ar[r] &
\Omega^\bullet_{Y/S}[2]
}
$$
은 $D(X, f^{-1}\mathcal{O}_S)$에서 가환한다.
\end{lemma}

\begin{proof}
더 정확히 말해, $\delta$를 이동된 짧은 완전열
$$
0 \to \Omega^\bullet_{X/S}[1] \to
\Omega^\bullet_{X/S}(\log Y)[1] \to
\Omega^\bullet_{Y/S} \to 0
$$
에 대응하는 경계 사상으로 정의한다. 각 삼각형이 가환함을 보이면 충분하다.
$\mathcal{L} = \mathcal{O}_X(-Y)$라 놓자. 선다발
$\pi : L \to X$를
$\pi_*\mathcal{O}_L = \bigoplus_{n \geq 0} \mathcal{L}^{\otimes n}$인
것으로 정의하자.

\medskip\noindent
왼쪽 위 삼각형의 가환성.
보조정리 \ref{derham-lemma-the-complex-for-L-star-gives-chern-class}에 의해
사상 $\xi'$는 보조정리 \ref{derham-lemma-the-complex-for-L-star}에서 주어진
삼각형의 경계 사상이다. 함자성에 의해 짧은 완전열 사이의 사상
$$
\xymatrix{
0 \ar[r] &
\Omega^\bullet_{X/S}[1] \ar[r] \ar[d] &
\Omega^\bullet_{L^\star/S, 0}[1] \ar[r] \ar[d] &
\Omega^\bullet_{X/S} \ar[r] \ar[d] &
0 \\
0 \ar[r] &
\Omega^\bullet_{X/S}[1] \ar[r] &
\Omega^\bullet_{X/S}(\log Y)[1] \ar[r] &
\Omega^\bullet_{Y/S} \ar[r] &
0
}
$$
이 존재함을 보이면 충분하다. 여기서 왼쪽과 오른쪽 수직 화살표는 자명한
것들이다. 가운데 수직 화살표는 다음 규칙으로 정의할 수 있다.
$$
\omega' + \text{d}\log(s) \wedge \omega \longmapsto
\omega' + \text{d}\log(f) \wedge \omega
$$
여기서 $\omega', \omega$는 $\Omega^\bullet_{X/S}$의 국소 단면이고,
$s$는 $\mathcal{L}$의 국소 생성원이며,
$f \in \mathcal{O}_X(-Y)$는 이에 대응하는 $Y$의 아이디얼 층을
$X$ 위에서 본 단면이다. 보조정리 \ref{derham-lemma-the-complex-for-L-star}와
\ref{derham-lemma-log-complex}에서 사상들을 구성한 방식이 정확히 일치하므로
이 정의가 잘 작동한다.

\medskip\noindent
오른쪽 아래 삼각형의 가환성. $\overline{L}$을 $L$의 $Y$에 대한
제한이라 하자. 보조정리
\ref{derham-lemma-the-complex-for-L-star-gives-chern-class}에 의해 사상
$\zeta'$는 선다발 $\overline{L}$을 $Y$ 위에서 사용하여 보조정리
\ref{derham-lemma-the-complex-for-L-star}에서 얻은 삼각형의 경계 사상이다.
함자성에 의해 짧은 완전열 사이의 사상
$$
\xymatrix{
0 \ar[r] &
\Omega^\bullet_{X/S}[1] \ar[r] \ar[d] &
\Omega^\bullet_{X/S}(\log Y)[1] \ar[r] \ar[d] &
\Omega^\bullet_{Y/S} \ar[r] \ar[d] &
0 \\
0 \ar[r] &
\Omega^\bullet_{Y/S}[1] \ar[r] &
\Omega^\bullet_{\overline{L}^\star/S, 0}[1] \ar[r] &
\Omega^\bullet_{Y/S} \ar[r] &
0 \\
}
$$
이 존재함을 보이면 충분하다. 여기서 왼쪽과 오른쪽 수직 화살표는 자명한
것들이다. 가운데 수직 화살표는 다음 규칙으로 정의할 수 있다.
$$
\omega' + \text{d}\log(f) \wedge \omega \longmapsto
\omega'|_Y + \text{d}\log(\overline{s}) \wedge \omega|_Y
$$
여기서 $\omega', \omega$는 $\Omega^\bullet_{X/S}$의 국소 단면이고,
$f$는 $\mathcal{O}_X(-Y)$의 국소 생성원을 $X$ 위의 함수로 본 것이며,
$\overline{s}$는 $f|_Y$를
$\mathcal{L}|_Y = \mathcal{O}_X(-Y)|_Y$의 단면으로 본 것이다.
보조정리 \ref{derham-lemma-the-complex-for-L-star}와
\ref{derham-lemma-log-complex}에서 사상들을 구성한 방식이 정확히 일치하므로
이 정의가 잘 작동한다.
\end{proof}

\begin{lemma}
\label{derham-lemma-log-complex-consequence}
$X \to S$를 스킴의 사상이라 하고 $Y \subset X$를 유효 카르티에 인자라
하자. $Y \subset X$에 대한 로그 극 드람 복합체가 $S$ 위에서 정의된다고
가정하자. $b \in H^m_{dR}(X/S)$를 $Y$에 제한하면 영이 되는 드람
코호몰로지류라 하자. 그러면
$c_1^{dR}(\mathcal{O}_X(Y)) \cup b = 0$은
$H^{m + 2}_{dR}(X/S)$에서 성립한다.
\end{lemma}

\begin{proof}
이는 보조정리 \ref{derham-lemma-gysin-via-log-complex}에서 곧바로 따른다. 실제로
$$
c_1^{dR}(\mathcal{O}_X(Y)) \cup b =
-c_1^{dR}(\mathcal{O}_X(-Y)) \cup b = -\xi'(b) = -\delta(b|_Y) = 0
$$
이므로 원하는 결론을 얻는다. 둘째 등식은 비고
\ref{derham-remark-cup-product-as-a-map}를 보라.
\end{proof}

\begin{lemma}
\label{derham-lemma-check-log-smooth}
$X \to T \to S$를 스킴의 사상들이라 하고 $Y \subset X$를 유효 카르티에
인자라 하자. $X \to T$와 $Y \to T$가 모두 매끄러우면
$Y \subset X$에 대한 로그 극 드람 복합체가 $S$ 위에서 정의된다.
\end{lemma}

\begin{proof}
$y \in Y$를 한 점이라 하자. 사상 심화의 보조정리
\ref{more-morphisms-lemma-etale-local-structure}에 의해 정수 $0 \geq m$과
가환 도식
$$
\xymatrix{
Y \ar[d] &
V \ar[l] \ar[d] \ar[r] &
\mathbf{A}^m_T
\ar[d]^{(a_1, \ldots, a_m) \mapsto (a_1, \ldots, a_m, 0)} \\
X &
U \ar[l] \ar[r]^-\pi &
\mathbf{A}^{m + 1}_T
}
$$
이 존재한다. 여기서 $U \subset X$는 열린집합이고, $V = Y \cap U$이며,
$\pi$는 에탈이고, $V = \pi^{-1}(\mathbf{A}^m_T)$이며, $y \in V$이다.
$z \in \mathbf{A}^m_T$를 $y$의 상이라 하자. 보조정리
\ref{derham-lemma-etale}에 의해
$$
\Omega^p_{X/S, y} = \Omega^p_{\mathbf{A}^{m + 1}_T/S, z}
\otimes_{\mathcal{O}_{\mathbf{A}^{m + 1}_T, z}} \mathcal{O}_{X, x}
$$
이다. 좌표함수들 $x_1, \ldots, x_{m + 1}$을
$\mathbf{A}^{m + 1}_T$ 위에서 생각하자. 정의
\ref{derham-definition-local-product}의 조건 (1), (2)는 국소 좌표의 선택에
의존하지 않으므로, $f$를 $x_{m + 1}$의, 평탄한 국소환 준동형
$\mathcal{O}_{\mathbf{A}^{m + 1}_T, z} \to \mathcal{O}_{X, x}$
아래의 상으로 둘 때 조건 (1), (2)를 확인하면
충분하다. 따라서 $\mathbf{A}^m_T \subset \mathbf{A}^{m + 1}_T$와 점
$z$에 대해 조건 (1), (2)를 확인하면 충분함을 알 수 있다. 이 경우를
증명하기 위해 $S = \Spec(A)$와 $T = \Spec(B)$가 아핀이라고 가정해도
된다. 환 사상 $A \to B$를 사상 $T \to S$에 대응하는 것으로 잡고,
$P = B[x_1, \ldots, x_{m + 1}]$라 놓자. 그러면
$\mathbf{A}^{m + 1}_T = \Spec(P)$이고
$$
\Omega_{P/A} = \Omega_{B/A} \otimes_B P \oplus
\bigoplus\nolimits_{j = 1, \ldots, m} P \text{d}x_j \oplus
P \text{d}x_{m + 1}
$$
이다. 따라서 사상 $P \to \Omega_{P/A}$,
$g \mapsto g \text{d}x_{m + 1}$은 분할 단사이고,
$x_{m + 1}$은 $\Omega^p_{P/A}$ 위에서 영인자가 아니며, 이는 모든
$p \geq 0$에 대해 성립한다. $z$에 대응하는
소 아이디얼에서 국소화하면 증명이 끝난다.
\end{proof}

\begin{remark}
\label{derham-remark-check-log-completion-1}
$S$를 국소 뇌터 스킴이라 하고 $X$를 $S$ 위에서 국소 유한형이라 하자.
$Y \subset X$를 유효 카르티에 인자라 하자. 사상
$$
\mathcal{O}_{X, y}^\wedge \longrightarrow \mathcal{O}_{Y, y}^\wedge
$$
이 모든 $y \in Y$에 대해 단면을 가지면,
$Y \subset X$에 대한 로그 극 드람 복합체가 $S$ 위에서
정의된다. 이 결과가 필요해지면 정확한 명제를 정식화하고 여기에 증명을
덧붙일 것이다.
\end{remark}

\begin{remark}
\label{derham-remark-check-log-completion-2}
$S$를 국소 뇌터 스킴이라 하고 $X$를 $S$ 위에서 국소 유한형이라 하자.
$Y \subset X$를 유효 카르티에 인자라 하자. 모든 $y \in Y$에 대해
$S$ 위의 스킴들의 도식
$$
X \xleftarrow{\varphi} U \xrightarrow{\psi} V
$$
을 찾을 수 있고 $\varphi$가 에탈이며
$\psi|_{\varphi^{-1}(Y)} : \varphi^{-1}(Y) \to V$도 에탈이면,
$Y \subset X$에 대한 로그 극 드람 복합체가 $S$ 위에서 정의된다.
특히 쌍 $(X, Y)$가 에탈 국소적으로
$(V \times \mathbf{A}^1, V \times \{0\})$과 같은 모양인 경우가 이에
해당한다. 이 결과가 필요해지면 정확한 명제를 정식화하고 여기에 증명을
덧붙일 것이다.
\end{remark}

















\section{계산}
\label{derham-section-calculations}

\noindent
이 절에서는 표준 블로업에 대한 몇 가지 호지 코호몰로지군과 드람
코호몰로지군을 계산한다.

\medskip\noindent
환 $R$를 고정하고 $S = \Spec(R)$라 놓자. 정수 $0 \leq m$과
$1 \leq n$을 고정하자. 닫힌 몰입
$$
Z = \mathbf{A}^m_S \longrightarrow \mathbf{A}^{m + n}_S = X,\quad
(a_1, \ldots, a_m) \mapsto (a_1, \ldots, a_m, 0, \ldots 0).
$$
을 생각하자. $L$을 $X$의 닫힌 부분스킴 $Z$를 따른 블로업이라
하자. 다음과 같이 쓰자.
$$
X =
\mathbf{A}^{m + n}_S =
\Spec(A)
\quad\text{단}\quad
A = R[x_1, \ldots, x_m, y_1, \ldots, y_n]
$$
$A = R[x_1, \ldots, x_m, y_1, \ldots, y_n]$를 등급 $R$-대수로
간주하되, $\deg(x_i) = 0$, $\deg(y_j) = 1$로 놓는다.
이 등급에 대해
$$
P =
\text{Proj}(A) =
\mathbf{A}^m_S \times_S \mathbf{P}^{n - 1}_S =
Z \times_S \mathbf{P}^{n - 1}_S =
\mathbf{P}^{n - 1}_Z
$$
이다. $Z$를 $X$ 안에서 잘라내는 아이디얼은 $A_+$임에 유의하자.
따라서 $L$은 리스 대수
$$
A \oplus A_+ \oplus (A_+)^2 \oplus \ldots =
\bigoplus\nolimits_{d \geq 0} A_{\geq d}
$$
의 Proj이다. 따라서 $L$은 사상 심화의 절
\ref{more-morphisms-section-proj-spec}에서 더 일반적으로 연구한 현상의
한 예이다. 거기서 얻은 관찰들을 더 언급하지 않고 사용하겠다. 특히 가환
도식
$$
\xymatrix{
P \ar[r]_0 \ar[d]_p &
L \ar[r]_-\pi \ar[d]^b &
P \ar[d]^p \\
Z \ar[r]^i &
X \ar[r] &
Z
}
$$
을 얻는다. 여기서 $\pi : L \to P$는
$P = Z \times_S \mathbf{P}^{n - 1}_S$ 위의 선다발이고, 영단면은
$0$이며, 그 상 $E = 0(P) \subset L$은 블로업 $b$의 예외 인자이다.

\begin{lemma}
\label{derham-lemma-comparison}
$a \geq 0$에 대해 다음이 성립한다.
\begin{enumerate}
\item 사상 $\Omega^a_{X/S} \to b_*\Omega^a_{L/S}$은 동형이다.
\item 사상 $\Omega^a_{Z/S} \to p_*\Omega^a_{P/S}$은 동형이다.
\item 사상 $Rb_*\Omega^a_{L/S} \to i_*Rp_*\Omega^a_{P/S}$은 차수가
$\geq 1$인 코호몰로지 층들 위에서 동형이다.
\end{enumerate}
\end{lemma}

\begin{proof}
먼저 (2)를 증명하자. $P = Z \times_S \mathbf{P}^{n - 1}_S$이므로
$$
\Omega^a_{P/S} = \bigoplus\nolimits_{a = r + s}
\text{pr}_1^*\Omega^r_{Z/S} \otimes
\text{pr}_2^*\Omega^s_{\mathbf{P}^{n - 1}_S/S}
$$
임을 알 수 있다. $p = \text{pr}_1$임을 상기하면, 사영 공식
(층 코호몰로지의 보조정리 \ref{cohomology-lemma-projection-formula})에 의해
$$
p_*\Omega^a_{P/S} = \bigoplus\nolimits_{a = r + s}
\Omega^r_{Z/S} \otimes
\text{pr}_{1, *}\text{pr}_2^*\Omega^s_{\mathbf{P}^{n - 1}_S/S}
$$
을 얻는다. 절 \ref{derham-section-projective-space}의 계산, 특히 보조정리
\ref{derham-lemma-hodge-cohomology-projective-space}의 증명에 의해
$\text{pr}_{1, *}\text{pr}_2^*\Omega^s_{\mathbf{P}^{n - 1}_S/S} = 0$이다.
단, $s = 0$인 경우에는
$\text{pr}_{1, *}\mathcal{O}_P = \mathcal{O}_Z$를 얻는다. 이로써 (2)가
증명되었다.

\medskip\noindent
절 \ref{derham-section-line-bundle}의 내용, 특히 보조정리
\ref{derham-lemma-push-omega-a}에 의해
$\pi_*\Omega^a_{L/S} = \Omega^a_{P/S} \oplus
\bigoplus_{k \geq 1} \Omega^a_{L/S, k}$.
위 도식에서 합성 $\pi \circ 0$은 $P$ 위의 항등사상이므로, (3)을
증명하려면 $\Omega^a_{L/S, k}$의 고차 코호몰로지가 $k > 0$일 때
영임을 보이면 충분하다. 보조정리
\ref{derham-lemma-the-complex-for-L-star}와 \ref{derham-lemma-push-omega-a}에 의해
짧은 완전열들
$$
0 \to \Omega^a_{P/S} \otimes \mathcal{O}_P(k)
\to \Omega^a_{L/S, k} \to
\Omega^{a - 1}_{P/S} \otimes \mathcal{O}_P(k) \to 0
$$
이 존재한다. 여기서 $\Omega^{a - 1}_{P/S} = 0$이며, 이는 $a = 0$인
경우이다.
$P = Z \times_S \mathbf{P}^{n - 1}_S$이므로
$$
\Omega^a_{P/S} = \bigoplus\nolimits_{i + j = a}
\Omega^i_{Z/S} \boxtimes \Omega^j_{\mathbf{P}^{n - 1}_S/S}
$$
이다. 이는 보조정리 \ref{derham-lemma-de-rham-complex-product}에서 따른다.
$\Omega^i_{Z/S}$가 유한 계수의 자유 가군이므로
$\mathcal{O}_Z \boxtimes \Omega^j_{\mathbf{P}^{n - 1}_S/S}(k)$의 고차
코호몰로지가 $k > 0$일 때 영임을 보이면 충분하다. 이는 보조정리
\ref{derham-lemma-twisted-hodge-cohomology-projective-space}를
$P = Z \times_S \mathbf{P}^{n - 1}_S = \mathbf{P}^{n - 1}_Z$에 적용하면
따른다. 이로써 (3)의 증명이 끝났다.

\medskip\noindent
이제 (1)을 증명해야 한다. $n = 1$이면 유효 카르티에 인자를 블로업하는
것이고 $b$는 동형이므로 (1)이 성립한다. $n > 1$이면 합성
$$
\Gamma(X, \Omega^a_{X/S})
\to
\Gamma(L, \Omega^a_{L/S})
\to
\Gamma(L \setminus E, \Omega^a_{L/S})
=
\Gamma(X \setminus Z, \Omega^a_{X/S})
$$
은 동형이다. 실제로 $\Omega^a_{X/S}$는 유한 자유이다(작은 세부 사항은
생략한다). 따라서 (1)이 성립하지 않을 수 있는 유일한 경우는
$\Gamma(L, \Omega^a_{L/S})$의 영이 아닌 원소 중 $E = 0(P)$의 바깥에서
영이 되는 것이 존재하는 경우이다. $L$은 $P$ 위의 선다발이고 영단면은
$0 : P \to L$이므로, 선다발 위에서 미분층의 영이 아닌 단면이 영단면의
바깥에서 항등적으로 영이 될 수 없음을 보이면 충분하다. 이는 국소 계산을
통해(이 방법이 바람직하다), 또는
$\Omega_{L/S, k} \subset \Omega_{L^\star/S, k}$임을 사용하면 알 수 있다
(위의 참조들을 보라).
\end{proof}

\noindent
이 지점에서 독자는 다음 절로 넘어가기를 권한다.

\begin{lemma}
\label{derham-lemma-comparison-bis}
$a \geq 0$에 대해 표준적인 사상들
$$
b^*\Omega^a_{X/S} \longrightarrow
\Omega^a_{L/S} \longrightarrow
b^*\Omega^a_{X/S} \otimes_{\mathcal{O}_L} \mathcal{O}_L((n - 1)E)
$$
이 존재하며, 그 합성은 포함
$\mathcal{O}_L \subset \mathcal{O}_L((n - 1)E)$로부터 유도된다.
\end{lemma}

\begin{proof}
표시된 공식의 첫째 화살표는 절 \ref{derham-section-de-rham-complex}에서
논하였다. 둘째 화살표를 얻으려면 $\Omega^a_{L/S}$의 국소 단면을
$b^*\Omega^a_{X/S}$의 ``유리형 단면''으로 보았을 때 그 극의 차수가
최대 $n - 1$이고, 그 극은 $E$를 따라 생김을 보여야 한다.
이를 위해 $L$ 위의 아핀
국소 좌표조각에서 계산하자. $L$은 아핀 블로업 대수
$A[\frac{I}{y_i}]$들의 스펙트럼으로 덮인다. 여기서 $I = A_{+}$는
$y_1, \ldots, y_n$으로 생성되는 아이디얼이다. 가환대수의 절
\ref{algebra-section-blow-up}과 인자의 보조정리
\ref{divisors-lemma-blowing-up-affine}를 보라. 대칭성에 의해 $i = 1$에
대응하는 좌표조각에서 계산하면 충분하다. 그러면
$$
A[\frac{I}{y_1}] = R[x_1, \ldots, x_m, y_1, t_2, \ldots, t_n]
$$
이고 $t_i = y_i/y_1$이다. 대수학 심화의 보조정리
\ref{more-algebra-lemma-blowup-regular-sequence}를 보라. 따라서 가군
$\Omega^1_{L/S}$은 대응하는 아핀 열린집합 위에서
$\text{d}x_1, \ldots, \text{d}x_m$, $\text{d}y_1$, 그리고
$\text{d}t_1, \ldots, \text{d}t_n$으로 자유롭게 생성된다.
물론 이 생성원들 중 처음 $m + 1$개는 $b^*\Omega^1_{X/S}$에서 오고,
나머지 $n - 1$개에 대해서는
$$
\text{d}t_j =
\text{d}\frac{y_j}{y_1} =
\frac{1}{y_1}\text{d}y_j - \frac{y_j}{y_1^2}\text{d}y_1 =
\frac{\text{d}y_j - t_j \text{d}y_1}{y_1}
$$
이다. 이 식은 차수 $1$의 극을 $E$를 따라 갖는다. 실제로 이 좌표조각에서
$E$는 $y_1$로 잘려 나온다. 이 원소들 가운데 $a$개를 취한 외승들이 이
좌표조각 위의 $\Omega^a_{L/S}$의 기저를 이루고, 그 가운데
최대 $n - 1$개만 $\text{d}t_j$일 수 있으므로 증명이 끝난다.
\end{proof}

\begin{lemma}
\label{derham-lemma-blowup-twist-same-cohomology}
$E = 0(P)$를 블로업 $b$의 예외 인자라 하자. 임의의 국소 자유
$\mathcal{O}_X$-가군 $\mathcal{E}$와 $0 \leq i \leq n - 1$에 대해 사상
$$
\mathcal{E}
\longrightarrow
Rb_*(b^*\mathcal{E} \otimes_{\mathcal{O}_L} \mathcal{O}_L(iE))
$$
은 $D(\mathcal{O}_X)$에서 동형이다.
\end{lemma}

\begin{proof}
사영 공식에 의해 $\mathcal{E} = \mathcal{O}_X$인 경우만 보이면 충분하다.
층 코호몰로지의 보조정리 \ref{cohomology-lemma-projection-formula}를 보라.
$X$가 아핀이므로 사상들
$$
H^0(X, \mathcal{O}_X) \to
H^0(L, \mathcal{O}_L) \to
H^0(L, \mathcal{O}_L(iE))
$$
이 동형이고 $H^j(X, \mathcal{O}_L(iE)) = 0$임을 보이면 충분하다. 여기서
$j > 0$이고 $0 \leq i \leq n - 1$이다. 스킴의 코호몰로지의 보조정리
\ref{coherent-lemma-quasi-coherence-higher-direct-images-application}를
보라. $\pi$가 아핀이므로 $\pi_*$를 취한 뒤 전역단면과 코호몰로지를
계산할 수 있다. 스킴의 코호몰로지의 보조정리
\ref{coherent-lemma-relative-affine-cohomology}를 보라. $n = 1$이면
$L \to X$는 동형이고 $i = 0$이므로 첫째 명제가 성립한다.
$n > 1$이면 합성
$$
H^0(X, \mathcal{O}_X) \to H^0(L, \mathcal{O}_L) \to
H^0(L, \mathcal{O}_L(iE)) \to H^0(L \setminus E, \mathcal{O}_L) =
H^0(X \setminus Z, \mathcal{O}_X)
$$
을 생각하자. 이 경우
$H^0(X \setminus Z, \mathcal{O}_X) = H^0(X, \mathcal{O}_X)$이다.
실제로 $Z$는 여차원 $n \geq 2$로 $X$ 안에 놓인다(세부 사항은 생략한다).
따라서 첫째 명제가 성립한다. 둘째 명제를 위해
$\mathcal{O}_L(E) = \mathcal{O}_L(-1)$임을 상기하자. 인자의 보조정리
\ref{divisors-lemma-blowing-up-gives-effective-Cartier-divisor}를 보라.
그러므로
$$
\pi_*\mathcal{O}_L(iE) =
\pi_*\mathcal{O}_L(-i) =
\bigoplus\nolimits_{k \geq -i} \mathcal{O}_P(k)
$$
이다. 이는 사상 심화의 절
\ref{more-morphisms-section-proj-spec}에서 논하였다.
$P = \mathbf{P}^{n - 1}_Z$ 위 구조층의 꼬임들의 코호몰로지가 영이라는
스킴의 코호몰로지의 보조정리
\ref{coherent-lemma-cohomology-projective-space-over-ring}를 적용하면
결론을 얻는다.
\end{proof}























\section{블로업과 드람 코호몰로지}
\label{derham-section-blowing-up}

\noindent
기저 스킴 $S$, 매끄러운 사상 $X \to S$, 그리고 닫힌 부분스킴
$Z \subset X$를 고정하자. 이 부분스킴도 $S$ 위에서 매끄럽다고 하자.
$b : X' \to X$를 $X$의 $Z$를
따른 블로업이라 하자. $E \subset X'$를 예외 인자라 하자. 그림은 다음과
같다.
\begin{equation}
\label{derham-equation-blowup}
\vcenter{
\xymatrix{
E \ar[r]_j \ar[d]_p & X' \ar[d]^b \\
Z \ar[r]^i & X
}
}
\end{equation}
이 절의 목표는 사상
$b^* : H_{dR}^*(X/S) \longrightarrow H_{dR}^*(X'/S)$
이 단사임을 증명하는 것이다(사실 이보다 훨씬 더 많은 것을 말할 수 있다).

\begin{lemma}
\label{derham-lemma-comparison-general}
사상 심화의 보조정리 \ref{more-morphisms-lemma-blowup}의 표기를
사용하자. $a \geq 0$에 대해 다음이 성립한다.
\begin{enumerate}
\item 사상
$\Omega^a_{X/S} \to b_*\Omega^a_{X'/S}$은 동형이다.
\item 사상 $\Omega^a_{Z/S} \to p_*\Omega^a_{E/S}$은 동형이다.
\item 사상 $Rb_*\Omega^a_{X'/S} \to i_*Rp_*\Omega^a_{E/S}$은 차수가
$\geq 1$인 코호몰로지 층들 위에서 동형이다.
\end{enumerate}
\end{lemma}

\begin{proof}
$\epsilon : X_1 \to X$를 전사 에탈 사상이라 하자. 사상 심화의
보조정리 \ref{more-morphisms-lemma-blowup}에서 생각한 대상들을 기저변환하여
얻는 사상과 부분스킴을 각각 $i_1 : Z_1 \to X_1$,
$b_1 : X'_1 \to X_1$, $E_1 \subset X'_1$, 그리고
$p_1 : E_1 \to Z_1$이라 하자. $i_1$은 $S$ 위에서 매끄러운 스킴들
사이의 닫힌 몰입이고, 인자의 보조정리
\ref{divisors-lemma-flat-base-change-blowing-up}에 의해 $b_1$은 중심이
$Z_1$인 블로업임에 유의하자. 사상 $b_1$, $p_1$, $i_1$에 대해 (1),
(2), (3)을 증명할 수 있다고 가정하자. 그러면 보조정리
\ref{derham-lemma-etale}에 의해 (1), (2), (3)의 사상들을 $\epsilon$으로 당겨
얻은 사상들이 동형이다. $\epsilon$은 전사 평탄 사상이므로 결론이 따른다.
따라서 에탈 국소적으로 작업하면, 사상 심화의 보조정리
\ref{more-morphisms-lemma-etale-local-structure}에 의해 절
\ref{derham-section-calculations}에서 논한 상황이라고 가정할 수 있다. 이 경우
이 보조정리는 보조정리 \ref{derham-lemma-comparison}과 같다.
\end{proof}

\begin{lemma}
\label{derham-lemma-distinguished-triangle-blowup}
사상 심화의 보조정리 \ref{more-morphisms-lemma-blowup}의 표기를
사용하고 구조 사상을 $f : X \to S$라 하자. 그러면 표준적인 완전 삼각형
$$
\Omega^\bullet_{X/S} \to
Rb_*(\Omega^\bullet_{X'/S}) \oplus i_*\Omega^\bullet_{Z/S} \to
i_*Rp_*(\Omega^\bullet_{E/S}) \to
\Omega^\bullet_{X/S}[1]
$$
$D(X, f^{-1}\mathcal{O}_S)$ 안에 존재한다. 여기서 네 사상
$$
\begin{matrix}
\Omega^\bullet_{X/S} & \to & Rb_*(\Omega^\bullet_{X'/S}), \\
\Omega^\bullet_{X/S} & \to & i_*\Omega^\bullet_{Z/S}, \\
Rb_*(\Omega^\bullet_{X'/S}) & \to & i_*Rp_*(\Omega^\bullet_{E/S}), \\
i_*\Omega^\bullet_{Z/S} & \to & i_*Rp_*(\Omega^\bullet_{E/S})
\end{matrix}
$$
은 표준적인 사상들이며(절 \ref{derham-section-de-rham-complex}), 그중 하나만
부호를 반대로 취한다.
\end{lemma}

\begin{proof}
완전 삼각형
$$
C \to Rb_*\Omega^\bullet_{X'/S} \oplus i_*\Omega^\bullet_{Z/S}
\to i_*Rp_*\Omega^\bullet_{E/S} \to C[1]
$$
을 $D(X, f^{-1}\mathcal{O}_S)$ 안에서 택하자. 표준적인 사상들과 양립하는
방식으로 $\Omega^\bullet_{X/S}$가 $C$와 동형임을 보이면 충분하다. 삼각
범주의 공리들에 의해 완전 삼각형들 사이의 사상
$$
\xymatrix{
C' \ar[r] \ar[d] &
b_*\Omega^\bullet_{X'/S} \oplus i_*\Omega^\bullet_{Z/S} \ar[r] \ar[d] &
i_*p_*\Omega^\bullet_{E/S} \ar[r] \ar[d] &
C'[1] \ar[d] \\
C \ar[r] &
Rb_*\Omega^\bullet_{X'/S} \oplus i_*\Omega^\bullet_{Z/S} \ar[r] &
i_*Rp_*\Omega^\bullet_{E/S} \ar[r] &
C[1]
}
$$
이 존재한다. 보조정리 \ref{derham-lemma-comparison-general}의 (3)과 유도 범주의 명제 \ref{derived-proposition-9}에 의해 $C' \to C$는
동형이다. 보조정리 \ref{derham-lemma-comparison-general}의 (2)에 의해 사상
$i_*\Omega^\bullet_{Z/S} \to i_*p_*\Omega^\bullet_{E/S}$은 동형이다.
따라서 유도 범주에서 $C' = b_*\Omega^\bullet_{X'/S}$이다. 마지막으로
보조정리 \ref{derham-lemma-comparison-general}의 (1)을 사용하면 이는
$\Omega^\bullet_{X/S}$와 같다. 이것이 표준적인 사상들과 양립함을
확인하는 일은 생략한다.
\end{proof}

\begin{proposition}
\label{derham-proposition-blowup-split}
사상 심화의 보조정리 \ref{more-morphisms-lemma-blowup}의 표기를
사용하자. 사상 $\Omega^\bullet_{X/S} \to Rb_*\Omega^\bullet_{X'/S}$은
$D(X, (X \to S)^{-1}\mathcal{O}_S)$ 안에서 분할된다.
\end{proposition}

\begin{proof}
보조정리 \ref{derham-lemma-distinguished-triangle-blowup}에서 구성한 삼각형을
생각하자. 사상
$$
Rb_*(\Omega^\bullet_{X'/S}) \oplus i_*\Omega^\bullet_{Z/S} \to
i_*Rp_*(\Omega^\bullet_{E/S})
$$
은 그 상이 직합 성분 $i_*\Omega^\bullet_{Z/S}$를 포함하도록 분할된다고
주장한다. 유도 범주의 보조정리 \ref{derived-lemma-split}에 의해,
이 주장은 삼각형의 첫째 화살표가 직합 성분
$i_*\Omega^\bullet_{Z/S}$ 위에서 영이 되는 분할을 가짐을 보이며, 따라서
보조정리를 증명한다. 이 주장을 증명하기 위해
$Rp_*\Omega^\bullet_{E/S}$를 직합으로 분해하겠다. 그 첫째 부분은
$\Omega^\bullet_{Z/S}$에 대응하고, 둘째 부분은
$Rb_*\Omega^\bullet_{X'/S}$을 통해 올릴 수 있다.

\medskip\noindent
주장의 증명. $X$를 $S$에 대한 상대 차원이 일정한 열리고 닫힌
부분스킴들로 분해할 수 있다. 스킴의 사상의 보조정리
\ref{morphisms-lemma-smooth-omega-finite-locally-free}를 보라. 이에 따라
유도 범주 $D(X, f^{-1}\mathcal{O})_S)$도 범주들의 곱으로 분해되므로,
$X$의 상대 차원이 일정한 $N$이라고 가정할 수 있으며, 이는 $S$에 대한
상대 차원이다. $Z = \coprod Z_m$을 상대 차원이 $m \geq 0$인 열리고
닫힌 부분스킴들로 분해할 수 있으며, 상대 차원은 $S$에 대해 취한다.
$i_m : Z_m \to X$를 $i$의 $Z_m$에 대한 제한이라 하자. 이는 여차원
$N - m$인 정칙 몰입이다. 인자의
보조정리
\ref{divisors-lemma-immersion-smooth-into-smooth-regular-immersion}를 보라.
$E = \coprod E_m$을 이에 대응하는 분해라 하자. 즉,
$E_m = p^{-1}(Z_m)$으로 놓는다. $p_m : E_m \to Z_m$이라 쓰자. 이는
$p$를 $E_m$에 제한한 사상이다. 그러면 표준적인 동형
$$
\tilde \xi_m :
\bigoplus\nolimits_{t = 0, \ldots, N - m - 1}
\Omega^\bullet_{Z_m/S}[-2t]
\longrightarrow
Rp_{m, *}\Omega^\bullet_{E_m/S}
$$
$D(Z_m, (Z_m \to S)^{-1}\mathcal{O}_S)$ 안에 존재한다. 여기서 차수
$0$에서는 표준적인 사상
$\Omega^\bullet_{Z_m/S} \to Rp_{m, *}\Omega^\bullet_{E_m/S}$을 얻는다.
주의 \ref{derham-remark-projective-space-bundle-formula}를 보라. 따라서 동형
$$
\tilde \xi :
\bigoplus\nolimits_m
\bigoplus\nolimits_{t = 0, \ldots, N - m - 1}
\Omega^\bullet_{Z_m/S}[-2t]
\longrightarrow
Rp_*(\Omega^\bullet_{E/S})
$$
$D(Z, (Z \to S)^{-1}\mathcal{O}_S)$ 안에 존재한다. 이 동형을 정의역의
직합 성분 $\Omega^\bullet_{Z/S} = \bigoplus \Omega^\bullet_{Z_m/S}$에
제한하면 표준적인 사상
$\Omega^\bullet_{Z/S} \to Rp_*(\Omega^\bullet_{E/S})$이 된다.
부분복합체 $M_m$과 $K_m$을 생각하자. 이들은 주의
\ref{derham-remark-projective-space-bundle-formula}에서 도입한 복합체
$\bigoplus\nolimits_{t = 0, \ldots, N - m - 1} \Omega^\bullet_{Z_m/S}[-2t]$
의 부분복합체들이다. 다음과 같이 놓자.
$$
M = \bigoplus M_m
\quad\text{및}\quad
K = \bigoplus K_m
$$
$M = K[-2]$이고, 구성에 의해 사상
$$
c_{E/Z} \oplus \tilde \xi|_M :
\Omega^\bullet_{Z/S} \oplus M
\longrightarrow
Rp_*(\Omega^\bullet_{E/S})
$$
은 동형이다(위에서 참조한 주의를 보라).

\medskip\noindent
사상
$$
\delta : \Omega^\bullet_{E/S}[-2] \longrightarrow \Omega^\bullet_{X'/S}
$$
을 생각하자. 이는 보조정리 \ref{derham-lemma-gysin-via-log-complex}에서 얻은
$D(X', (X' \to S)^{-1}\mathcal{O}_S)$ 안의 사상이며, 합성
$$
\Omega^\bullet_{E/S}[-2] \longrightarrow \Omega^\bullet_{X'/S}
\longrightarrow
\Omega^\bullet_{E/S}
$$
은 주의 \ref{derham-remark-cup-product-as-a-map}의 사상 $\theta'$이며, 여기서는
$c_1^{dR}(\mathcal{O}_{X'}(-E))|_E) = c_1^{dR}(\mathcal{O}_E(1))$이다. 주의
\ref{derham-remark-projective-space-bundle-formula}의 마지막 명제에 의해 도식
$$
\xymatrix{
K[-2] \ar[d]_{(\tilde \xi|_K)[-2]} \ar[r]_{\text{id}} &
M \ar[d]^{\tilde x|_M} \\
Rp_*\Omega^\bullet_{E/S}[-2] \ar[r]^-{Rp_*\theta'} &
Rp_*\Omega^\bullet_{E/S}
}
$$
은 가환한다. 따라서 주장에 필요한 분할은 다음 사상으로 얻어진다.
\begin{align*}
Rp_*(\Omega^\bullet_{E/S})
& \xrightarrow{(c_{E/Z} \oplus \tilde \xi|_M)^{-1}}
\Omega^\bullet_{Z/S} \oplus M \\
& \xrightarrow{\text{id} \oplus \text{id}^{-1}}
\Omega^\bullet_{Z/S} \oplus K[-2] \\
& \xrightarrow{\text{id} \oplus (\tilde \xi|_K)[-2]}
\Omega^\bullet_{Z/S} \oplus Rp_*\Omega^\bullet_{E/S}[-2] \\
& \xrightarrow{\text{id} \oplus Rb_*\delta}
\Omega^\bullet_{Z/S} \oplus Rb_*\Omega^\bullet_{X'/S}
\end{align*}
위에서 서술한 $\theta'$와 $\delta$의 관계 및 $\theta'$,
$\tilde \xi|_K$, $\tilde \xi|_M$이 등장하는 위의 가환 도식은 이것이
표준적인 사상
$\Omega^\bullet_{Z/S} \oplus Rb_*(\Omega^\bullet_{X'/S}) \to
Rp_*(\Omega^\bullet_{E/S})$의 분할 사상임을 보이는 데 정확히 필요한
내용이다. 이로써 주장의 증명이 끝난다.
\end{proof}

\noindent
보조정리 \ref{derham-lemma-splitting-on-omega-a}는 호지 코호몰로지 위의 분할을
만드는 일이 명제 \ref{derham-proposition-blowup-split}의 결과보다 훨씬
간단함을 보여 준다. 독자에게 다음 절로 건너뛸 것을 강하게 권한다.

\begin{lemma}
\label{derham-lemma-ext-zero}
$i : Z \to X$가 여차원 $c$인 정칙 닫힌 몰입이라고 하자. 그러면
$\Ext^q_{\mathcal{O}_X}(i_*\mathcal{F}, \mathcal{E}) = 0$이다. 여기서
$q < c$이고, $\mathcal{E}$는 $X$ 위에서 국소 자유이며,
$\mathcal{F}$는 임의의 $\mathcal{O}_Z$-가군이다.
\end{lemma}

\begin{proof}
$\Ext$의 국소-대역 스펙트럼 열에 의해 이를 $X$ 위에서 아핀 국소적으로
증명하면 충분하다. 층 코호몰로지의 절 \ref{cohomology-section-ext}를 보라.
따라서 $X = \Spec(A)$라고 가정할 수 있고, 정칙열
$f_1, \ldots, f_c$가 $A$ 안에 존재하여
$Z = \Spec(A/(f_1, \ldots, f_c))$가 된다고 할 수 있다. 또한
$c \geq 1$이라고 가정할 수 있다. 그러면
$f_1 : \mathcal{E} \to \mathcal{E}$은 단사이다. $i_*\mathcal{F}$는
$f_1$에 의해 소멸되므로, 보조정리는 $i = 0$일 때 성립하고 전사
$$
\Ext^{q - 1}_{\mathcal{O}_X}(i_*\mathcal{F}, \mathcal{E}/f_1\mathcal{E})
\longrightarrow
\Ext^q_{\mathcal{O}_X}(i_*\mathcal{F}, \mathcal{E})
$$
을 얻는다. 따라서 이 화살표의 정의역이 영임을 보이면 충분하다. 이제
이 논증을 반복하자. $c \geq 2$이면 사상
$f_2 : \mathcal{E}/f_1\mathcal{E} \to \mathcal{E}/f_1\mathcal{E}$은
단사이다. $i_*\mathcal{F}$는 $f_2$에 의해 소멸되므로, 보조정리는
$q = 1$일 때 성립하고 전사
$$
\Ext^{q - 2}_{\mathcal{O}_X}(i_*\mathcal{F},
\mathcal{E}/f_1\mathcal{E} + f_2\mathcal{E})
\longrightarrow
\Ext^{q - 1}_{\mathcal{O}_X}(i_*\mathcal{F}, \mathcal{E}/f_1\mathcal{E})
$$
을 얻는다. 이와 같이 계속하면 보조정리가 증명된다.
\end{proof}

\begin{lemma}
\label{derham-lemma-splitting-on-omega-a}
사상 심화의 보조정리 \ref{more-morphisms-lemma-blowup}의 표기를
사용하자. $a \geq 0$에 대해 유일한 화살표
$Rb_*\Omega^a_{X'/S} \to \Omega^a_{X/S}$가 $D(\mathcal{O}_X)$ 안에
존재하며, 이를
$\Omega^a_{X/S} \to Rb_*\Omega^a_{X'/S}$와 합성하면
$\Omega^a_{X/S}$ 위의 항등사상이 된다.
\end{lemma}

\begin{proof}
$X$를 $S$에 대한 상대 차원이 일정한 열리고 닫힌 부분스킴들로 분해할
수 있다. 스킴의 사상의 보조정리
\ref{morphisms-lemma-smooth-omega-finite-locally-free}를 보라. 이에 따라
유도 범주 $D(X, f^{-1}\mathcal{O})_S)$도 범주들의 곱으로 분해되므로,
$X$의 상대 차원이 일정한 $N$이라고 가정할 수 있으며, 이는 $S$에 대한
상대 차원이다. $Z = \coprod Z_m$을 상대 차원이 $m \geq 0$인 열리고
닫힌 부분스킴들로 분해할 수 있으며, 상대 차원은 $S$에 대해 취한다.
$i_m : Z_m \to X$를 $i$의 $Z_m$에 대한 제한이라 하자. 이는 여차원
$N - m$인 정칙 몰입이다. 인자의 보조정리
\ref{divisors-lemma-immersion-smooth-into-smooth-regular-immersion}를 보라.
$E = \coprod E_m$을 이에 대응하는 분해라 하자. 즉,
$E_m = p^{-1}(Z_m)$으로 놓는다. 다음과 같은 자연스러운 사상들이
존재한다고 주장한다.
$$
b^*\Omega^a_{X/S} \to \Omega^a_{X'/S} \to
b^*\Omega^a_{X/S} \otimes_{\mathcal{O}_{X'}}
\mathcal{O}_{X'}(\sum (N - m - 1)E_m)
$$
그 합성은 포함
$\mathcal{O}_{X'} \to \mathcal{O}_{X'}(\sum (N - m - 1)E_m)$로부터
유도된다. 이를 증명하려면 첫째 화살표의 여핵이 $X'$ 위에서 국소적으로
유효 카르티에 인자 $\sum (N - m - 1)E_m$의 국소 방정식에 의해 소멸됨을
보이면 충분하다. 이를 확인하기 위해 보조정리
\ref{derham-lemma-comparison-general}의 증명에서처럼 $X$ 위에서 에탈 국소적으로
작업하고 보조정리 \ref{derham-lemma-comparison-bis}를 적용할 수 있다. 보조정리
\ref{derham-lemma-blowup-twist-same-cohomology}를 사용하여 에탈 국소적으로
계산하면, 유도된 합성
$$
\Omega^a_{X/S} \to Rb_*\Omega^a_{X'/S} \to
Rb_*\left(b^*\Omega^a_{X/S} \otimes_{\mathcal{O}_{X'}}
\mathcal{O}_{X'}(\sum (N - m - 1)E_m)\right)
$$
은 $D(\mathcal{O}_X)$ 안에서 동형이다. 이로써 보조정리의 사상이
존재함을 얻는다.

\medskip\noindent
유일성을 보이려면 $\tau_{\geq 1}Rb_*\Omega_{X'/S}$에서
$\Omega^a_{X/S}$으로 가는 $D(\mathcal{O}_X)$ 안의 영이 아닌 사상이
없음을 보이면 충분하다. 이를 위해서는 다시
$R^qb_*\Omega_{X'/s}[-q]$에서 $\Omega^a_{X/S}$으로 가는
$D(\mathcal{O}_X)$ 안의 영이 아닌 사상이 $q \geq 1$에 대해 없음을
보이면 충분하다(세부 사항은 생략한다). 보조정리
\ref{derham-lemma-comparison-general}에 의해
$R^qb_*\Omega_{X'/s} \cong i_*R^qp_*\Omega^a_{E/S}$는
$Z = \coprod Z_m$ 위 가군의 밀어앞이다. 더욱이
$R^qp_*\Omega^a_{E/S}$를 $Z_m$에 제한한 것은 $q < N - m$일 때만
영이 아닐 수 있다. 실제로 $E_m \to Z_m$의 올들은 차원이
$N - m - 1$이고, 스킴의 극한의 보조정리
\ref{limits-lemma-higher-direct-images-zero-above-dimension-fibre}를 적용할
수 있다. 따라서 원하는 소멸은 보조정리 \ref{derham-lemma-ext-zero}에서 따른다.
\end{proof}














\section{미분형식 층의 비교}
\label{derham-section-quasi-finite-syntomic}

\noindent
이 절의 목표는 모든 국소 준유한 신토믹 스킴 사상 $f : Y \to X$에 대해
사상들
$$
c^p_{Y/X} :
\Omega^p_{Y/\mathbf{Z}}
\longrightarrow
f^*\Omega^p_{X/\mathbf{Z}} \otimes_{\mathcal{O}_Y} \det(\NL_{Y/X})
$$
을 모든 $p \geq 0$에 대해 구성하는 것이다. 이 사상들은 다음 성질들을
만족한다.
\begin{enumerate}
\item
\label{derham-item-degree-zero}
$c_{Y/X}^0(1) = \delta(\NL_{Y/X})$이다. 판별식과 디퍼런트 아이디얼의 절
\ref{discriminant-section-tate-map}을 보라.
\item
\label{derham-item-multiplicative}
$c_{Y/X}^{q + p}(\omega \wedge \eta) = \omega \wedge c_{Y/X}^p(\eta)$이다.
여기서 $\omega$는 $f^*\Omega^q_{X/\mathbf{Z}}$의 국소 단면이고,
$\eta$는 $\Omega^p_{Y/\mathbf{Z}}$의 국소 단면이다.
\item
\label{derham-item-base-change}
$c^p_{Y/X}$의 구성은 열린집합으로의 제한 및 기저변환과 양립한다
(보조정리 \ref{derham-lemma-base-change-Garel-upstairs}에서 서술한 의미에서).
\end{enumerate}
이 조건들로부터 사상 $c^p_{Y/X}$이 $\mathcal{O}_Y$-선형이고, 이를
$f^*\Omega^p_{X/\mathbf{Z}} \to \Omega^p_{Y/\mathbf{Z}}$와 합성하면
$\delta(\NL_{Y/X})$를 곱하는 사상이 됨에 유의하자. 사실 성질 (1),
(2), (3)을 갖는 연산자들의 모음 $c^p_{Y/X}$이 유일하게 존재함을
보이겠다. 이 결과는 $f$가 유한일 때 드람 복합체 위의 대각합 사상을
구성하기 위해 절 \ref{derham-section-trace}에서 적용할 것이다.

\begin{lemma}
\label{derham-lemma-funny-map}
$R$를 환이라 하고 다음 가환 도식을 생각하자.
$$
\xymatrix{
0 \ar[r] &
K^0 \ar[r] &
L^0 \ar[r] &
M^0 \ar[r] & 0 \\
& & L^{-1} \ar[u]_\partial \ar@{=}[r] &
M^{-1} \ar[u]
}
$$
이는 $R$-가군들의 도식이고, 위 행은 완전하며, $M^0$과 $M^{-1}$은
같은 계수의 유한 자유 가군이다. $p \geq 0$에 대해 표준적인 사상들
$$
c^p : 
\wedge^p(H^0(L^\bullet))
\longrightarrow
\wedge^p(K^0) \otimes_R \det(M^\bullet)
$$
이 존재하며 다음 성질들을 갖는다.
\begin{enumerate}
\item $c^0(1) = \delta(M^\bullet)$이다.
\item $c^{q + p}(\omega \wedge \eta) = \omega \wedge c^p(\eta)$
이다. 여기서 $\omega \in \wedge^q(K^0)$이고
$\eta \in \wedge^p(H^0(L^\bullet))$이다.
\end{enumerate}
특히 $c^p$를 사상
$\wedge^p(K^0) \to \wedge^p(H^0(L^\bullet))$와 합성하면
$\delta(M^\bullet)$를 곱하는 사상이 된다.
\end{lemma}

\begin{proof}
$M^0$과 $M^{-1}$이 계수 $n$의 자유 가군이라고 하자. 모든
$p \geq 0$에 대해 유일한 전사
$$
\pi_p :
\wedge^{p + n}(L^0)
\longrightarrow
\wedge^p(K^0) \otimes \wedge^n(M^0)
$$
가 존재하며 다음 두 성질을 갖는다. (i)
$
\pi_p(k_1 \wedge \ldots \wedge k_p \wedge l_1 \wedge \ldots \wedge l_n) =
k_1 \wedge \ldots \wedge k_p \otimes
m_1 \wedge \ldots \wedge m_n
$
이다. 여기서 $k_1, \ldots, k_p \in K^0$이고,
$l_1, \ldots, l_n \in L^0$의 상들이
$m_1, \ldots, m_n \in M^0$이다. (ii) $\pi_p$의 핵은 외적
$l_1 \wedge \ldots \wedge l_{p + n}$들 가운데 $> p$개의 $l_j$가
$K^0$에 속하는 것들로 생성되는 부분가군이다.
$e_1, \ldots, e_n$을 기저로 택하자. 이는 $L^{-1} = M^{-1}$의
기저이다. 다음 규칙으로 사상을 정의한다.
$$
\eta \mapsto
\pi_p(\tilde \eta \wedge \partial(e_1) \wedge \ldots \wedge \partial(e_n))
\otimes (e_1 \wedge \ldots \wedge e_n)^{\otimes -1}
$$
여기서 $\eta \in \wedge^p(H^0(L^\bullet))$이고,
$\tilde \eta \in \wedge^p(L^0)$는 $\eta$의 올림이다. 우변의 식은
$\det(M^\bullet) = \wedge^n(M^0) \otimes \wedge^n(M^{-1})^{\otimes -1}$
이므로 의미가 있다. 우변은 $\tilde \eta$의 선택과 무관하다. 실제로
$\partial(\xi) \wedge \partial(e_1) \wedge \ldots \wedge \partial(e_n) = 0$
은 임의의 $\xi \in L^{-1}$에 대해 성립한다. 직접 계산하면 이 사상이
$L^{-1} = M^{-1}$의 기저 선택과도 무관함을 알 수 있다. 대수학 심화의 절 \ref{more-algebra-section-determinants-complexes}에 나오는
$\delta(M^\bullet)$의 정의로부터
$c^0(1) = \delta(M^\bullet)$임은 즉시 따른다. 마지막으로 (2)의 규칙은
정의에서 직접 확인할 수 있다.
\end{proof}

\begin{lemma}
\label{derham-lemma-funny-map-functorial}
$R_1 \to R_2$를 환 준동형이라 하자. $i = 1, 2$에 대해 다음 가환
도식들을 생각하자.
$$
\xymatrix{
0 \ar[r] &
K_i^0 \ar[r] &
L_i^0 \ar[r] &
M_i^0 \ar[r] & 0 \\
& & L_i^{-1} \ar[u]_{\partial_i} \ar@{=}[r] &
M_i^{-1} \ar[u]
}
$$
이는 보조정리 \ref{derham-lemma-funny-map}에서와 같은 $R_i$-가군들의 도식이다.
사상들 $K_1^0 \to K_2^0$, $L_1^j \to L_2^j$,
$M_1^j \to M_2^j$가 있다고 가정하자. 이들은 주어진 환 준동형
$R_1 \to R_2$ 및 표시된 도식들의 사상들과 양립한다고 하자. 사상
$M_1^j \otimes_{R_1} R_2 \to M_2^j$가 동형이면 도식들
$$
\xymatrix{ 
\wedge^p(H^0(L_1^\bullet)) \ar[d]
\ar[r]_-{c^p} &
\wedge^p(K_1^0) \otimes_{R_1} \det(M_1^\bullet) \ar[d] \\
\wedge^p(H^0(L_2^\bullet))
\ar[r]^-{c^p} &
\wedge^p(K_2^0) \otimes_{R_2} \det(M_2^\bullet)
}
$$
은 가환한다.
\end{lemma}

\begin{proof}
이는 보조정리 \ref{derham-lemma-funny-map}의 증명에서 제시한 사상의 명시적
기술로부터 따른다. $M_1^j \otimes_{R_1} R_2 \to M_2^j$가 동형이라는
가정은 $M_1^{-1}$에 사용한 기저가 $M_2^{-1}$의 기저로 감을 보이는 데
필요함에 유의하자. 따라서 두 구성의 양립성을 확인할 때 ``같은'' 기저를
사용할 수 있다.
\end{proof}

\begin{remark}
\label{derham-remark-local-description}
$A$를 환이라 하자. $P = A[x_1, \ldots, x_n]$이라 하자.
$f_1, \ldots, f_n \in P$를 택하고 $B = P/(f_1, \ldots, f_n)$로 놓자.
$A \to B$가 준유한이라고 가정하자. 그러면 $B$는 $A$ 위의 상대 대역
완전 교차이다(가환대수의 정의
\ref{algebra-definition-relative-global-complete-intersection}). 또한
가환대수의 보조정리
\ref{algebra-lemma-relative-global-complete-intersection-conormal}에 의해
$(f_1, \ldots, f_n)/(f_1, \ldots, f_n)^2$는 류
$\overline{f}_i$들로 생성되는 자유 가군이다. 다음 도식을 생각하자.
$$
\xymatrix{
\Omega_{A/\mathbf{Z}} \otimes_A B \ar[r] &
\Omega_{P/\mathbf{Z}} \otimes_P B \ar[r] &
\Omega_{P/A} \otimes_P B \\
&
(f_1, \ldots, f_n)/(f_1, \ldots, f_n)^2 \ar[u] \ar@{=}[r] &
(f_1, \ldots, f_n)/(f_1, \ldots, f_n)^2 \ar[u]
}
$$
오른쪽 열은 $\NL_{B/A}$를 $D(B)$에서 나타내므로 코호몰로지
$\Omega_{B/A}$를 차수 $0$에서 갖는다. 위 행은 분할 짧은 완전열
$0 \to \Omega_{A/\mathbf{Z}} \otimes_A B \to
\Omega_{P/\mathbf{Z}} \otimes_P B \to \Omega_{P/A} \otimes_P B \to 0$이다.
가운데 열은 코호몰로지 $\Omega_{B/\mathbf{Z}}$를 차수 $0$에서 갖는다.
이는 가환대수의 보조정리 \ref{algebra-lemma-differential-seq}에서 따른다. 따라서
보조정리 \ref{derham-lemma-funny-map}에 의해 표준적인 $B$-가군 사상들
$$
\Omega^p_{B/\mathbf{Z}} \longrightarrow
\Omega^p_{A/\mathbf{Z}} \otimes_A \det(\NL_{B/A})
$$
을 얻는다. 이 사상들은 $p \geq 0$에 대해 성질 (1), (2)를 만족하며,
특히 이를 $\Omega^p_{A/\mathbf{Z}} \to \Omega^p_{B/\mathbf{Z}}$와
합성하면 $\delta(\NL_{B/A})$를 곱하는 사상이 된다.
\end{remark}

\begin{lemma}
\label{derham-lemma-base-change-Garel-upstairs}
다음 가환 도식을 생각하자.
$$
\xymatrix{
Y' \ar[d]_{f'} \ar[r]_b & Y \ar[d]^f \\
X' \ar[r]^a & X
}
$$
이 스킴들의 도식은 $Y'$를 $X' \times_X Y$의 열린 부분스킴과 동형으로
만든다. $f$가 국소 준유한이고 신토믹이라고 가정하자. 또한
$c^p_{Y/X} : \Omega^p_{Y/\mathbf{Z}} \to
f^*\Omega^p_{X/\mathbf{Z}} \otimes_{\mathcal{O}_Y} \det(\NL_{Y/X})$가
$p \geq 0$에 대해 주어져 있고
(\ref{derham-item-degree-zero})와 (\ref{derham-item-multiplicative})를 만족한다고 하자.
그러면 사상들의 모음
$c^p_{Y'/X'} : \Omega^p_{Y'/\mathbf{Z}} \to
(f')^*\Omega^p_{X'/\mathbf{Z}} \otimes_{\mathcal{O}_{Y'}} \det(\NL_{Y'/X'})$
은 $p \geq 0$에 대해 많아야 하나 존재한다. 여기서는
(\ref{derham-item-degree-zero})와 (\ref{derham-item-multiplicative})를 만족하고 도식들
$$
\xymatrix{
b^*\Omega^p_{Y/\mathbf{Z}} \ar[rr]_-{b^*c^p_{Y/X}} \ar[d] & &
b^*(f^*\Omega^p_{X/\mathbf{Z}} \otimes_{\mathcal{O}_Y} \det(\NL_{Y/X}))
\ar@{=}[r] &
(f')^*a^*\Omega^p_{X/\mathbf{Z}} \otimes_{\mathcal{O}_Y} b^*\det(\NL_{Y/X}))
\ar[d] \\
\Omega^p_{Y'/\mathbf{Z}} \ar[rrr]^-{c^p_{Y'/X'}} & & &
(f')^*\Omega^p_{X'/\mathbf{Z}} \otimes_{\mathcal{O}_{Y'}} \det(\NL_{Y'/X'})
}
$$
이 모든 $p \geq 0$에 대해 가환할 것을 요구한다. 여기서 수직 화살표에는
절 \ref{derham-section-de-rham-complex}의 사상들
$b^*\Omega^p_{Y/\mathbf{Z}} \to \Omega^p_{Y'/\mathbf{Z}}$ 및
$a^*\Omega^p_{X/\mathbf{Z}} \to \Omega^p_{X'/\mathbf{Z}}$와,
판별식과 디퍼런트 아이디얼의 절 \ref{discriminant-section-tate-map}의 식별
$b^*\det(\NL_{Y/X}) = \det(\NL_{Y'/X'})$을 사용한다.
\end{lemma}

\begin{proof}
사상
$$
(f')^*\Omega_{X'/\mathbf{Z}} \oplus b^*\Omega_{Y/\mathbf{Z}}
\longrightarrow
\Omega_{Y'/\mathbf{Z}}
$$
은 $Y'$가 올곱의 열린 부분스킴이므로 전사이다. 이에 대응하는 대수 명제는
$\Omega_{B \otimes_A C/\mathbf{Z}}$가 $\Omega_{B/\mathbf{Z}}$와
$\Omega_{C/\mathbf{Z}}$의 상들에 의해 가군으로 생성된다는 것이다.
따라서 모든 $p$에 대해 사상
$$
\bigoplus\nolimits_{i = 0, \ldots, p}
(f')^*\Omega^i_{X'/\mathbf{Z}} \otimes
b^*\Omega^{p - i}_{Y/\mathbf{Z}}
\longrightarrow
\Omega^p_{Y'/\mathbf{Z}}
$$
은 전사이다. 조건 (\ref{derham-item-degree-zero}),
(\ref{derham-item-multiplicative})와 보조정리의 도식이 가환한다는 사실을 함께
사용하면 사상 $c^p_{Y'/X'}$이 결정됨을 알 수 있다(그러나 존재성은
즉시 따르지 않는다).
\end{proof}

\begin{lemma}
\label{derham-lemma-existence-base-change-Garel-upstairs}
다음 가환 도식을 생각하자.
$$
\xymatrix{
Y' \ar[d]_{f'} \ar[r]_b & Y \ar[d]^f \\
X' \ar[r]^a & X
}
$$
이 스킴들의 도식은 $Y'$를 $X' \times_X Y$의 열린 부분스킴과 동형으로
만든다. $f$가 국소 준유한이고 신토믹이라고 가정하자. 그러면 모든
$y' \in Y'$에 대해 열린집합들 $V' \subset Y'$, $V \subset Y$,
$U' \subset X$, $U \subset X$를 찾을 수 있다. 이들은 $y' \in V'$,
$f'(V') \subset U'$, $b(V') \subset V$, $a(U') \subset U$,
$f(V) \subset U$를 만족한다. 또한 $p \geq 0$에 대해
(\ref{derham-item-degree-zero})와 (\ref{derham-item-multiplicative})를 만족하는 사상들
$c^p_{V/U}$와 $c^p_{V'/U'}$가 존재하며, 도식
$$
\xymatrix{
V' \ar[d] \ar[r]_b & V \ar[d] \\
U' \ar[r] & U
}
$$
과 보조정리 \ref{derham-lemma-base-change-Garel-upstairs}에서 설명한 의미로
양립한다.
\end{lemma}

\begin{proof}
이를 증명할 때 $Y'$를 $X' \times_X Y$로 바꾸어도 됨은 분명하다.
판별식과 디퍼런트 아이디얼의 보조정리
\ref{discriminant-lemma-syntomic-quasi-finite}의 (5)와 같이 아핀
열린집합 $V \subset Y$와 $U \subset X$를 택하되, $V$가 $y'$의 상을
포함하도록 하자. $U' \subset X'$를 $y'$의 상을 포함하고 $U$ 안으로
가는 아핀 열린집합으로 택하자. $X, Y, X', Y'$를 각각
$U, V, U', U' \times_U V$로 바꾸면 다음 문단의 상황이라고 가정할 수
있다.

\medskip\noindent
$X' = \Spec(A')$, $X = \Spec(A)$, $Y = \Spec(B)$이고
$Y' = \Spec(A' \otimes_A B)$라고 가정하자. 여기서
$B = A[x_1, \ldots, x_n]/(f_1, \ldots, f_n)$는 $A$ 위의 상대 대역
완전 교차이다. 그러면
$B' = A'[x_1, \ldots, x_n]/(f'_1, \ldots, f'_n)$도 $A'$ 위의 상대 대역
완전 교차이다. 여기서 $f'_i$는 $f_i$의
$A'[x_1, \ldots, x_n]$ 안에서의 상이다. 가환대수의 보조정리
\ref{algebra-lemma-base-change-relative-global-complete-intersection}를 보라.
주의 \ref{derham-remark-local-description}의 구성은 사상들 $c^p_{V/U}$와
$c^p_{V'/U'}$를 준다. 이 사상들은 보조정리
\ref{derham-lemma-funny-map-functorial} 및 $B$와 $B'$의 $A$와 $A'$ 위 표시들의
양립성에 의해 서로 양립한다. 몇 가지 세부 사항은 생략한다.
\end{proof}

\begin{lemma}
\label{derham-lemma-Garel-upstairs}
모든 국소 준유한 신토믹 스킴 사상 $f : Y \to X$에 $\mathcal{O}_Y$-가군
사상들
$$
c^p_{Y/X} :
\Omega^p_{Y/\mathbf{Z}}
\longrightarrow
f^*\Omega^p_{X/\mathbf{Z}} \otimes_{\mathcal{O}_Y} \det(\NL_{Y/X})
$$
을 대응시키는 유일한 규칙이 존재한다. 이 사상들은 $p \geq 0$에 대해
(\ref{derham-item-degree-zero}), (\ref{derham-item-multiplicative}),
(\ref{derham-item-base-change})를 만족한다. 특히 $c^p_{Y/X}$을
$f^*\Omega^p_{X/\mathbf{Z}} \to \Omega^p_{Y/\mathbf{Z}}$와 합성하면
$\delta(\NL_{Y/X})$를 곱하는 사상이 된다.
\end{lemma}

\begin{proof}
이 증명은 판별식과 디퍼런트 아이디얼의 명제
\ref{discriminant-proposition-tate-map}의 증명과 매우 유사하다. 독자는
먼저 그 증명을 살펴보기를 권한다.

\medskip\noindent
명제를 다시 서술해 보자. 범주 $\mathcal{C}$를 다음과 같이 정의한다.
그 대상은 $Y/X$로 표기하는 국소 준유한 신토믹 스킴 사상
$f : Y \to X$이고, 사상 $b/a : Y'/X' \to Y/X$은 가환 도식
$$
\xymatrix{
Y' \ar[d]_{f'} \ar[r]_b & Y \ar[d]^f \\
X' \ar[r]^a & X
}
$$
으로 주어지며, 이 도식은 $Y'$를 $X' \times_X Y$의 열린 부분스킴과
동형으로 만든다. 보조정리의 뜻은 다음과 같다. $Y/X$를
$\mathcal{C}$의 임의의 대상이라 하면 성질 (\ref{derham-item-degree-zero})와
(\ref{derham-item-multiplicative})를 갖는 사상들 $c^p_{Y/X}$,
$p \geq 0$가 있다. 또한 $b/a : Y'/X' \to Y/X$를
$\mathcal{C}$의 임의의 사상이라 하면 $c^p_{Y/X}$와 $c^p_{Y'/X'}$는
보조정리 \ref{derham-lemma-base-change-Garel-upstairs}에서와 같이 양립한다.

\medskip\noindent
$Y/X$가 $\mathcal{C}$에 속하고 $y \in Y$가 주어졌다고 하자. 그러면 아핀
열린집합 $V \subset Y$와 $U \subset X$를 $f(V) \subset U$가 되도록
찾을 수 있고, 사상들
$$
\Omega^p_{Y/\mathbf{Z}}|_V
\longrightarrow
\left(
f^*\Omega^p_{X/\mathbf{Z}} \otimes_{\mathcal{O}_Y} \det(\NL_{Y/X})
\right)|_V
$$
이 성질 (\ref{derham-item-degree-zero})와 (\ref{derham-item-multiplicative})를
만족하도록 할 수 있다. 이는 판별식과 디퍼런트 아이디얼의 보조정리
\ref{discriminant-lemma-syntomic-quasi-finite}의 (5)에서와 같이 아핀
열린집합들을 택하고 주의 \ref{derham-remark-local-description}의 구성을
사용하면 따른다.

\medskip\noindent
$f$의 에탈 점들은 정확히 $\delta(\NL_{Y/X})$가 영이 되지 않는 점들이다.
판별식과 디퍼런트 아이디얼의 절 \ref{discriminant-section-tate-map}의 논의를 보라.
특히 사상
$f^*\Omega^p_{X/\mathbf{Z}} \to \Omega^p_{Y/\mathbf{Z}}$은
$\delta(\NL_{Y/X})$를 가역화한 뒤 동형이 된다. 따라서
$\Omega^p_{X/\mathbf{Z}}$가 유한 국소 자유이고
$\delta(\NL_{Y/X})$의 $\mathcal{O}_Y$ 안에서의 소멸자가 영이면, 앞
문단에서 구성한 국소 사상들은
유일하며 자동으로 이어붙는다.

\medskip\noindent
$\mathcal{C}_{nice} \subset \mathcal{C}$를 다음 조건을 만족하는 $Y/X$들로
이루어진 충만한 부분범주라 하자.
\begin{enumerate}
\item $\Omega_{X/\mathbf{Z}}$는 국소 자유이다.
\item $\delta(\NL_{Y/X})$의 $\mathcal{O}_Y$ 안에서의 소멸자는 영이다.
\end{enumerate}
앞 문단의 관찰에 의해, 임의의
대상 $Y/X$가 $\mathcal{C}_{nice}$에 속할 때 조건
(\ref{derham-item-degree-zero})와 (\ref{derham-item-multiplicative})를 만족하는 유일한
사상들 $c^p_{Y/X}$, $p \geq 0$가 존재한다. 사상
$b/a : Y'/X' \to Y/X$가 $\mathcal{C}_{nice}$에 속하면, 사상들
$c^p_{Y/X}$와 $c^p_{Y'/X'}$는 보조정리
\ref{derham-lemma-base-change-Garel-upstairs}에서와 같이 양립한다. 실제로 보조정리
\ref{derham-lemma-existence-base-change-Garel-upstairs}에 의해 국소적으로 양립하는
사상들이 존재하며, 방금 언급한 유일성에 의해 이 사상들은
$c^p_{Y/X}$ 및 $c^p_{Y'/X'}$와 일치한다. 다시 말해 충만한 부분범주
$\mathcal{C}_{nice}$ 위에서 문제를 해결하였다. $Y/X$가
$\mathcal{C}_{nice}$에 속할 때, 방금 얻은 해를 계속 $c^p_{Y/X}$로
표기하겠다.

\medskip\noindent
사실 좀 더 일반적으로, $b/a : Y'/X' \to Y/X$가 $\mathcal{C}$의
사상으로 주어져 있고 $Y/X$가
$\mathcal{C}_{nice}$의 대상이라고 하자. 이번에는 보조정리
\ref{derham-lemma-base-change-Garel-upstairs}의 유일성을 사용하는 같은 논증에
의해, 조건 (\ref{derham-item-degree-zero})와 (\ref{derham-item-multiplicative})를
만족하는 사상들 $c^p_{Y'/X'}$, $p \geq 0$가 존재하며 이미 구성한
사상들 $c^p_{Y/X}$와 양립한다. 이를 기저변환
$(b/a)^*c^p_{Y/X}$라 부르자.

\medskip\noindent
사상들
$$
Y_1/X_1 \xleftarrow{b_1/a_1} Y/X \xrightarrow{b_2/a_2} Y_2/X_2
$$
을 $\mathcal{C}$ 안에서 생각하자. 여기서 $Y_1/X_1$과 $Y_2/X_2$는
$\mathcal{C}_{nice}$의 대상이다.
{\bf 주장.} 두 기저변환 $(b_i/a_i)^*c^p_{Y_i/X_i}$,
$i = 1, 2$는 같다. 먼저 이 주장이 보조정리를 함의함을 보인 다음 주장을
증명하겠다.

\medskip\noindent
$d, n \geq 1$이라 하고, 판별식과 디퍼런트 아이디얼의 예
\ref{discriminant-example-universal-quasi-finite-syntomic}에서 구성한 국소
준유한 신토믹 사상 $Y_{n, d} \to X_{n, d}$를 생각하자. 그러면
$Y_{n, d}$와 $Y_{n, d}$는 $\mathbf{Z}$ 위에서 유한형이고 매끄러운
기약 스킴이다. 실제로 $X_{n, d}$는 $\mathbf{Z}$ 위 다항식환의
스펙트럼이고, $Y_{n, d}$는 그러한 스펙트럼의 열린 부분스킴이다. 사상
$Y_{n, d} \to X_{n, d}$는 국소 준유한 신토믹이고 조밀한 열린집합
위에서 에탈이다. 판별식과 디퍼런트 아이디얼의 보조정리
\ref{discriminant-lemma-universal-quasi-finite-syntomic-etale}를 보라.
따라서 $\delta(\NL_{Y_{n, d}/X_{n, d}})$는 영이 아니다. 예를 들어
판별식과 디퍼런트 아이디얼의 주의
\ref{discriminant-remark-local-description-delta}에는
$\delta(\NL_{Y/X})$의 국소 기술이 있고, 스킴의 사상의 보조정리
\ref{morphisms-lemma-etale-at-point}의 (8)에는 에탈 사상의 국소 기술이
있다. 이제 기약 정칙 스킴 위 가역 가군의 영이 아닌 단면은 소멸자가
영이다. 따라서 $Y_{n, d}/X_{n, d}$는 $\mathcal{C}_{nice}$의 대상이다.

\medskip\noindent
$Y/X$를 $\mathcal{C}$의 임의의 대상이라 하고 $y \in Y$라 하자.
판별식과 디퍼런트 아이디얼의 보조정리
\ref{discriminant-lemma-locally-comes-from-universal}에 의해
$n, d \geq 1$과 사상들
$$
Y/X \leftarrow V/U \xrightarrow{b/a} Y_{n, d}/X_{n, d}
$$
을 $\mathcal{C}$ 안에서 찾을 수 있다. 여기서 $V \subset Y$와
$U \subset X$는 열린집합이다. 그러면 기저변환
$c^p_{V/U} = (b/a)^*c^p_{Y_{n, d}/X_{n, d}}$를 얻는다. 주장에 의해
이렇게 국소적으로 구성한 사상들 $c^p_{V/U}$는 이어붙는다. 따라서 성질
(\ref{derham-item-degree-zero})와 (\ref{derham-item-multiplicative})를 갖는 잘 정의된
대역 사상들 $c^p_{Y/X}$을 얻는다. 마지막으로
$b/a : Y'/X' \to Y/X$를 $\mathcal{C}$의 임의의 사상이라 하자. 이
문단에서 구성한 사상들 $c^p_{Y'/X'}$, $p \geq 0$와
$c^p_{Y/X}$, $p \geq 0$가 있다. 이들이 보조정리
\ref{derham-lemma-base-change-Garel-upstairs}에서와 같이 양립함을 확인할 때는
$Y'/X'$와 $Y/X$ 위에서 국소적으로 작업해도 된다. 따라서 사상
$Y/X \to Y_{n, d}/X_{n, d}$가 존재한다고 가정할 수 있다. 구성에 의해
모음 $c^p_{Y'/X'}$와 모음 $c^p_{X/Y}$는 둘 다
$Y_{n, d}/X_{n, d}$으로 가는 사상과 양립한다. 이 사실과 보조정리
\ref{derham-lemma-base-change-Garel-upstairs}의 유일성으로부터
$c^p_{Y'/X'}$가 $c^p_{Y/X}$와 양립함은 쉽게 따른다. 이제 주장을
증명하는 일만 남았다.

\medskip\noindent
남은 증명에서는 주장을 증명한다. 점 $y \in Y$를 택하고 이 사상들이
$y$의 열린 근방에서 일치함을 보이면 된다. 따라서 $Y_1$, $Y_2$를
$y$의 $Y_1$, $Y_2$ 안에서의 상들의 열린 근방으로 바꿀 수 있다. 그러므로
$Y, X, Y_1, X_1, Y_2, X_2$가 아핀이라고 가정할 수 있고, 이들이 각각
환 $B, A, B_1, A_1, B_2, A_2$의 스펙트럼이라고 하자. 그림은 다음과
같다.
$$
\xymatrix{
B_1 \ar[r] & B & B_2 \ar[l] \\
A_1 \ar[u] \ar[r] & A \ar[u] & A_2 \ar[l] \ar[u]
}
$$
가정에 의해 $B$의 스펙트럼은 $A \otimes_{A_1} B_1$의 스펙트럼과
$A \otimes_{A_2} B_2$의 스펙트럼 모두의 아핀 열린집합이다. 더 축소하면
원소들 $g_i \in A \otimes_{A_i} B_i$가 존재하고 우리의 사상들이
동형들 $(A \otimes_{A_i} B_i)_{g_i} = B$를 준다고 가정할 수 있다.
스킴의 성질의 보조정리 \ref{properties-lemma-standard-open-two-affines}를
보라. $x_\alpha$, $y_\beta$를 충분히 많은 변수들의 모음으로 택하여
전사들
$$
A'_1 = A_1[x_\alpha] \to A
\quad\text{및}\quad
A'_2 = A_2[y_\beta] \to A
$$
을 각각 $A_1$-대수 및 $A_2$-대수의 사상으로 택할 수 있게 하자. 그러면
$h_i \in A'_i \otimes_{A_i} B_i$를 택할 수 있다. 이는
$g_i \in A \otimes_{A_i} B_i$의 올림이다. 이제 도식
$$
\xymatrix{
(A'_1 \otimes_{A_1} B_1)_{h_1} \ar[r] & B &
(A'_2 \otimes_{A_1} B_2)_{h_2} \ar[l] \\
A'_1 \ar[u] \ar[r] & A \ar[u] & A'_2 \ar[l] \ar[u]
}
$$
을 생각하자. 구성에 의해 두 사각형은 쌍대카르테시안이다. 다음 환 사상을
생각하자.
$$
A' = A'_1 \times_A A'_2 \longrightarrow
B' =
(A'_1 \otimes_{A_1} B_1)_{h_1} \times_B
(A'_2 \otimes_{A_1} B_2)_{h_2}
$$
대수학 심화의 보조정리
\ref{more-algebra-lemma-module-over-fibre-product}에 의해
$A'_1 \otimes_{A'} B' = (A'_1 \otimes_{A_1} B_1)_{h_1}$이고
$A'_2 \otimes_{A'} B' = (A'_2 \otimes_{A_2} B_2)_{h_2}$이다. 특히 사상
$Y' = \Spec(B') \to \Spec(A') = X'$의 올들은 사상들
$Y_i \to X_i$의 올들을 기저변환한 것의 열린 부분스킴이므로 유한 완전
국소 교차이다(판별식과 디퍼런트 아이디얼의 보조정리
\ref{discriminant-lemma-syntomic-quasi-finite}를 보라). 대수학 심화의
보조정리 \ref{more-algebra-lemma-properties-algebras-over-fibre-product}에
의해 환 사상 $A' \to B'$는 평탄하고 유한 표시이다. 따라서
판별식과 디퍼런트 아이디얼의 보조정리
\ref{discriminant-lemma-syntomic-quasi-finite}의 (3)에 의해
$Y'/X'$는 $\mathcal{C}$의 대상이다. 이제 가환 도식
$$
\xymatrix{
& & Y/X \ar[ld] \ar@/_2pc/[lldd]_{b_1/a_1} \ar[rd] \ar@/^2pc/[rrdd]^{b_2/a_2} \\
& Y'_1/X'_1 \ar[rd] \ar[ld] & & Y'_2/X'_2 \ar[ld] \ar[rd] \\
Y_1/X_1 & & Y'/X' & & Y_2/X_2
}
$$
을 생각하자. 여기서 $Y'_i/X'_i$는
$A'_i \to (A'_i \otimes_{A_i} B_i)_{h_i}$에 대응한다.
$Y'_i/X'_i$는 $\mathcal{C}_{nice}$의 대상임에 유의하자. 실제로 무한
차원 아핀 공간과 $Y_i/X_i$의 곱의 열린 부분스킴을 취하여 얻어진다.
작은 세부 사항은 생략한다. 특히 $c^p_{Y_i/X_i}$를 $(b_i/a_i)$로
당긴 것은 $c^p_{Y'_i/X'_i}$를 $Y/X \to Y'_i/X'_i$로 당긴 것과 같다.
이제 $Y'/X'$가 $\mathcal{C}_{nice}$의 대상이라면 증명이 끝난다. 그러나
이는 거의 언제나 성립하지 않는다. 실제로 이 경우
$c^p_{Y'/X'}$를 사각형의 양쪽 변을 따라 당기면 원하는 결론을 얻는다.
$Y'/X'$가 $\mathcal{C}_{nice}$에 속하지 않는 문제를 피하기 위해 위의
논증을 사용한다. 필요하면 스킴들 $X, Y, X'_1, Y'_1, X'_2, Y'_2, X', Y'$
모두를 축소한 뒤 어떤 $n, d \geq 1$을 찾아 도식을 다음과 같이 확장할
수 있다.
$$
\xymatrix{
& Y/X \ar[ld] \ar[rd] \\
Y'_1/X'_1 \ar[rd] & & Y'_2/X'_2 \ar[ld] \\
& Y'/X' \ar[d] \\
& Y_{n, d}/X_{n, d}
}
$$
그런 다음 $c^p_{Y_{n, d}/X_{n, d}}$에서 당겨 오는 앞의 논증을 사용할
수 있다. 이로써 증명이 끝난다.
\end{proof}








\section{드람 복합체 위의 대각합 사상}
\label{derham-section-trace}

\noindent
이 절의 내용 일부에 대한 참고문헌은 \cite{Garel}이다. $S$를 스킴이라
하자. $f : Y \to X$를 $S$ 위 스킴들의 유한 국소 자유 사상이라 하자.
그러면 대각합 사상 $\text{Trace}_f : f_*\mathcal{O}_Y \to \mathcal{O}_X$가
존재한다. 판별식과 디퍼런트 아이디얼의 절 \ref{discriminant-section-discriminant}를
보라. 이 상황에서 드람 복합체 위의 대각합 사상이란 복합체들의 사상
$$
\Theta_{Y/X} : f_*\Omega^\bullet_{Y/S} \longrightarrow \Omega^\bullet_{X/S}
$$
으로서 $\Theta_{Y/X}$가 $\text{Trace}_f$와 차수 $0$에서 같고 다음을
만족하는 것이다.
$$
\Theta_{Y/X}(\omega \wedge \eta) = \omega \wedge \Theta_{Y/X}(\eta)
$$
여기서 $\omega$는 $\Omega^\bullet_{X/S}$의 국소 단면이고, $\eta$는
$f_*\Omega^\bullet_{Y/S}$의 국소 단면이다. 이러한 대각합 사상
$\Theta_{Y/X}$가 모든 유한 국소 자유 사상 $Y \to X$에 대해 존재하는지는
분명하지 않다. 반례나 증명이 있다면
\href{mailto:stacks.project@gmail.com}{stacks.project@gmail.com}
으로 이메일을 보내기 바란다.

\begin{example}
\label{derham-example-no-trace}
드람 복합체 위의 대각합 사상이 존재하지 않는 한 예를 들자.
$\mathbf{C}$-대수 $B = \mathbf{C}[x, y]$를 생각하되,
$G = \{\pm 1\}$가 $x \mapsto -x$ 및 $y \mapsto -y$로 작용한다고 하자.
불변식들 $A = B^G$는 $\mathbf{C}$ 위 유한형 정규 정역을 이루며
$x^2, xy, y^2$로 생성된다. 포함 $A \subset B$에 대해서는 적절한
대각합 사상
$\Omega_{B/\mathbf{C}} \to \Omega_{A/\mathbf{C}}$
이 $1$-형식 위에 존재하지 않는다고 주장한다. 실제로 원소
$\omega = x \text{d} y \in \Omega_{B/\mathbf{C}}$를 생각하자.
$\omega$는 $G$의 작용 아래 불변이므로 ``적절한'' 대각합 사상이
존재한다면 $2\omega$는
$\Omega_{A/\mathbf{C}} \to \Omega_{B/\mathbf{C}}$의 상에 속해야 한다.
그러나 그렇지 않다. $2\omega$를 $\text{d}(x^2)$, $\text{d}(xy)$,
$\text{d}(y^2)$의 선형결합으로 쓸 수 없으며, 이는 계수를 $B$에서
택하도록 허용해도 마찬가지다. 이 예는 \cite{Zannier}의
주요 정리와 모순된다.
\end{example}

\begin{lemma}
\label{derham-lemma-Garel}
모든 유한 신토믹 스킴 사상 $f : Y \to X$에 $\mathcal{O}_X$-가군
사상들
$$
\Theta^p_{Y/X} :
f_*\Omega^p_{Y/\mathbf{Z}}
\longrightarrow
\Omega^p_{X/\mathbf{Z}}
$$
을 대응시키는 유일한 규칙이 존재하며, 다음 성질들을 만족한다.
\begin{enumerate}
\item 다음 사상과 합성하면
$\Omega^p_{X/\mathbf{Z}} \otimes_{\mathcal{O}_X} f_*\mathcal{O}_Y
\to f_*\Omega^p_{Y/\mathbf{Z}}$
$\text{id} \otimes \text{Trace}_f$와 같아진다. 여기서
$\text{Trace}_f : f_*\mathcal{O}_Y \to \mathcal{O}_X$는
판별식과 디퍼런트 아이디얼의 절 \ref{discriminant-section-discriminant}의 사상이다.
\item 이 규칙은 기저변환과 양립한다.
\end{enumerate}
\end{lemma}

\begin{proof}
먼저 $X$가 국소 뇌터라고 가정하자. 보조정리
\ref{derham-lemma-Garel-upstairs}에 의해 표준적인 사상
$$
c^p_{Y/X} : \Omega_{Y/\mathbf{Z}}^p
\longrightarrow
f^*\Omega_{X/\mathbf{Z}}^p \otimes_{\mathcal{O}_Y} \det(\NL_{Y/X})
$$
을 얻는다. 판별식과 디퍼런트 아이디얼의 명제 \ref{discriminant-proposition-tate-map}에
의해 표준적인 동형
$$
c_{Y/X} : \det(\NL_{Y/X}) \to \omega_{Y/X}
$$
을 얻으며, 이는 $\delta(\NL_{Y/X})$를 $\tau_{Y/X}$로 보낸다. 이 사상들을
결합하면
$$
c^p_{Y/X} \otimes c_{Y/X} :
\Omega_{Y/\mathbf{Z}}^p
\longrightarrow
f^*\Omega_{X/\mathbf{Z}}^p \otimes_{\mathcal{O}_Y} \omega_{Y/X}
$$
을 얻는다. 판별식과 디퍼런트 아이디얼의 절 \ref{discriminant-section-finite-morphisms}에
의하면 이는 사상
$$
\Theta_{Y/X}^p :
f_*\Omega_{Y/\mathbf{Z}}^p
\longrightarrow
\Omega_{X/\mathbf{Z}}^p
$$
과 같은 것이다. $c^p_{Y/X} \otimes c_{Y/X}$와 $\Theta_{Y/X}^p$의
관계에는 평가 사상 $f_*\omega_{Y/X} \to \mathcal{O}_X$가 쓰이며, 이
사상은 $\tau_{Y/X}$를 $\text{Trace}_f(1)$로 보낸다는 것을 상기하자.
판별식과 디퍼런트 아이디얼의 절 \ref{discriminant-section-finite-morphisms}를 보라.
따라서 성질 (1)이 성립한다. $X' \to X$에 의한 기저변환에서 $X'$가
국소 뇌터이면, $c^p_{Y/X}$와 $c_{Y/X}$가 모두 이러한 기저변환과
양립하므로 성질 (2)도 성립한다. $f : Y \to X$가 유한 신토믹이고
$X$가 국소 뇌터인 경우, 방금 얻은 해를 계속 $\Theta^p_{Y/X}$로
표기하겠다.

\medskip\noindent
유일성. 유한 신토믹 사상 $f: Y \to X$가 주어졌다고 하자. 여기서
$X$는 $\Spec(\mathbf{Z})$ 위에서 매끄럽고, $f$는 $X$의 조밀한
열린집합 위에서 에탈이라고 하자. 이 경우 $\Theta^p_{Y/X}$는 성질
(1)에 의해 유일하게 결정된다고 주장한다. 실제로 사상들
$$
\Omega^p_{X/\mathbf{Z}} \otimes_{\mathcal{O}_X} f_*\mathcal{O}_Y \to
f_*\Omega^p_{Y/\mathbf{Z}} \to
\Omega^p_{X/\mathbf{Z}}
$$
을 생각하자. 층 $\Omega^p_{X/\mathbf{Z}}$는 가정한 매끄러움에 의해
꼬임이 없다. 따라서 $\Theta^p_{Y/X}$를 $f$가 에탈인 조밀한 열린집합에
제한한 것이 유일하게 결정됨을 확인하면 충분하다. 즉, $f$가 에탈이라고
가정할 수 있다. 그런데 $f$가 에탈이면
$f^*\Omega_{X/\mathbf{Z}} = \Omega_{Y/\mathbf{Z}}$이므로 표시된 식의 첫째
화살표는 동형이다. 합성을 이미 결정했으므로 유일성이 따른다.

\medskip\noindent
$f : Y \to X$를 국소 뇌터 스킴들의 유한 신토믹 사상이라 하자.
$x \in X$라 하자. 판별식과 디퍼런트 아이디얼의 보조정리
\ref{discriminant-lemma-locally-comes-from-universal-finite}
에 의해 $d \geq 1$과 가환 도식
$$
\xymatrix{
Y \ar[d] &
V \ar[d] \ar[l] \ar[r] &
V_d \ar[d] \\
X &
U \ar[l] \ar[r] &
U_d
}
$$
을 찾을 수 있다. 여기서 $x \in U \subset X$는 열린집합이고,
$V = f^{-1}(U)$이며, $V = U \times_{U_d} V_d$이다. 따라서
$\Theta^p_{Y/X}|_V$는 사상 $\Theta^p_{V_d/U_d}$를 당겨 얻는다. 그러나
위의 유일성 논의와 판별식과 디퍼런트 아이디얼의 보조정리들
\ref{discriminant-lemma-universal-finite-syntomic-smooth}와
\ref{discriminant-lemma-universal-finite-syntomic-etale}
에 의해 사상 $\Theta^p_{V_d/U_d}$는 조건 (1)에 따라 유일하게 결정된다.
따라서 유일성이 성립한다.

\medskip\noindent
이제 모든 유한 신토믹 사상 $Y \to X$ 중 $X$가 국소 뇌터인 것에 대해
존재성과 유일성을 얻었다. 보조정리 \ref{derham-lemma-Garel-upstairs}의 증명과
비슷한 논증을 제시하여 일반적인 $X$로 확장할 수도 있다. 그러나 절대
뇌터 근사를 직접 사용하여 증명을 끝낼 수 있다. 실제로
$\Theta^p_{Y/X}$을 구성할 때는 $X$ 위에서 자리스키 국소적으로 구성하면
충분하다(유일성도 보인다는 전제 아래). 따라서 $X$가 아핀이라고 가정할
수 있다(작은 세부 사항은 생략한다). 그러면 $X = \lim_{i \in I} X_i$를
뇌터 아핀 스킴들의 유향 집합 $I$ 위 극한으로 쓸 수 있다. 가환대수의
보조정리 \ref{algebra-lemma-colimit-category-fp-algebras}에 의해
$0 \in I$와 아핀 스킴들의 유한 표시 사상 $f_0 : Y_0 \to X_0$를 찾아,
이를 $X$로 기저변환하면 $Y \to X$가 되게 할 수 있다. $0$을 증가시킨
뒤에는 $Y_0 \to X_0$가 유한 신토믹이라고 가정할 수 있다. 가환대수의
보조정리 \ref{algebra-lemma-colimit-lci}와
\ref{algebra-lemma-colimit-finite}를 보라. $i \geq 0$에 대해서도 기저변환
$f_i : Y_i = Y_0 \times_{X_0} X_i \to X_i$는 유한 신토믹이다. 이제
$$
\Gamma(X, f_*\Omega^p_{Y/\mathbf{Z}}) =
\Gamma(Y, \Omega^p_{Y/\mathbf{Z}}) =
\colim_{i \geq 0} \Gamma(Y_i, \Omega^p_{Y_i/\mathbf{Z}}) =
\colim_{i \geq 0} \Gamma(X_i, f_{i, *}\Omega^p_{Y_i/\mathbf{Z}})
$$
따라서 $\Theta^p_{Y/X}$을 사상들 $\Theta^p_{Y_i/X_i}$의 쌍대극한으로
정의할 수 있고, 또 그렇게 정의해야 한다. 이 사상은 임의의 카르테시안
도식
$$
\xymatrix{
Y' \ar[r] \ar[d] & Y \ar[d] \\
X' \ar[r] & X
}
$$
과 양립한다. 여기서 $X'$는 아핀이다. 위의 논증으로 뇌터 아핀 스킴의
경우에 이를 알고 있기 때문이다. 작은 세부 사항은 생략한다. 힌트로,
$X' = \lim_{j \in J} X'_j$라고도 쓰면 모든 $i \in I$에 대해
$j \in J$와 사상 $X'_j \to X_i$가 존재하며, 이는 사상
$X' \to X$와 양립한다. 이로써 증명이 끝난다.
\end{proof}

\begin{proposition}
\label{derham-proposition-Garel}
\begin{reference}
\cite{Garel}
\end{reference}
$f : Y \to X$를 스킴들의 유한 신토믹 사상이라 하자. 보조정리
\ref{derham-lemma-Garel}의 사상들 $\Theta^p_{Y/X}$는 복합체들의 사상
$$
\Theta_{Y/X} :
f_*\Omega^\bullet_{Y/\mathbf{Z}}
\longrightarrow
\Omega^\bullet_{X/\mathbf{Z}}
$$
을 정의하며, 다음 성질들을 갖는다.
\begin{enumerate}
\item 차수 $0$에서
$\text{Trace}_f : f_*\mathcal{O}_Y \to \mathcal{O}_X$를 얻는다.
판별식과 디퍼런트 아이디얼의 절 \ref{discriminant-section-discriminant}를 보라.
\item 다음이 성립한다.
$\Theta_{Y/X}(\omega \wedge \eta) = \omega \wedge \Theta_{Y/X}(\eta)$
여기서 $\omega$는 $\Omega^\bullet_{X/\mathbf{Z}}$에 속하고, $\eta$는
$f_*\Omega^\bullet_{Y/\mathbf{Z}}$에 속한다.
\item $f$가 기저 스킴 $S$ 위의 사상이면, $\Theta_{Y/X}$는 복합체들의
사상 $f_*\Omega^\bullet_{Y/S} \to \Omega^\bullet_{X/S}$을 유도한다.
\end{enumerate}
\end{proposition}

\begin{proof}
판별식과 디퍼런트 아이디얼의 보조정리
\ref{discriminant-lemma-locally-comes-from-universal-finite}에 의해, 모든
$x \in X$에 대해 $d \geq 1$과 가환 도식
$$
\xymatrix{
Y \ar[d] &
V \ar[d] \ar[l] \ar[r] &
V_d \ar[d] \ar[r] &
Y_d = \Spec(B_d) \ar[d] \\
X &
U \ar[l] \ar[r] &
U_d \ar[r] &
X_d = \Spec(A_d)
}
$$
을 찾을 수 있다. 여기서 $x \in U \subset X$는 아핀 열린집합이고,
$V = f^{-1}(U)$이며 $V = U \times_{U_d} V_d$이다.
$U = \Spec(A)$ 및 $V = \Spec(B)$라고 쓰자.
$B = A \otimes_{A_d} B_d$임에 유의하고,
$B_d = A_d e_1 \oplus \ldots \oplus A_d e_d$임을 상기하자.
$a_1, \ldots, a_r \in A$와 $b_1, \ldots, b_s \in B$가 주어졌다고 하자.
$b_j = \sum a_{j, l} e_l$로 쓸 수 있으며, 여기서
$a_{j, l} \in A$이다. $N = r + sd$로 놓고 다음 분해들을 생각하자.
$$
\xymatrix{
V \ar[r] \ar[d] &
V' = \mathbf{A}^N \times V_d \ar[r] \ar[d] &
V_d \ar[d] \\
U \ar[r]&
U' = \mathbf{A}^N \times U_d \ar[r] &
U_d
}
$$
여기서 오른쪽 아래 수평 화살표는 앞 도식의 사상 $U \to U_d$와 사상
$U \to \mathbf{A}^N$에 의해 정해진다. 뒤의 사상은
$a_1, \ldots, a_r, a_{1, 1}, \ldots, a_{s, d}$로 주어진다. 그러면 함수들
$a_1, \ldots, a_r$는
$\Gamma(U', \mathcal{O}_{U'}) \to \Gamma(U, \mathcal{O}_U)$의 상에
속하고, 함수들 $b_1, \ldots, b_s$는
$\Gamma(V', \mathcal{O}_{V'}) \to \Gamma(V, \mathcal{O}_V)$의 상에
속함을 알 수 있다. 따라서 임의의 유한한 원소 모음\footnote{결국 이
원소들은 모두 다음 꼴 원소들의 유한합이다.
$a_0 \text{d}a_1 \wedge \ldots \wedge \text{d}a_i$이며
$a_0, \ldots, a_i \in A$인 경우와, 다음 꼴 원소들의 유한합인 경우가 있다.
$b_0 \text{d}b_1 \wedge \ldots \wedge \text{d}b_j$이며
$b_0, \ldots, b_j \in B$이다.}을 다음 군들에서 택했을 때
$$
\Gamma(V, \Omega^i_{Y/\mathbf{Z}}),\quad i = 0, 1, 2, \ldots
\quad\text{및}\quad
\Gamma(U, \Omega^j_{X/\mathbf{Z}}),\quad j = 0, 1, 2, \ldots
$$
위와 같은 분해들 $V \to V' \to V_d$ 및 $U \to U' \to U_d$를 찾을 수
있다. 여기서 $V' = \mathbf{A}^N \times V_d$이고
$U' = \mathbf{A}^N \times U_d$이며, 이 단면들은 다음 군에 속하는
단면들을 당겨서 얻어진다.
$$
\Gamma(V', \Omega^i_{V'/\mathbf{Z}}),\quad i = 0, 1, 2, \ldots
\quad\text{및}\quad
\Gamma(U', \Omega^j_{U'/\mathbf{Z}}),\quad j = 0, 1, 2, \ldots
$$
따라서 $\text{d} \circ \Theta_{Y/X} = \Theta_{Y/X} \circ \text{d}$를
확인하려면 $\Theta_{V'/U'}$에 대해 이를 확인하면 충분하다. 보조정리의
성질 (2)에 대해서도 마찬가지다.

\medskip\noindent
판별식과 디퍼런트 아이디얼의 보조정리
\ref{discriminant-lemma-universal-finite-syntomic-smooth}와
\ref{discriminant-lemma-universal-finite-syntomic-etale}에 의해 스킴
$U_d$는 매끄럽고, 사상 $V_d \to U_d$는 $U_d$의 조밀한 열린집합 위에서
에탈이다. 따라서 사상 $V' \to U'$도 같은 성질을 갖는다.
$\Omega_{U'/\mathbf{Z}}$가 국소 자유이므로
$\Omega^p_{U'/\mathbf{Z}}$는 꼬임이 없다. 따라서 $V'$가 유한 에탈인
열린집합에 제한한 뒤 원하는 관계들을 확인하면 충분하다. 그 뒤 전사
에탈 기저변환을 하고 관계들을 확인해도 된다. 따라서 유한 에탈 덮개를
분할하여 다음 꼴 사상을 생각한다고 가정할 수 있다.
$$
\coprod\nolimits_{i = 1, \ldots, d} W \longrightarrow W
$$
여기서 $W$는 $\mathbf{Z}$ 위에서 매끄럽다. 이 경우 우리 구성의 국소적
성질들은 참이기만 하면 자명하게 확인된다. 이로써 (1), (2)의 증명이
끝난다.

\medskip\noindent
마지막으로, $\Omega_{Y/S}$는 $\Omega_{Y/\mathbf{Z}}$를
$\Omega_{S/\mathbf{Z}}$의 국소 단면들을 당겨 얻은 것들이 생성하는
부분가군으로 나눈 몫이므로 (3)은 (2)에서 따른다.
\end{proof}

\begin{example}
\label{derham-example-Garel}
$A$를 환이라 하자.
$f = x^d + \sum_{0 \leq i < d} a_{d - i} x^i \in A[x]$라 하자.
$B = A[x]/(f)$라 하자. 명제 \ref{derham-proposition-Garel}에 의해 복합체들의
사상
$$
\Theta_{B/A} : \Omega^\bullet_B \longrightarrow \Omega^\bullet_A
$$
을 얻는다. 특히 $t \in B$가 $x \in A[x]$의 상을 나타낸다면 원소들
$$
\Theta_{B/A}(t^i\text{d}t) \in \Omega^1_A,\quad i = 0, \ldots, d - 1
$$
을 생각할 수 있다. 이 원소들은 무엇인가? 명제
\ref{derham-proposition-Garel}의 증명에서 사용한 것과 같은 원리에 의해 대역적인
경우, 즉 $A = \mathbf{Z}[a_1, \ldots, a_d]$인 경우에 계산하면 충분하다.
나아가 $A$를 분수체 $\mathbf{Q}(a_1, \ldots, a_d)$로 바꾸어도 된다.
기호적으로
$$
f = \prod\nolimits_{i = 1, \ldots, d} (x - \alpha_i)
$$
라고 쓰면 $\mathbf{Q}(\alpha_1, \ldots, \alpha_d)$ 위에서 대수 $B$는
분할된다.
$$
\mathbf{Q}(a_0, \ldots, a_{d - 1})[x]/(f) 
\longrightarrow
\prod\nolimits_{i = 1, \ldots, d} \mathbf{Q}(\alpha_1, \ldots, \alpha_d),
\quad
t \longmapsto (\alpha_1, \ldots, \alpha_d)
$$
따라서 예를 들어
$$
\Theta(\text{d}t) = \sum \text{d} \alpha_i = - \text{d}a_1
$$
이다. 다음으로
$$
\Theta(t\text{d}t) = \sum \alpha_i \text{d}\alpha_i =
a_1 \text{d} a_1 - \text{d}a_2
$$
이다. 다음으로
$$
\Theta(t^2\text{d}t) = \sum \alpha_i^2 \text{d}\alpha_i =
- a_1^2 \text{d} a_1 + a_1 \text{d}a_2 + a_2 \text{d}a_1 - \text{d}a_3
$$
이다(계산 오류 가능성은 제외한다). 이와 같이 계속된다. 이로부터
$f(x) = x^d - a$이면
$$
\Theta_{B/A}(t^i\text{d}t) =
\left\{
\begin{matrix}
0 & \text{이면} & i = 0, \ldots, d - 2 \\
\text{d}a & \text{이면} & i = d - 1
\end{matrix}
\right.
$$
가 $\Omega_A$ 안에서 성립할 것으로 보인다. 이는 참이다. 이 특별한
경우에는 확대 $\mathbf{Q}(a)[x]/(x^d - a)$에 대해 계산하여 직접 확인할
수 있기 때문이다.
\end{example}

\begin{lemma}
\label{derham-lemma-Garel-map-frobenius-smooth-char-p}
$p$를 소수라 하자. $X \to S$를 상대 차원 $d$인 매끄러운 사상이라
하고, 이 스킴들의 표수가 $p$라고 하자. 상대 프로베니우스
$F_{X/S} : X \to X^{(p)}$를 생각하자. 이는 $X/S$의 상대 프로베니우스이다
(대수다양체의 정의
\ref{varieties-definition-relative-frobenius})는 유한 신토믹이고, 이에
대응하는 사상
$$
\Theta_{X/X^{(p)}} :
F_{X/S, *}\Omega^\bullet_{X/S} \to \Omega^\bullet_{X^{(p)}/S}
$$
은 차수가 $d$일 때 전사를 정의하고, 그 외 모든 차수에서 영이다.
\end{lemma}

\begin{proof}
대수다양체의 보조정리 \ref{varieties-lemma-relative-frobenius-finite}에 의해
$F_{X/S}$는 유한 사상이다. $F_{X/S}$가 평탄임을 증명하려면, 올들 사이의
사상 $F_{X/S, s} : X_s \to X^{(p)}_s$가 모든 $s \in S$에 대해
평탄임을 보이면 충분하다. 사상 심화의 정리 \ref{more-morphisms-theorem-criterion-flatness-fibre}를
보라. $X_s \to X^{(p)}_s$의 평탄성은 가환대수의 보조정리
\ref{algebra-lemma-CM-over-regular-flat}와 이미 증명한 유한성에서 따른다.
사상 심화의 보조정리
\ref{more-morphisms-lemma-lci-permanence}에 의해 사상 $F_{X/S}$는 완전
국소 교차 사상이다. 따라서 $F_{X/S}$는 유한 신토믹이다(사상 심화의 보조정리 \ref{more-morphisms-lemma-flat-lci}를 보라).

\medskip\noindent
모든 점 $x \in X$에 대해 가환 도식
$$
\xymatrix{
X \ar[d] & U \ar[l] \ar[d]_\pi \\
S & \mathbf{A}^d_S \ar[l]
}
$$
을 택할 수 있다. 여기서 $\pi$는 에탈이고 $x \in U$는 $X$의
열린집합이다. 스킴의 사상의 보조정리
\ref{morphisms-lemma-smooth-etale-over-affine-space}를 보라. 사상
$\mathbf{A}^d_S \to \mathbf{A}^d_S$,
$(x_1, \ldots, x_d) \mapsto (x_1^p, \ldots, x_d^p)$은
$\mathcal{A}^d_S$의 $S$ 위 상대 프로베니우스임에 유의하자. 가환 도식
$$
\xymatrix{
U \ar[d]_\pi \ar[r]_{F_{X/S}} & U^{(p)} \ar[d]^{\pi^{(p)}} \\
\mathbf{A}^d_S \ar[r]^{x_i \mapsto x_i^p} & \mathbf{A}^d_S
}
$$
은 $\pi : U \to \mathbf{A}^d_S$에 대한 대수다양체의 보조정리
\ref{varieties-lemma-relative-frobenius-endomorphism-identity}의 도식이며,
스킴의 에탈 사상의 보조정리
\ref{etale-lemma-relative-frobenius-etale}에 의해 카르테시안이다.
$\Theta$의 구성이 기저변환과 양립하고
$\Omega_{U/S} = \pi^*\Omega_{\mathbf{A}^d_S/S}$이므로(보조정리
\ref{derham-lemma-etale}), $\mathbf{A}^d_S$에 대해 보조정리를 보이면 충분하다.

\medskip\noindent
$A$를 표수 $p$인 환이라 하자. 유일한 $A$-대수 준동형
$A[y_1, \ldots, y_d] \to A[x_1, \ldots, x_d]$을 생각하자. 이는
$y_i$를 $x_i^p$로 보낸다. 위의 논증에
따라 사상
$$
\Theta^i : \Omega^i_{A[x_1, \ldots, x_d]/A} \to
\Omega^i_{A[y_1, \ldots, y_d]/A}
$$
을 계산하면 된다. 독자가 이 경우의 계산을 직접 해 보기를 권한다.
예 \ref{derham-example-Garel}에서와 같이 대역적인 경우의 $\Theta^i$에 대한
공식을 계산하는 문제로 환원할 수 있다.
$$
\mathbf{Z}[y_1, \ldots, y_d] \to \mathbf{Z}[x_1, \ldots, x_d],\quad
y_i \mapsto x_i^p
$$
다시 $\Theta^i$의 공식은 복소수의 경우, 즉 $\mathbf{C}$-대수 사상
$$
\mathbf{C}[y_1, \ldots, y_d] \to \mathbf{C}[x_1, \ldots, x_d],\quad
y_i \mapsto x_i^p
$$
에 대해 계산하면 찾을 수 있다. 나아가
$x_1, \ldots, x_d$와 $y_1, \ldots, y_d$를 가역화해도 된다. 이 경우
$\text{d}x_i = p^{-1} x_i^{- p + 1}\text{d}y_i$이다. 따라서
\begin{align*}
\Theta^i(
x_1^{e_1} \ldots x_d^{e_d} \text{d}x_1 \wedge \ldots \wedge \text{d}x_i)
& =
p^{-i} \Theta^i(
x_1^{e_1 - p + 1} \ldots x_i^{e_i - p + 1} x_{i + 1}^{e_{i + 1}} \ldots
x_d^{e_d} \text{d}y_1 \wedge \ldots \wedge \text{d}y_i ) \\
& =
p^{-i} \text{Trace}(x_1^{e_1 - p + 1} \ldots x_i^{e_i - p + 1}
x_{i + 1}^{e_{i + 1}} \ldots x_d^{e_d})
\text{d}y_1 \wedge \ldots \wedge \text{d}y_i
\end{align*}
이다. 이는 $\Theta^i$의 성질들에서 따른다. 초등적인 계산에 의해 위 식의
대각합은 $e_1, \ldots, e_i$가 $-1$과 법 $p$로 합동이고
$e_{i + 1}, \ldots, e_d$가 $p$로 나누어떨어지는 경우를 제외하면
영이다. 더욱이 그 경우에는
$$
p^{d - i} y_1^{(e_1 - p + 1)/p} \ldots y_i^{(e_i - p + 1)/p}
y_{i + 1}^{e_{i + 1}/p} \ldots y_d^{e_d/p}
\text{d}y_1 \wedge \ldots \wedge \text{d}y_i
$$
을 얻는다. 따라서 표수 $p$에서는 $d = i$인 경우를 제외하면 영이고,
그 경우에는 가능한 모든 $d$-형식을 얻는다.
\end{proof}








\section{푸앵카레 쌍대성}
\label{derham-section-poincare-duality}

\noindent
이 절에서는 체 위 고유 매끄러운 스킴의 드람 코호몰로지에 대한 푸앵카레
쌍대성을 증명한다. 먼저 호지 코호몰로지에서 이것이 어떻게 성립하는지
설명하자.

\begin{lemma}
\label{derham-lemma-duality-hodge}
$k$를 체라 하자. $X$를 $k$ 위의 공집합이 아닌 매끄러운 고유 스킴이라
하고 차원이 $d$인 등차원 스킴이라고 하자. 그러면 $k$-선형 사상
$$
t : H^d(X, \Omega^d_{X/k}) \longrightarrow k
$$
이 존재한다. 이는 $H^0(X, \mathcal{O}_X)$의 가역원을 곱하는 사상을
앞에서 합성하는 것까지 유일하며, 모든 $p, q$에 대해 쌍
$$
H^q(X, \Omega^p_{X/k}) \times H^{d - q}(X, \Omega^{d - p}_{X/k})
\longrightarrow
k, \quad
(\xi, \xi') \longmapsto t(\xi \cup \xi')
$$
은 완전하다.
\end{lemma}

\begin{proof}
스킴의 쌍대성의 보조정리
\ref{duality-lemma-duality-proper-over-field}에 의해
$\omega_X^\bullet = \Omega^d_{X/k}[d]$이다. $\Omega_{X/k}$는 계수
$d$의 국소 자유 가군이므로(스킴의 사상의 보조정리
\ref{morphisms-lemma-smooth-omega-finite-locally-free}),
$$
\Omega^d_{X/k} \otimes_{\mathcal{O}_X} (\Omega^p_{X/k})^\vee
\cong
\Omega^{d - p}_{X/k}
$$
이다. 따라서 스킴의 쌍대성의 보조정리
\ref{duality-lemma-duality-proper-over-field-perfect}에 의해 명제가
성립하도록 하는 $k$-선형 사상 $t : H^d(X, \Omega^d_{X/k}) \to k$를
얻는다. 특히 쌍
$H^0(X, \mathcal{O}_X) \times H^d(X, \Omega^d_{X/k}) \to k$
은 완전하다. 따라서 임의의 $k$-선형 사상
$t' : H^d(X, \Omega^d_{X/k}) \to k$는 $\xi \mapsto t(g\xi)$ 꼴이며,
여기서 $g \in H^0(X, \mathcal{O}_X)$이다.
물론 $t'$가 $H^0(X, \mathcal{O}_X)$와 $H^d(X, \Omega^d_{X/k})$ 사이의
쌍대성을 여전히 주려면 $g$가 가역원이어야 한다.
$\langle -, - \rangle_{p, q}$를 $t$를 사용하여 구성한 쌍이라 하고,
$\langle -, - \rangle'_{p, q}$를 $t'$를 사용하여 구성한 쌍이라 하자.
분명히
$$
\langle \xi, \xi' \rangle'_{p, q} =
\langle g\xi, \xi' \rangle_{p, q}
$$
이다. 여기서 $\xi \in H^q(X, \Omega^p_{X/k})$이고
$\xi' \in H^{d - q}(X, \Omega^{d - p}_{X/k})$이다. $g$는 가역원이므로,
$t'$를 사용하고 $t$를 사용하지 않아도 모든 $p, q$에 대해 완전한 쌍을
얻는다.
\end{proof}

\begin{lemma}
\label{derham-lemma-bottom-part-degenerates}
$k$를 체라 하자. $X$를 $k$ 위 매끄러운 고유 스킴이라 하자. 사상
$$
\text{d} : H^0(X, \mathcal{O}_X) \to H^0(X, \Omega^1_{X/k})
$$
은 영이다.
\end{lemma}

\begin{proof}
$X$는 $k$ 위에서 매끄러우므로 $k$ 위에서 기하적으로 축약이다.
대수다양체의 보조정리 \ref{varieties-lemma-smooth-geometrically-normal}를
보라. 따라서 $H^0(X, \mathcal{O}_X) = \prod k_i$는 유한 분리 체 확대
$k_i/k$들의 유한 곱이다. 대수다양체의 보조정리
\ref{varieties-lemma-proper-geometrically-reduced-global-sections}를 보라.
그러므로 $\Omega_{H^0(X, \mathcal{O}_X)/k} = \prod \Omega_{k_i/k} = 0$이다
(예를 들어 가환대수의 보조정리
\ref{algebra-lemma-characterize-separable-algebraic-field-extensions}를 보라).
보조정리의 사상은
$$
H^0(X, \mathcal{O}_X) \to
\Omega_{H^0(X, \mathcal{O}_X)/k} \to
H^0(X, \Omega_{X/k})
$$
으로 분해된다. 이는 드람 복합체의 함자성에서 따른다(절
\ref{derham-section-de-rham-complex}를 보라). 따라서 결론을 얻는다.
\end{proof}

\begin{lemma}
\label{derham-lemma-top-part-degenerates}
$k$를 체라 하자. $X$를 $k$ 위 매끄러운 고유 스킴이라 하고, 차원이
$d$인 등차원 스킴이라고 하자. 사상
$$
\text{d} : H^d(X, \Omega^{d - 1}_{X/k}) \to H^d(X, \Omega^d_{X/k})
$$
은 영이다.
\end{lemma}

\begin{proof}
보조정리 \ref{derham-lemma-bottom-part-degenerates}와
\ref{derham-lemma-duality-hodge}를 결합하면 이 결과가 따른다고 생각하기 쉽다.
그러나 이는 통하지 않는다. 사상들
$\mathcal{O}_X \to \Omega^1_{X/k}$와
$\Omega^{d - 1}_{X/k} \to \Omega^d_{X/k}$는 $\mathcal{O}_X$-선형이
아니기 때문이다. 따라서 스킴의 쌍대성의 주의
\ref{duality-remark-coherent-duality-proper-over-field}에서 논한 함자성을
사용하여 보조정리 \ref{derham-lemma-bottom-part-degenerates}의 사상이 이
보조정리의 사상과 쌍대라고 결론내릴 수 없다.

\medskip\noindent
$X$를 $X$의 한 연결 성분으로 바꿀 수 있다. 따라서 $X$가 기약이라고
가정할 수 있다. 대수다양체의 보조정리
\ref{varieties-lemma-smooth-geometrically-normal}와
\ref{varieties-lemma-proper-geometrically-reduced-global-sections}에 의해
$k' = H^0(X, \mathcal{O}_X)$는 유한 분리 확대 $k'/k$이다.
$\Omega_{k'/k} = 0$이므로(예를 들어 가환대수의 보조정리
\ref{algebra-lemma-characterize-separable-algebraic-field-extensions}를 보라),
$\Omega_{X/k} = \Omega_{X/k'}$이다(스킴의 사상의 보조정리
\ref{morphisms-lemma-triangle-differentials}를 보라). 따라서 $k$를 $k'$로
바꾸어 $H^0(X, \mathcal{O}_X) = k$라고 가정할 수 있다.

\medskip\noindent
$H^0(X, \mathcal{O}_X) = k$라고 가정하자. 보조정리
\ref{derham-lemma-duality-hodge}에 의해
$\dim H^d(X, \Omega^d_{X/k}) = 1$임을 얻는다. 먼저 $k$의 표수가 소수
$p$라고 가정하자. $F_{X/k} : X \to X^{(p)}$를 $X$의 $k$ 위 상대
프로베니우스라 하자. 보조정리
\ref{derham-lemma-Garel-map-frobenius-smooth-char-p}에서 이 사상에 대해 증명한
사실들을 염두에 두기 바란다. 가환 도식
$$
\xymatrix{
H^d(X, \Omega^{d - 1}_{X/k}) \ar[d] \ar[r] &
H^d(X^{(p)}, F_{X/k, *}\Omega^{d - 1}_{X/k}) \ar[d] \ar[r]_{\Theta^{d - 1}} &
H^d(X^{(p)}, \Omega^{d - 1}_{X^{(p)}/k}) \ar[d] \\
H^d(X, \Omega^d_{X/k}) \ar[r] &
H^d(X^{(p)}, F_{X/k, *}\Omega^d_{X/k}) \ar[r]^{\Theta^d} &
H^d(X^{(p)}, \Omega^d_{X^{(p)}/k})
}
$$
을 생각하자. $F_{X/k}$가 유한이므로 왼쪽의 두 수평 화살표는 동형이다.
스킴의 코호몰로지의 보조정리
\ref{coherent-lemma-relative-affine-cohomology}를 보라. 오른쪽 사각형은
$\Theta_{X^{(p)}/X}$가 복합체들의 사상이고 $\Theta^{d - 1}$가
영이므로 가환한다. 따라서 $\Theta^d$가 영이 아님을 보이면 충분하다.
실제로 위의 논의에 의해 사상 $\Theta^d$의 정의역 차원은 $1$이다. 한편
$$
\Theta^d : F_{X/k, *}\Omega^d_{X/k} \to \Omega^d_{X^{(p)}/k}
$$
은 전사임을 알고 있다. 따라서 오른쪽 완전 함자 $H^d(X^{(p)}, -)$를
적용한 뒤에도 전사이다. 이 함자가 오른쪽 완전인 것은 $d$보다 높은
차수에서 코호몰로지가 소멸하기 때문이며, 이는 층 코호몰로지의 명제
\ref{cohomology-proposition-vanishing-Noetherian}에서 따른다. 마지막으로
$H^d(X^{(d)}, \Omega^d_{X^{(d)}/k})$는 영이 아니다. 예를 들어 이는
$H^0(X^{(d)}, \mathcal{O}_{X^{(p)}})$와 쌍대이다. 여기서는 보조정리
\ref{derham-lemma-duality-hodge}를 $X^{(p)}$에 $k$ 위에서 적용하였다. 이로써
이 경우의 증명이 끝난다.

\medskip\noindent
마지막으로 $k$의 표수가 $0$이라고 가정하자. $k$를 유한형
$\mathbf{Z}$-부분대수들 $R$의 여과 쌍대극한으로 쓸 수 있다. 이들 중
하나에 대해 스킴들의 카르테시안 도식
$$
\xymatrix{
X \ar[d] \ar[r] & Y \ar[d] \\
\Spec(k) \ar[r] & \Spec(R)
}
$$
을 찾을 수 있으며, $Y \to \Spec(R)$은 상대 차원 $d$인 매끄러운 고유
사상이다. 스킴의 극한의 보조정리들
\ref{limits-lemma-descend-finite-presentation},
\ref{limits-lemma-descend-smooth}, \ref{limits-lemma-descend-dimension-d},
\ref{limits-lemma-eventually-proper}를 보라. 가군들
$M^{i, j} = H^j(Y, \Omega^i_{Y/R})$은 유한 $R$-가군이다. 스킴의 코호몰로지의 보조정리
\ref{coherent-lemma-proper-over-affine-cohomology-finite}를 보라. 따라서
$R$을 한 국소화로 바꾼 뒤에는 이 가군들이 모두 유한 자유라고 가정할
수 있다. 평탄 기저변환에 의해
$M^{i, j} \otimes_R k = H^j(X, \Omega^i_{X/k})$이다(스킴의 코호몰로지의 보조정리 \ref{coherent-lemma-flat-base-change-cohomology}를 보라).
따라서 $M^{d - 1, d} \to M^{d, d}$가 영임을 보이면 충분하다. 이는
정역 위 유한 자유 가군들 사이의 사상이므로, 소아이디얼들
$\mathfrak p \subset R$의 조밀한 집합을 찾아 $\kappa(\mathfrak p)$와
텐서한 뒤 영이 되도록 하면 충분하다. $R$은 $\mathbf{Z}$ 위 유한형이므로
잉여체의 표수가 양수인 소아이디얼들 $\mathfrak p$의 모음을 택할 수 있다
(세부 사항은 생략한다). 다음에 유의하자.
$$
M^{d - 1, d} \otimes_R \kappa(\mathfrak p) =
H^d(Y_{\kappa(\mathfrak p)},
\Omega^{d - 1}_{Y_{\kappa(\mathfrak p)}/\kappa(\mathfrak p)})
$$
예를 들어 스킴의 극한의 보조정리
\ref{limits-lemma-higher-direct-images-zero-above-dimension-fibre}에 의해
이 식이 성립한다. $M^{d, d}$에 대해서도 마찬가지다. 따라서
$M^{d - 1, d} \otimes_R \kappa(\mathfrak p) \to
M^{d, d} \otimes_R \kappa(\mathfrak p)$
는 위에서 다룬 양의 표수의 경우에 의해 영이다.
\end{proof}

\begin{proposition}
\label{derham-proposition-poincare-duality}
$k$를 체라 하자. $X$를 $k$ 위의 공집합이 아닌 매끄러운 고유 스킴이라
하고 차원이 $d$인 등차원 스킴이라고 하자. 그러면 $k$-선형 사상
$$
t : H^{2d}_{dR}(X/k) \longrightarrow k
$$
이 존재한다. 이는 $H^0(X, \mathcal{O}_X)$의 가역원을 곱하는 사상을
앞에서 합성하는 것까지 유일하며, 모든 $i$에 대해 쌍
$$
H^i_{dR}(X/k) \times H_{dR}^{2d - i}(X/k)
\longrightarrow
k, \quad
(\xi, \xi') \longmapsto t(\xi \cup \xi')
$$
은 완전하다.
\end{proposition}

\begin{proof}
호지-드람 스펙트럼 열(절 \ref{derham-section-hodge-to-de-rham}),
$\Omega^i_{X/k}$가 $i > d$일 때 소멸한다는 사실, 층 코호몰로지의 명제
\ref{cohomology-proposition-vanishing-Noetherian}의 소멸, 그리고 보조정리
\ref{derham-lemma-bottom-part-degenerates}와
\ref{derham-lemma-top-part-degenerates}의 결과에 의해
$H^0_{dR}(X/k) = H^0(X, \mathcal{O}_X)$ 및
$H^d(X, \Omega^d_{X/k}) = H_{dR}^{2d}(X/k)$임을 알 수 있다. 더
정확히 말하면 이 식별들은 복합체들의 사상
$$
\Omega^\bullet_{X/k} \to \mathcal{O}_X[0]
\quad\text{및}\quad
\Omega^d_{X/k}[-d] \to \Omega^\bullet_{X/k}
$$
으로부터 나온다. $t : H_{dR}^{2d}(X/k) \to k$를 택하자. 이는 이
식별을 통해 보조정리 \ref{derham-lemma-duality-hodge}의 $t$에 대응한다. 그러면 어느
경우든 명제에 표시된 쌍은 $i = 0$일 때 완전하다.

\medskip\noindent
$\underline{k}$를 값이 $k$인 상수층이라 하자. 이는 $X$ 위의 층이다. 간단히
$\Omega^\bullet = \Omega^\bullet_{X/k}$라고 쓰자. 우리 상황에서
(\ref{derham-equation-wedge})의 사상은
$$
\wedge :
\text{Tot}(\Omega^\bullet \otimes_{\underline{k}} \Omega^\bullet)
\longrightarrow
\Omega^\bullet
$$
과 같다. 모든 정수 $p = 0, 1, \ldots, d$에 대해 이 사상은 차수상의
이유로 부분복합체
$\text{Tot}(\sigma_{> p} \Omega^\bullet \otimes_{\underline{k}}
\sigma_{\geq d - p} \Omega^\bullet)$를 소멸시킨다. 따라서 $\wedge$를
부분복합체
$\text{Tot}(\Omega^\bullet \otimes_{\underline{k}}
\sigma_{\geq d - p}\Omega^\bullet)$에 제한한 사상은 복합체들의 사상
$$
\gamma_p :
\text{Tot}(\sigma_{\leq p} \Omega^\bullet \otimes_{\underline{k}}
\sigma_{\geq d - p} \Omega^\bullet)
\longrightarrow
\Omega^\bullet
$$
을 통해 분해된다. 절 \ref{derham-section-cup-product}에서와 같은 절차를 사용하면
컵곱들
$$
H^i(X, \sigma_{\leq p} \Omega^\bullet) \times
H^{2d - i}(X, \sigma_{\geq d - p}\Omega^\bullet)
\longrightarrow
H_{dR}^{2d}(X, \Omega^\bullet)
$$
을 얻는다. $p$에 대한 귀납법으로, 이 컵곱들을 $t$와 합성하면
$H^i(X, \sigma_{\leq p} \Omega^\bullet)$와
$H^{2d - i}(X, \sigma_{\geq d - p}\Omega^\bullet)$ 사이에 완전한 쌍이
유도됨을 증명하겠다. $p = d$인 경우가 명제의
주장이다.

\medskip\noindent
기초 단계는 $p = 0$이다. 이 경우 보조정리
\ref{derham-lemma-duality-hodge}의 $H^i(X, \mathcal{O}_X)$와
$H^{d - i}(X, \Omega^d)$ 사이의 쌍을 그대로 얻으므로 결과가 성립한다.

\medskip\noindent
귀납 단계. 결과가 $p$에 대해 성립한다고 하자. 그러면 완전 삼각형
$$
\Omega^{p + 1}[-p - 1] \to
\sigma_{\leq p + 1}\Omega^\bullet \to
\sigma_{\leq p}\Omega^\bullet \to
\Omega^{p + 1}[-p]
$$
과 완전 삼각형
$$
\sigma_{\geq d - p}\Omega^\bullet \to
\sigma_{\geq d - p - 1}\Omega^\bullet \to
\Omega^{d - p - 1}[-d + p + 1] \to
(\sigma_{\geq d - p}\Omega^\bullet)[1]
$$
을 생각하자. 둘 다 $\underline{k}$-가군 층들의 복합체의 호모토피 범주
안에서 완전 삼각형임에 유의하자. 특히 사상들
$\sigma_{\leq p}\Omega^\bullet \to \Omega^{p + 1}[-p]$ 및
$\Omega^{d - p - 1}[-d + p + 1] \to (\sigma_{\geq d - p}\Omega^\bullet)[1]$
은 실제 복합체 사상들로 주어지며, 구체적으로 미분
$\Omega^p \to \Omega^{p + 1}$ 및 미분
$\Omega^{d - p - 1} \to \Omega^{d - p}$를 사용한다. 이 완전 삼각형들에
딸린 긴 완전 코호몰로지 열들을 생각하자.
$$
\xymatrix{
H^{i - 1}(X, \sigma_{\leq p}\Omega^\bullet) \ar[d]_a \\
H^i(X, \Omega^{p + 1}[-p - 1]) \ar[d]_b \\
H^i(X, \sigma_{\leq p + 1}\Omega^\bullet) \ar[d]_c \\
H^i(X, \sigma_{\leq p}\Omega^\bullet) \ar[d]_d \\
H^{i + 1}(X, \Omega^{p + 1}[-p - 1])
}
\quad\quad
\xymatrix{
H^{2d - i  + 1}(X, \sigma_{\geq d - p}\Omega^\bullet) \\
H^{2d - i}(X, \Omega^{d - p - 1}[-d + p + 1]) \ar[u]_{a'} \\
H^{2d - i}(X, \sigma_{\geq d - p - 1}\Omega^\bullet) \ar[u]_{b'} \\
H^{2d - i}(X, \sigma_{\geq d - p}\Omega^\bullet) \ar[u]_{c'} \\
H^{2d - i - 1}(X, \Omega^{d - p - 1}[-d + p + 1]) \ar[u]_{d'}
}
$$
귀납 가정과 보조정리 \ref{derham-lemma-duality-hodge}에 의해, 위에서 구성한
첫째, 둘째, 넷째, 다섯째 행의 $k$-벡터공간들 사이의 쌍들은 완전하다.
$5$-보조정리에 의해 가운데 행의 코호몰로지 군들 사이의 쌍이 완전함을
보이려면 쌍들 $(a, a')$, $(b, b')$, $(c, c')$, $(d, d')$가 주어진
쌍들과 양립함을 보이면 충분하다(아래를 보라).

\medskip\noindent
먼저 쌍 $(c, c')$에 대해 이를 증명하자. 다음 가환 도식이 있음을
관찰하면 된다.
$$
\xymatrix{
\text{Tot}(\sigma_{\leq p} \Omega^\bullet \otimes_{\underline{k}}
\sigma_{\geq d - p} \Omega^\bullet) \ar[d]_{\gamma_p} &
\text{Tot}(\sigma_{\leq p + 1} \Omega^\bullet \otimes_{\underline{k}}
\sigma_{\geq d - p} \Omega^\bullet) \ar[l] \ar[d] \\
\Omega^\bullet &
\text{Tot}(\sigma_{\leq p + 1} \Omega^\bullet \otimes_{\underline{k}}
\sigma_{\geq d - p - 1} \Omega^\bullet) \ar[l]_-{\gamma_{p + 1}}
}
$$
따라서 $\alpha \in H^i(X, \sigma_{\leq p + 1}\Omega^\bullet)$ 및
$\beta \in H^{2d - i}(X, \sigma_{\geq d - p}\Omega^\bullet)$가 주어지면,
컵곱의 함자성에 의해
$\gamma_p(\alpha \cup c'(\beta)) = \gamma_{p + 1}(c(\alpha) \cup \beta)$
를 얻는다.

\medskip\noindent
마찬가지로 쌍 $(b, b')$에 대해서는 가환 도식
$$
\xymatrix{
\text{Tot}(\sigma_{\leq p + 1} \Omega^\bullet \otimes_{\underline{k}}
\sigma_{\geq d - p - 1} \Omega^\bullet) \ar[d]_{\gamma_{p + 1}} &
\text{Tot}(\Omega^{p + 1}[-p - 1] \otimes_{\underline{k}}
\sigma_{\geq d - p - 1} \Omega^\bullet) \ar[l] \ar[d] \\
\Omega^\bullet &
\Omega^{p + 1}[-p - 1]
\otimes_{\underline{k}}
\Omega^{d - p - 1}[-d + p + 1] \ar[l]_-\wedge
}
$$
을 사용하고 같은 방식으로 논증한다.

\medskip\noindent
쌍 $(d, d')$에 대해서는 가환 도식
$$
\xymatrix{
\Omega^{p + 1}[-p] \otimes_{\underline{k}}
\Omega^{d - p - 1}[-d + p] \ar[d] &
\text{Tot}(\sigma_{\leq p}\Omega^\bullet \otimes_{\underline{k}}
\Omega^{d - p - 1}[-d + p]) \ar[l] \ar[d] \\
\Omega^\bullet &
\text{Tot}(\sigma_{\leq p}\Omega^\bullet \otimes_{\underline{k}}
\sigma_{\geq d - p}\Omega^\bullet) \ar[l]
}
$$
을 사용하고 코호몰로지 류들
$H^i(X, \sigma_{\leq p}\Omega^\bullet)$ 및
$H^{2d - i}(X, \Omega^{d - p - 1}[-d + p])$을 살펴본다.
$i$를 $i - 1$로 바꾸면 쌍 $(a, a')$에 대한 결과를 얻는다. 이로써
우리의 쌍들이 완전하다는 증명이 끝난다.

\medskip\noindent
$t$가 $H^0(X, \mathcal{O}_X)$의 가역원을 곱하는 사상을 앞에서
합성하는 것까지 유일함을 보이는 논증은 생략한다.
\end{proof}







\section{천 특성류}
\label{derham-section-chern-classes}

\noindent
지금까지 증명한 결과들은 베유 코호몰로지 이론의 절
\ref{weil-section-chern}의 논의를 사용하여 드람 코호몰로지의 천
특성류를 만들기에 충분하다.

\begin{lemma}
\label{derham-lemma-chern-classes}
모든 준콤팩트 준분리 스킴 $X$에 전체 천 특성류
$$
c^{dR} :
K_0(\textit{Vect}(X))
\longrightarrow
\prod\nolimits_{i \geq 0} H^{2i}_{dR}(X/\mathbf{Z})
$$
를 대응시키는 유일한 규칙이 존재하며, 다음 성질들을 갖는다.
\begin{enumerate}
\item $c^{dR}(\alpha + \beta) = c^{dR}(\alpha) c^{dR}(\beta)$이다. 여기서
$\alpha, \beta \in K_0(\textit{Vect}(X))$이다.
\item $f : X \to X'$가 준콤팩트 준분리 스킴들의 사상이면
$c^{dR}(f^*\alpha) = f^*c^{dR}(\alpha)$이다.
\item $\mathcal{L} \in \Pic(X)$가 주어지면
$c^{dR}([\mathcal{L}]) = 1 + c_1^{dR}(\mathcal{L})$이다.
\end{enumerate}
\end{lemma}

\noindent
이 구성은 모든 스킴으로 쉽게 확장할 수 있지만, 그러려면 베유 코호몰로지 이론의 절 \ref{weil-section-chern}의 논의를 약간 보강해야 한다.

\begin{proof}
베유 코호몰로지 이론의 명제 \ref{weil-proposition-chern-class}를
적용하겠다.

\medskip\noindent
$\mathcal{C}$를 모든 준콤팩트 준분리 스킴의 범주라 하자. 이 범주는
분명히 베유 코호몰로지 이론의 절 \ref{weil-section-chern}의 조건
(1), (2), (3)(a), (b), (c)를 만족한다.

\medskip\noindent
등급 대수의 범주로 가는 우리의 반변 함자 $A$는 $\mathcal{C}$에서
$X$를
$A(X) = \bigoplus_{i \geq 0} H_{dR}^{2i}(X/\mathbf{Z})$
로 보내고, 여기에는 컵곱을 준다. 함자성은 절
\ref{derham-section-de-rham-cohomology}에서, 컵곱은 절
\ref{derham-section-cup-product}에서 논하였다. 가법 사상 $c_1^A$로는 절
\ref{derham-section-first-chern-class}에서 구성한 $c_1^{dR}$을 택한다.

\medskip\noindent
실제로 보조정리 \ref{derham-lemma-cup-product-graded-commutative}에 의해 가환
대수들을 얻으며, 이는 $A$에 대해 공리 (1)이 성립함을 보인다.

\medskip\noindent
공리 (2)를 $A$에 대해 확인하려면
$H^*_{dR}(X \coprod Y/\mathbf{Z}) =  H^*_{dR}(X/\mathbf{Z}) \times
H^*_{dR}(Y/\mathbf{Z})$.
임을 확인하면 충분하다. 이는 드람 코호몰로지가 자리스키 위상에서 미분
등급 대수 층의 코호몰로지를 취하여 구성된다는 사실에서 따른다.

\medskip\noindent
$A$에 대한 공리 (3)은 가역 가군의 제1 천 특성류를 취하는 연산이
당김과 양립한다는 명제이다. 이는 더 일반적인 보조정리
\ref{derham-lemma-pullback-c1}에서 따른다.

\medskip\noindent
$A$에 대한 공리 (4)는 명제
\ref{derham-proposition-projective-space-bundle-formula}에서 증명한 사영공간
다발 공식이다.

\medskip\noindent
공리 (5). $X$를 준콤팩트 준분리 스킴이라 하고,
$\mathcal{E} \to \mathcal{F}$를 유한 국소 자유 $\mathcal{O}_X$-가군
사이의 전사라 하자. 그 계수들은 각각 $r + 1$과 $r$이다. 이에 대응하는 포함 사상을
$i : P' = \mathbf{P}(\mathcal{F}) \to \mathbf{P}(\mathcal{E}) = P$라 하자.
이는 $X$ 위 매끄러운 사영 스킴들 사이의 사상이며 $P'$를 $P$ 위의 유효
카르티에 인자로 나타낸다. 따라서 보조정리 \ref{derham-lemma-check-log-smooth}에
의해 $P' \subset P$에 대한 $\mathbf{Z}$ 위 로그 극 복합체가 정의된다.
그러므로 $a \in A(P)$와 $i^*a = 0$에 대해 보조정리
\ref{derham-lemma-log-complex-consequence}에 의해
$a \cup c_1^A(\mathcal{O}_P(P')) = 0$이다. 이로써 증명이 끝난다.
\end{proof}

\begin{remark}
\label{derham-remark-splitting-principle}
베유 코호몰로지 이론의 보조정리
\ref{weil-lemma-splitting-principle}(분할 원리)와
\ref{weil-lemma-chern-classes-E-tensor-L}(텐서곱의 천 특성류)에 대응하는
명제들은 준콤팩트 준분리 스킴 위 드람 천 특성류에 대해서도 성립한다.
보조정리 \ref{derham-lemma-chern-classes}의 증명에서 베유 코호몰로지 이론의
절 \ref{weil-section-chern}의 모든 공리가 성립함을 보였으므로 이는
분명하다.
\end{remark}

\noindent
$\mathbf{Q}$ 위의 스킴을 다루면 천 지표를 구성할 수 있다.

\begin{lemma}
\label{derham-lemma-chern-character}
모든 준콤팩트 준분리 스킴 $X$ 중 $\mathbf{Q}$ 위에 놓인 것에 ``천 지표''
$$
ch^{dR} : K_0(\textit{Vect}(X)) \longrightarrow
\prod\nolimits_{i \geq 0} H_{dR}^{2i}(X/\mathbf{Q})
$$
를 대응시키는 유일한 규칙이 존재하며, 다음 성질들을 갖는다.
\begin{enumerate}
\item $ch^{dR}$은 모든 $X$에 대해 환 준동형이다.
\item $f : X' \to X$가 $\mathbf{Q}$ 위 준콤팩트 준분리 스킴들의
사상이면 $f^* \circ ch^{dR} =  ch^{dR} \circ f^*$이다.
\item $\mathcal{L} \in \Pic(X)$가 주어지면
$ch^{dR}([\mathcal{L}]) = \exp(c_1^{dR}(\mathcal{L}))$이다.
\end{enumerate}
\end{lemma}

\noindent
이 구성은 $\mathbf{Q}$ 위의 모든 스킴으로 쉽게 확장할 수 있지만,
그러려면 베유 코호몰로지 이론의 절 \ref{weil-section-chern}의 논의를
약간 보강해야 한다.

\begin{proof}
보조정리 \ref{derham-lemma-chern-classes}의 증명에서와 똑같이,
$\mathbf{Q}$ 위 준콤팩트 준분리 스킴의 범주와 함자
$A^*(X) = \bigoplus_{i \geq 0} H_{dR}^{2i}(X/\mathbf{Q})$가 베유 코호몰로지 이론의 절 \ref{weil-section-chern}의 공리들을 만족함을
보일 수 있다. 더욱이 이 경우 $A(X)$는 $\mathbf{Q}$-대수이며, 이는 모든
$X$에 대해 성립한다. 따라서 보조정리는 베유 코호몰로지 이론의
명제 \ref{weil-proposition-chern-character}에서 따른다.
\end{proof}







\section{베유 코호몰로지 이론}
\label{derham-section-weil}

\noindent
$k$를 체라 하자. 그 표수는 $0$이라고 하자. 이 절에서는 함자
$$
X \longmapsto H^*_{dR}(X/k)
$$
가 베유 코호몰로지 이론의 정의
\ref{weil-definition-weil-cohomology-theory}에서 정의한 의미로 $k$ 위의
계수체 $k$인 베유 코호몰로지 이론을 정의함을 증명한다. 이 장 앞부분의
구성들이 베유 코호몰로지 이론의 절 \ref{weil-section-c1}의 공리
(A1)--(A9)를 만족하는 데이터 (D0), (D1), (D2')을 줌을 확인하겠다.

\medskip\noindent
이 절의 나머지 부분에서는 체 $k$를 고정하자. 그 표수는 $0$이고, $F = k$로
놓는다. 다음 데이터를 택한다.
\begin{enumerate}
\item[(D0)] $1$차원 $F$-벡터공간 $F(1)$로 $F(1) = F = k$를 택한다.
\item[(D1)] 함자 $H^*$로는 매끄러운 사영 스킴 $X$를 택하여, 이 스킴이
$k$ 위에 놓일 때 $H^*_{dR}(X/k)$로 보내는 함자를 택한다. 함자성은 절
\ref{derham-section-de-rham-cohomology}에서, 컵곱은 절
\ref{derham-section-cup-product}에서 논하였다. 보조정리
\ref{derham-lemma-cup-product-graded-commutative}에 의해 등급 가환
$F$-대수들을 얻는다.
\item[(D2')] 사상 $c_1^H : \Pic(X) \to H^2(X)(1)$로는 절
\ref{derham-section-first-chern-class}에서 도입한 드람 제1 천 특성류를 사용한다.
\end{enumerate}
공리 (A1)--(A9)가 성립함을 보이겠다.

\medskip\noindent
이 문단에서는 베유 코호몰로지 이론의 보조정리
\ref{weil-lemma-check-over-extension}를 사용하여 공리의 확인을 $k$가
대수적으로 닫힌 경우로 환원한다. $k'$를 $k$의 대수적 폐포라 하자.
$F' = k'$로 놓자. 위와 정확히 같은 방식으로 $k'$ 위에서 계수체가
$F'$인 데이터 (D0), (D1), (D2')을 얻는다. 보조정리
\ref{derham-lemma-proper-smooth-de-Rham}에 의해 함자적 동형들
$$
H_{dR}^{2d}(X/k) \otimes_k k'
\longrightarrow
H_{dR}^{2d}(X_{k'}/k')
$$
이 매끄러운 사영 스킴 $X$에 대해 존재하며, 이 스킴은 $k$ 위에 놓인다.
더욱이 도식
$$
\xymatrix{
\Pic(X) \ar[r]_{c^{dR}_1} \ar[d] & H_{dR}^2(X/k) \ar[d] \\
\Pic(X_{k'}) \ar[r]^{c^{dR}_1} & H_{dR}^2(X_{k'}/k')
}
$$
은 보조정리 \ref{derham-lemma-pullback-c1}에 의해 가환한다. 이로써 환원의
증명이 끝난다.

\medskip\noindent
$k$가 대수적으로 닫힌 체이고 표수가 영이라고 가정하자. 위에서 주어진
데이터 (D0), (D1), (D2')에 대해 공리 (A1)--(A9)를 보이겠다.

\medskip\noindent
공리 (A1). 여기서는
$H^*_{dR}(X \coprod Y/k) =  H^*_{dR}(X/k) \times H^*_{dR}(Y/k)$.
임을 확인해야 한다. 이는 드람 코호몰로지가 자리스키 위상에서 미분 등급
대수 층의 코호몰로지를 취하여 구성된다는 사실에서 따른다.

\medskip\noindent
공리 (A2). 이는 가역 가군의 제1 천 특성류를 취하는 연산이 당김과
양립한다는 명제이다. 더 일반적인 보조정리 \ref{derham-lemma-pullback-c1}에서
따른다.

\medskip\noindent
공리 (A3). 더 일반적인 명제
\ref{derham-proposition-projective-space-bundle-formula}에서 따른다.

\medskip\noindent
공리 (A4). 더 일반적인 보조정리
\ref{derham-lemma-log-complex-consequence}에서 따른다.

\medskip\noindent
이 단계에서 이미 베유 코호몰로지 이론의 보조정리
\ref{weil-lemma-chern-classes}와 \ref{weil-lemma-cycle-classes}를 사용하면
천 지표와 순환류 사상
$$
\gamma :
\CH^*(X)
\longrightarrow
\bigoplus\nolimits_{i \geq 0} H^{2i}_{dR}(X/k)
$$
을 얻는다. 이는 매끄러운 사영 스킴 $X$가 $k$ 위에 놓일 때 정의되는
등급 환 준동형이며, 매끄러운 사영 스킴들의 사상 $f : X \to Y$가
$k$ 위에 놓일 때 그 당김과 양립한다.

\medskip\noindent
공리 (A5). $H_{dR}^*(\Spec(k)/k) = k = F$가 차수 $0$에 놓인다. 보조정리
\ref{derham-lemma-kunneth-de-rham}에 의해 두 매끄러운 사영 $k$-스킴의 곱에
대한 퀸네트 공식이 성립한다. 명제의 유도 텐서곱들은 체 $k$ 위에서
텐서하므로 문제가 되지 않음에 유의하자.

\medskip\noindent
공리 (A7). 명제 \ref{derham-proposition-blowup-split}에서 따른다.

\medskip\noindent
공리 (A8). $X$를 매끄러운 사영 스킴이라 하되 $k$ 위에 놓인다고 하자.
베유 코호몰로지 이론의 절 \ref{weil-section-c1}에서 이 공리를 설명한
글에 의해 $k' = H^0(X, \mathcal{O}_X)$는 유한 분리 $k$-대수이다.
따라서 $H_{dR}^*(\Spec(k')/k) = k'$가 차수 $0$에 놓인다. 이는
$\Omega_{k'/k} = 0$이기 때문이다. 보조정리
\ref{derham-lemma-bottom-part-degenerates}에 의해서도
$H_{dR}^0(X, \mathcal{O}_X) = k'$이고, 이로써 공리를 얻는다.

\medskip\noindent
공리 (A6). $X$를 공집합이 아닌 매끄러운 사영 스킴이라 하고 $k$ 위에
놓인다고 하자. 또한 차원이 $d$인 등차원 스킴이라고 하자.
$\Delta : X \to X \times_{\Spec(k)} X$를 $X$의 $k$ 위 대각 사상이라
하자. $k$-선형 사상
$$
\lambda : H_{dR}^{2d}(X/k) \longrightarrow k
$$
이 존재하여 $(1 \otimes \lambda)\gamma([\Delta]) = 1$이
$H^0_{dR}(X/k)$에서 성립함을 보여야 한다. 다음과 같이 쓰자.
$$
\gamma = \gamma([\Delta]) = \gamma_0 + \ldots + \gamma_{2d}
$$
여기서 $\gamma_i \in H_{dR}^i(X/k) \otimes_k H_{dR}^{2d - i}(X/k)$는
퀸네트 성분이다. 문제는 선형 사상
$\lambda : H_{dR}^{2d}(X/k) \to k$가 존재하여
$(1 \otimes \lambda)\gamma_0 = 1$이 $H^0_{dR}(X/k)$에서 성립함을
보이는 것이다.

\medskip\noindent
$X = \coprod X_i$를 $X$의 연결 성분들, 따라서 기약 성분들로의
분해라 하자. 이에 대응하여 $\Delta = \coprod \Delta_i$이고
$\Delta_i \subset X_i \times X_i$이다. 따라서
$$
\gamma([\Delta]) = \sum \gamma([\Delta_i])
$$
이다. 더욱이 $\gamma([\Delta_i])$는 $\Delta_i \subset X_i \times X_i$의
류에 대응하며, 이 대응은 분해
$$
H^*_{dR}(X \times X) = \prod\nolimits_{i, j} H^*_{dR}(X_i \times X_j)
$$
를 통해 주어진다. 세부 사항은 생략한다. 이를 보이는 한 방법은
$\CH^0(X \times X)$ 안의 멱등원들 $e_{i, j}$를 사용하는 것이다. 이들은
열리고 닫힌 부분스킴 $X_i \times X_j$에 대응한다. 또한
$\gamma$가 $e_{i, j}$를 표시된 코호몰로지의 곱 분해에서 대응하는
멱등원으로 보내는 환 준동형임을 사용하는 것이다. 만일
$\lambda_i : H_{dR}^{2d}(X_i/k) \to k$를 찾아
$(1 \otimes \lambda_i)\gamma([\Delta_i]) = 1$이
$H^0_{dR}(X_i/k)$에서 성립하게 할 수 있다면,
$\lambda = \sum \lambda_i$를 택하여 $X$에 대한 문제를 해결할 수 있다.
따라서 $X$가 기약이라고 가정할 수 있고, 그렇게 가정한다.

\medskip\noindent
$X$가 기약인 경우의 공리 (A6)의 증명. $k$는 대수적으로 닫혀 있으므로
$H^0_{dR}(X/k) = k$이다. 실제로 $H^0(X, \mathcal{O}_X) = k$이다.
이는 $X$가 대수적으로 닫힌 체 위의 사영 다양체이기 때문이다. 예를 들어 대수다양체의
보조정리 \ref{varieties-lemma-proper-geometrically-reduced-global-sections}를
보라. $x \in X$를 임의의 닫힌점이라 하고 카르테시안 도식
$$
\xymatrix{
x \ar[d] \ar[r] & X \ar[d]^\Delta \\
X \ar[r]^-{x \times \text{id}} & X \times_{\Spec(k)} X
}
$$
을 생각하자. $\gamma$가 당김과 양립하므로 $\gamma([\Delta])$의 상은
$\gamma([x])$이며, 이는 $H_{dR}^{2d}(X/k)$ 안의 원소이다. 다시 말해
$\gamma_0 = 1 \otimes \gamma([x])$이다. 이로부터 두 가지를 얻는다.
(a) 류 $\gamma([x])$는 $x$와 무관하다. (b) 류 $\gamma([x])$가 영이
아님을 보이면 충분하다. 따라서 (c) 영차원 순환 $\alpha$를 $X$ 위에서
찾아 $\gamma(\alpha) \not = 0$가 되게 하면 충분하다. 이를 위해 유한
사상
$$
f : X \longrightarrow \mathbf{P}^d_k
$$
을 택하자. 이러한 사상이 존재함은 교차 이론의 절
\ref{intersection-section-projection}, 특히 보조정리
\ref{intersection-lemma-projection-generically-finite}를 보라. $f$는 유한
신토믹임에 유의하자. 사상 심화의 보조정리
\ref{more-morphisms-lemma-lci-permanence}에 의해 완전 국소 교차 사상이고,
가환대수의 보조정리 \ref{algebra-lemma-CM-over-regular-flat}에 의해
평탄하기 때문이다. 명제 \ref{derham-proposition-Garel}에 의해 대각합 사상
$$
\Theta_f :
f_*\Omega^\bullet_{X/k}
\longrightarrow
\Omega^\bullet_{\mathbf{P}^d_k/k}
$$
을 얻으며, 이를 표준적인 사상
$$
\Omega^\bullet_{\mathbf{P}^d_k/k}
\longrightarrow
f_*\Omega^\bullet_{X/k}
$$
과 합성하면 $f$의 차수를 곱하는 사상이 된다. 따라서 사상
$$
\Theta : H_{dR}^{2d}(X/k) \to H_{dR}^{2d}(\mathbf{P}^d_k/k)
$$
을 얻고, $\Theta \circ f^*$는 양의 정수를 곱하는 사상이다. 따라서
$\mathbf{P}^d_k$ 위에서 그 류가 영이 아닌 영차원 순환을 찾을 수 있다면,
$\gamma$가 당김과 양립함을 사용하여 결론을 얻는다. 이는 보조정리
\ref{derham-lemma-de-rham-cohomology-projective-space}에 의해 참이다. 이로써
공리 (A6)의 증명이 끝난다.

\medskip\noindent
아래에서는 다음 사실들을 더 언급하지 않고 사용한다. 첫째, 베유 코호몰로지 이론의 주의 \ref{weil-remark-trace}에 의해 사상
$\lambda_X : H^{2d}_{dR}(X/k) \to k$는 유일하다. 둘째, 공리 (A6)의
증명에서 $\lambda_X(\gamma([x])) = 1$임을 보았으며, 이때 $X$는 기약이다.
즉, 순환류 사상 $\gamma : \CH^d(X) \to H_{dR}^{2d}(X/k)$을
$\lambda_X$와 합성하면 차수 사상이 된다.

\medskip\noindent
공리 (A9). $Y \subset X$를 공집합이 아닌 매끄러운 인자라 하자. 여기서
$X$는 $k$ 위의 공집합이 아닌 매끄러운 등차원 사영 스킴이고 차원은
$d$이다. 도식
$$
\xymatrix{
H_{dR}^{2d - 2}(X/k)
\ar[rrr]_{c^{dR}_1(\mathcal{O}_X(Y)) \cap -} \ar[d]_{restriction} & & &
H_{dR}^{2d}(X) \ar[d]^{\lambda_X} \\
H_{dR}^{2d - 2}(Y/k) \ar[rrr]^-{\lambda_Y} & & & k
}
$$
이 가환함을 보여야 한다. 여기서 $\lambda_X$와 $\lambda_Y$는 공리
(A6)에서와 같다. 위에서 $X = \coprod X_i$를 연결 성분들, 동치로 기약
성분들로 분해하면 이에 대응하여 $\lambda_X = \sum \lambda_{X_i}$임을
보았다. 마찬가지로 $Y = \coprod Y_j$를 연결 성분들, 동치로 기약
성분들로 분해하면 $\lambda_Y = \sum \lambda_{Y_j}$이다. 더욱이 이 경우
$\mathcal{O}_X(Y) = \otimes_j \mathcal{O}_X(Y_j)$이고, 따라서
$$
c_1^{dR}(\mathcal{O}_X(Y)) = \sum\nolimits_j
c^{dR}_1(\mathcal{O}_X(Y_j))
$$
가 $H_{dR}^2(X/k)$ 안에서 성립한다. 직접 도식을 추적하면 $X$와 $Y$가
둘 다 기약인 경우에 도식의 가환성을 증명하면 충분함을 알 수 있다.
그러면 $H_{dR}^{2d - 2}(Y/k)$는 $1$차원이다. 이는 베유 코호몰로지 이론의 보조정리 \ref{weil-lemma-poincare-duality}에 의해 $Y$에 대한
푸앵카레 쌍대성이 성립하기 때문이다. 공리 (A4)에 의해
제한 사상, 즉 왼쪽 수직 화살표의 핵은
$c^{dR}_1(\mathcal{O}_X(Y))$와의 컵곱의 핵에 포함된다. 따라서 코호몰로지
류 $a \in H_{dR}^{2d - 2}(X)$를 하나 찾아, 이를 $Y$에 제한한 것이 영이
아니고 이 $a$에 대해 도식이 가환하도록 하면 충분하다. 임의의 풍부한
가역 가군 $\mathcal{L}$을 택하고
$$
a = c^{dR}_1(\mathcal{L})^{d - 1}
$$
$a|_Y = c^{dR}_1(\mathcal{L}|_Y)^{d - 1}$임을 알고 있으므로
$$
\lambda_Y(a|_Y) = \deg(c_1(\mathcal{L}|_Y)^{d - 1} \cap [Y])
$$
이다. 이는 위에서 기술한 $\lambda_Y$에서 따른다. 이 값은 Chow 호몰로지의 보조정리 \ref{chow-lemma-degrees-and-numerical-intersections}와
대수다양체의 보조정리 \ref{varieties-lemma-ample-positive}에 의해 양의
정수이다. 마찬가지로
$$
\lambda_X(c^{dR}_1(\mathcal{O}_X(Y)) \cap a) =
\deg(c_1(\mathcal{O}_X(Y)) \cap c_1(\mathcal{L})^{d - 1} \cap [X])
$$
를 얻는다. 정의에 의해 사실상
$c_1(\mathcal{O}_X(Y)) \cap [X] = [Y]$임을 알고 있으므로 영차원 순환들의
등식
$$
(Y \to X)_*\left(c_1(\mathcal{L}|_Y)^{d - 1} \cap [Y]\right) =
c_1(\mathcal{O}_X(Y)) \cap c_1(\mathcal{L})^{d - 1} \cap [X]
$$
을 $X$ 위에서 얻는다. 따라서 이 순환들은 차수가 같고 증명이 끝난다.

\begin{proposition}
\label{derham-proposition-de-rham-is-weil}
$k$를 표수 영인 체라 하자. 매끄러운 사영 스킴 $X$가 $k$ 위에 놓일 때
$H_{dR}^*(X/k)$로 보내는 함자는 베유 코호몰로지 이론의 정의
\ref{weil-definition-weil-cohomology-theory}의 의미에서 베유 코호몰로지
이론이다.
\end{proposition}

\begin{proof}
위의 논의에서 우리의 데이터 (D0), (D1), (D2')이 베유 코호몰로지 이론의 절 \ref{weil-section-c1}의 공리 (A1)--(A9)를 만족함을
보였다. 따라서 베유 코호몰로지 이론의 명제
\ref{weil-proposition-get-weil}에 의해 결론을 얻는다.

\medskip\noindent
다음 내용은 읽지 않기를 바란다. 위의 주장들을 증명할 때 보조정리들
\ref{derham-lemma-proper-smooth-de-Rham},
\ref{derham-lemma-pullback-c1},
\ref{derham-lemma-log-complex-consequence},
\ref{derham-lemma-kunneth-de-rham},
\ref{derham-lemma-bottom-part-degenerates}와
\ref{derham-lemma-de-rham-cohomology-projective-space},
명제들
\ref{derham-proposition-projective-space-bundle-formula},
\ref{derham-proposition-blowup-split}와
\ref{derham-proposition-Garel},
베유 코호몰로지 이론의 보조정리들
\ref{weil-lemma-check-over-extension},
\ref{weil-lemma-chern-classes},
\ref{weil-lemma-cycle-classes}와
\ref{weil-lemma-poincare-duality},
베유 코호몰로지 이론의 주의 \ref{weil-remark-trace}, 대수다양체의
보조정리들
\ref{varieties-lemma-proper-geometrically-reduced-global-sections}와
\ref{varieties-lemma-ample-positive},
교차 이론의 절 \ref{intersection-section-projection}과 보조정리
\ref{intersection-lemma-projection-generically-finite}, 사상 심화의
보조정리 \ref{more-morphisms-lemma-lci-permanence}, 가환대수의 보조정리
\ref{algebra-lemma-CM-over-regular-flat}, 그리고 Chow 호몰로지의 보조정리
\ref{chow-lemma-degrees-and-numerical-intersections}도 사용하였다.
\end{proof}

\begin{remark}
\label{derham-remark-hodge-cohomology-is-weil}
위와 정확히 같은 방식으로, 호지 코호몰로지
$X \mapsto H_{Hodge}^*(X/k)$가 $c_1^{Hodge}$를 갖추면 베유 코호몰로지
이론을 결정함을 보일 수 있다. 이 사실이 필요해지면 여기서 정확히
서술하고 증명하겠다. 이로부터
다음과 같은 흥미로운 결과가 나온다. 베유 코호몰로지 이론의 베티 수가
선택한 베유 코호몰로지 이론과 무관하다면, 우리의 체 $k$ 위에서 그
표수가 $0$일 때 호지-드람 스펙트럼 열은 $E_1$에서 퇴화한다. 물론 호지-드람
스펙트럼 열의 퇴화는 알려져 있다. 예를 들어 뛰어난 대수적 증명은
\cite{Deligne-Illusie}를 보라. 그러나 이는 결코 쉬운 결과가 아니다.
이는 베티 수의 무관성을 증명하는 일도 어려운 문제임을 시사하며, 우리가
아는 한 아직 미해결 문제이다. 관련 질문은 베유 코호몰로지 이론의
주의 \ref{weil-remark-betti-numbers-in-some-sense}를 보라.
\end{remark}







\section{닫힌 몰입에 대한 Gysin 사상}
\label{derham-section-gysin}

\noindent
이 절에서는 닫힌 몰입에 대한 Gysin 사상을 정의한다.

\begin{remark}
\label{derham-remark-gysin-equations}
$X \to S$를 스킴의 사상이라 하자.
$f_1, \ldots, f_c \in \Gamma(X, \mathcal{O}_X)$라 하자. $Z \subset X$를
$f_1, \ldots, f_c$로 정의되는 닫힌 부분스킴이라 하자. 아래에서는
다음 {\it Gysin 사상}을 연구한다.
\begin{equation}
\label{derham-equation-gysin}
\gamma^p_{f_1, \ldots, f_c} :
\Omega^p_{Z/S}
\longrightarrow
\mathcal{H}_Z^c(\Omega^{p + c}_{X/S})
\end{equation}
이 사상은 다음과 같이 정의된다. $\omega$를 $\Omega^p_{Z/S}$의 국소
단면이라 하고, 이것이 단면 $\tilde \omega$인 $\Omega^p_{X/S}$의 단면의
제한이면 다음과 같이 둔다.
$$
\gamma^p_{f_1, \ldots, f_c}(\omega) =
c_{f_1, \ldots, f_c}(\tilde \omega|_Z) \wedge
\text{d}f_1 \wedge \ldots \wedge \text{d}f_c
$$
여기서 $c_{f_1, \ldots, f_c} : \Omega^p_{X/S} \otimes \mathcal{O}_Z \to
\mathcal{H}_Z^c(\Omega^p_{X/S})$는 스킴의 유도 범주의 주의
\ref{perfect-remark-supported-map-c-equations}에서 구성한 사상이다.
이는 잘 정의된다. $\omega$가 주어졌을 때 $\tilde \omega$의 선택을
다음 꼴의 원소만큼 바꿀 수 있는데,
$\sum f_i \omega'_i + \sum \text{d}(f_i) \wedge \omega''_i$
이들은 위 구성에 의해 영으로 보내진다.
\end{remark}

\begin{lemma}
\label{derham-lemma-gysin-differential}
Gysin 사상 (\ref{derham-equation-gysin})은 $\Omega^\bullet_{X/S}$와
$\Omega^\bullet_{Z/S}$ 위의 드람 미분과 양립한다.
\end{lemma}

\begin{proof}
이를 올바르게 해석하면 거의 자명한 계산에서 결론이 따른다. 먼저
$\mathcal{H}^c_Z$를 $\mathcal{O}_X$-가군의 범주에서 계산한 함자는
$X$ 위 아벨 층의 범주에서 마찬가지로 정의한 함자와 일치함을 상기하자.
층 코호몰로지의 보조정리
\ref{cohomology-lemma-sections-support-abelian-unbounded}를 보라.
따라서 미분 $\text{d} : \Omega^p_{X/S} \to \Omega^{p + 1}_{X/S}$은 사상
$\mathcal{H}^c_Z(\Omega^p_{X/S}) \to \mathcal{H}^c_Z(\Omega^{p + 1}_{X/S})$을
유도한다. 더 나아가 스킴의 유도 범주의 주의
\ref{perfect-remark-support-c-equations}에 나오는 확장 교대 {\v C}ech
복합체의 구성은 아벨 층의 범주에서도 작동한다. 사상
$$
\Coker\left(\bigoplus \mathcal{F}_{1 \ldots \hat i \ldots c} \to
\mathcal{F}_{1 \ldots c}\right)
\longrightarrow
i_*\mathcal{H}^c_Z(\mathcal{F})
$$
은 스킴의 유도 범주의 주의
\ref{perfect-remark-supported-map-c-equations}에서
$c_{f_1, \ldots, f_c}$를 구성할 때 사용되었으며, $X$ 위 모든 아벨 층의
범주에서 잘 정의되고 함자적이다. 따라서 다음의 자명한 등식으로부터
보조정리가 따름을 알 수 있다.
$$
\text{d}\left(
\frac{\tilde \omega \wedge \text{d}f_1 \wedge \ldots \wedge
\text{d}f_c}{f_1 \ldots f_c}\right) =
\frac{\text{d}(\tilde \omega) \wedge
\text{d}f_1 \wedge \ldots \wedge \text{d}f_c}{f_1 \ldots f_c}
$$
\end{proof}

\begin{lemma}
\label{derham-lemma-gysin-global}
$X \to S$를 스킴의 사상이라 하자. $Z \to X$를 유한 표시 닫힌 몰입이라
하고, 그 여법층 $\mathcal{C}_{Z/X}$가 계수 $c$인 국소 자유층이라고 하자.
그러면 표준 사상
$$
\gamma^p : \Omega^p_{Z/S} \to \mathcal{H}^c_Z(\Omega^{p + c}_{X/S})
$$
이 존재하며, 국소적으로는 주의 \ref{derham-remark-gysin-equations}의 사상들
$\gamma^p_{f_1, \ldots, f_c}$로 주어진다.
\end{lemma}

\begin{proof}
가정에 의해 $x \in Z \subset X$가 주어지면 열린 근방 $U$가 존재하여
$x$를 포함하고, $Z$는 $c$개의 원소
$f_1, \ldots, f_c \in \mathcal{O}_X(U)$로 정의된다. 따라서
$f_1, \ldots, f_c$와 $g_1, \ldots, g_c$가 $\mathcal{O}_X(U)$ 안에서
$Z \cap U$를 정의할 때 사상 $\gamma^p_{f_1, \ldots, f_c}$와
$\gamma^p_{g_1, \ldots, g_c}$가 같음을 보이면 충분하다. 이를 위해
$U$를 줄인 뒤 $g_j = \sum a_{ji} f_i$라고 가정할 수 있으며, 여기서
$a_{ji} \in \mathcal{O}_X(U)$이다. 그러면 스킴의 유도 범주의 보조정리 \ref{perfect-lemma-supported-map-determinant}에 의해
$c_{f_1, \ldots, f_c} = \det(a_{ji}) c_{g_1, \ldots, g_c}$이다.
한편 다음 합동식이 성립한다.
$$
\text{d}(g_1) \wedge \ldots \wedge \text{d}(g_c) \equiv
\det(a_{ji}) \text{d}(f_1) \wedge \ldots \wedge \text{d}(f_c)
\bmod (f_1, \ldots, f_c)\Omega^c_{X/S}
$$
이 관계들을 결합하면 간단한 계산으로 원하는 등식을 얻는다.
\end{proof}

\begin{lemma}
\label{derham-lemma-gysin-differential-global}
$X \to S$와 $i : Z \to X$가 보조정리 \ref{derham-lemma-gysin-global}에서와
같다고 하자. Gysin 사상 $\gamma^p$는 $\Omega^\bullet_{X/S}$와
$\Omega^\bullet_{Z/S}$ 위의 드람 미분과 양립한다.
\end{lemma}

\begin{proof}
이는 국소적으로 확인할 수 있으며, 그러면 보조정리
\ref{derham-lemma-gysin-differential}에서 따른다.
\end{proof}

\begin{lemma}
\label{derham-lemma-gysin-projection}
$X \to S$와 $i : Z \to X$가 보조정리 \ref{derham-lemma-gysin-global}에서와
같다고 하자. $\alpha \in H^q(X, \Omega^p_{X/S})$가 주어지면
$\gamma^p(\alpha|_Z) = i^{-1}\alpha \wedge \gamma^0(1)$가
$H^q(Z, \mathcal{H}^c_Z(\Omega^{p + c}_{X/S}))$에서 성립한다.
표기는 증명을 보라.
\end{lemma}

\begin{proof}
제한 $\alpha|_Z$는 호지 코호몰로지의 함자성으로 주어지는
$H^q(Z, \Omega^p_{Z/S})$의 원소이다. 다음 사상을 사용하여 코호몰로지의
함자성을 적용하면
$\gamma^p : \Omega^p_{Z/S} \to \mathcal{H}^c_Z(\Omega^{p + c}_{X/S})$
$\gamma^p(\alpha|_Z)$를
$H^q(Z, \mathcal{H}^c_Z(\Omega^{p + c}_{X/S}))$의 원소로 얻는다.
이로써 공식의 좌변이 설명된다.

\medskip\noindent
우변을 설명하기 위해 먼저 환 달린 공간의 사상
$i : (Z, i^{-1}\mathcal{O}_X) \to (X, \mathcal{O}_X)$으로 당겨 원소
$i^{-1}\alpha \in H^q(Z, i^{-1}\Omega^p_{X/S})$를 얻는다.
$\gamma^0(1) \in H^0(Z, \mathcal{H}_Z^c(\Omega^c_{X/S}))$를
$1 \in H^0(Z, \mathcal{O}_Z) = H^0(Z, \Omega^0_{Z/S})$의
$\gamma^0$에 의한 상이라 하자.
컵곱을 사용하면 원소
$$
i^{-1}\alpha \cup \gamma^0(1)
\in
H^{q + c}(Z, 
i^{-1}\Omega^p_{X/S} \otimes_{i^{-1}\mathcal{O}_X}
\mathcal{H}^c_Z(\Omega^c_{X/S}))
$$
층 코호몰로지의 주의 \ref{cohomology-remark-support-cup-product}와 쐐기곱을
사용하면 표준 사상들
$$
i^{-1}\Omega^p_{X/S} \otimes_{i^{-1}\mathcal{O}_X}^\mathbf{L}
R\mathcal{H}_Z(\Omega^c_{X/S}) \to
R\mathcal{H}_Z(\Omega^p_{X/S} \otimes_{\mathcal{O}_X}^\mathbf{L}
\Omega^c_{X/S}) \to
R\mathcal{H}_Z(\Omega^{p + c}_{X/S})
$$
이 존재한다. 스킴의 유도 범주의 보조정리
\ref{perfect-lemma-supported-trivial-vanishing}에 의해 대상
$R\mathcal{H}_Z(\Omega^j_{X/S})$의 차수 $> c$인 코호몰로지 층들은
소멸한다. 따라서 차수 $c$의 코호몰로지 층 위에서 사상
$$
i^{-1}\Omega^p_{X/S} \otimes_{i^{-1}\mathcal{O}_X}
\mathcal{H}^c_Z(\Omega^c_{X/S}) \longrightarrow
\mathcal{H}^c_Z(\Omega^{p + c}_{X/S})
$$
을 얻는다. 식 $i^{-1}\alpha \wedge \gamma^0(1)$은 코호몰로지의
함자성에 따른 컵곱 $i^{-1}\alpha \cup \gamma^0(1)$의 상이다.

\medskip\noindent
이와 같이 공식의 내용을 설명했으므로, 이제 컵곱의 일반 성질
(층 코호몰로지의 절 \ref{cohomology-section-cup-product})에 의해 다음
그림이 가환임을 보이면 충분하다.
$$
\xymatrix{
i^{-1}\Omega^p_X \otimes \Omega^0_Z \ar[rr]_{\text{id} \otimes \gamma^0}
\ar[d] & &
i^{-1}\Omega^p_X \otimes \mathcal{H}^c_Z(\Omega^c_X) \ar[d]^\wedge \\
\Omega^p_Z \otimes \Omega^0_Z \ar[r]^\wedge &
\Omega^p_Z \ar[r]^{\gamma^p} &
\mathcal{H}^c_Z(\Omega^{p + c}_X)
}
$$
여기서 가환성은 $Z$ 위 층의 범주에서 말하며, 표기를 자명하게
남용하였다. 이는 사상들 $\gamma^j_{f_1, \ldots, f_c}$에 대한 간단한
계산으로 귀착되므로 생략한다. 실제로 이 사상들은 정확히 이 계산이
성립하고 $1$이
$\frac{\text{d}f_1 \wedge \ldots \wedge \text{d}f_c}{f_1 \ldots f_c}$로
보내지도록 선택되었다.
\end{proof}

\begin{lemma}
\label{derham-lemma-gysin-transverse}
$c \geq 0$을 정수라 하자. 다음을 스킴의 가환 그림이라 하자.
$$
\xymatrix{
Z' \ar[d]_h \ar[r] & X' \ar[d]_g \ar[r] & S' \ar[d] \\
Z \ar[r] & X \ar[r] & S
}
$$
다음을 가정한다.
\begin{enumerate}
\item $Z \to X$와 $Z' \to X'$는 보조정리
\ref{derham-lemma-gysin-global}의 가정을 만족한다.
\item 그림의 왼쪽 정사각형은 데카르트이다.
\item $h^*\mathcal{C}_{Z/X} \to \mathcal{C}_{Z'/X'}$
(스킴의 사상의 보조정리 \ref{morphisms-lemma-conormal-functorial})는
동형이다.
\end{enumerate}
그러면 그림
$$
\xymatrix{
h^*\Omega^p_{Z/S} \ar[rr]_-{h^{-1}\gamma^p} \ar[d] & &
\mathcal{O}_{X'}|_{Z'} \otimes_{h^{-1}\mathcal{O}_X|_Z}
h^{-1}\mathcal{H}^c_Z(\Omega^{p + c}_{X/S}) \ar[d] \\
\Omega^p_{Z'/S'} \ar[rr]^{\gamma^p} & &
\mathcal{H}^c_{Z'}(\Omega^{p + c}_{X'/S'})
}
$$
은 가환이다. 왼쪽 수직 화살표는 미분 가군의 함자성이고, 오른쪽 수직
화살표에는 층 코호몰로지의 주의
\ref{cohomology-remark-support-functorial}을 사용한다.
\end{lemma}

\begin{proof}
더 정확히는 다음 합성을 생각하자.
\begin{align*}
\mathcal{O}_{X'}|_{Z'} \otimes_{h^{-1}\mathcal{O}_X|_Z}^\mathbf{L}
h^{-1}R\mathcal{H}_Z(\Omega^{p + c}_{X/S})
& \to
R\mathcal{H}_{Z'}(Lg^*\Omega^{p + c}_{X/S}) \\
& \to
R\mathcal{H}_{Z'}(g^*\Omega^{p + c}_{X/S}) \\
& \to
R\mathcal{H}_{Z'}(\Omega^{p + c}_{X'/S'})
\end{align*}
첫 번째 화살표는 층 코호몰로지의 주의
\ref{cohomology-remark-support-functorial}에서 주어지고, 마지막 화살표는
미분의 함자성에서 주어진다. 스킴의 유도 범주의 보조정리
\ref{perfect-lemma-supported-trivial-vanishing}에 의해 차수 $> c$에서
코호몰로지 층이 소멸하므로, 이 합성은 오른쪽 수직 화살표를 유도한다.
가환성은 국소적으로 확인할 수 있다. 따라서 $Z$가
$f_1, \ldots, f_c \in \Gamma(X, \mathcal{O}_X)$로 정의된다고 가정할 수
있다. 그러면 $Z'$는 $f'_i = g^\sharp(f_i)$로 정의된다.
사상들 $c_{f_1, \ldots, f_c}$와 $c_{f'_1, \ldots, f'_c}$는 다음 가환
그림을 이룬다.
$$
\xymatrix{
h^*i^*\Omega^p_{X/S} \ar[rr]_-{h^{-1}c_{f_1, \ldots, f_c}} \ar[d] & &
\mathcal{O}_{X'}|_{Z'} \otimes_{h^{-1}\mathcal{O}_X|_Z}
h^{-1}\mathcal{H}^c_Z(\Omega^p_{X/S}) \ar[d] \\
(i')^*\Omega^p_{X'/S'} \ar[rr]^{c_{f'_1, \ldots, f'_c}} & &
\mathcal{H}^c_{Z'}(\Omega^p_{X'/S'})
}
$$
스킴의 유도 범주의 주의
\ref{perfect-remark-supported-functorial}을 보라. $p$-형식 $\omega$가
$Z$ 위에 주어지면, $\gamma^p(\omega)$를 정의하기 위해 ($X$와 $Z$
위에서 국소적으로) $p$-형식 $\tilde \omega$를 $X$ 위에서 선택하여
$\omega$를 올리고 다음과 같이 두었음을 상기하자.
$\gamma^p(\omega) =
c_{f_1, \ldots, f_c}(\tilde \omega) \wedge
\text{d}f_1 \wedge \ldots \wedge \text{d}f_c$.
형식 $\text{d}f_1 \wedge \ldots \wedge \text{d}f_c$를 당기면
$\text{d}f'_1 \wedge \ldots \wedge \text{d}f'_c$가 되므로 결론을 얻는다.
\end{proof}

\begin{remark}
\label{derham-remark-how-to-use}
$X \to S$, $i : Z \to X$, $c \geq 0$이 보조정리
\ref{derham-lemma-gysin-global}에서와 같다고 하자.
$p \geq 0$이라 하고 $\mathcal{H}^i_Z(\Omega^{p + c}_{X/S}) = 0$이
$i = 0, \ldots, c - 1$에 대해 성립한다고 가정하자. $X \to S$가 매끄럽고
$Z \to X$가 코쥘 정칙 몰입이면 이 소멸이 성립한다. 스킴의 유도 범주의 보조정리 \ref{perfect-lemma-supported-vanishing}을 보라.
그러면 다음 사상
$$
\gamma^{p, q} :
H^q(Z, \Omega^p_{Z/S})
\longrightarrow
H^{q + c}(X, \Omega^{p + c}_{X/S})
$$
을 얻는다. 먼저 사상
$\gamma^p : \Omega^p_{Z/S} \to \mathcal{H}^c_Z(\Omega^{p + c}_{X/S})$를
사용하여
$$
H^q(Z, \mathcal{H}^c_Z(\Omega^{p + c}_{X/S})) =
H^q(Z, R\mathcal{H}_Z(\Omega^{p + c}_{X/S})[c]) =
H^q(X, i_*R\mathcal{H}_Z(\Omega^{p + c}_{X/S})[c])
$$
으로 보낸 다음, 수반 사상
$i_*R\mathcal{H}_Z(\Omega^{p + c}_{X/S}) \to \Omega^{p + c}_{X/S}$을
사용하여 원하는 호지 코호몰로지 가군으로 계속 보낸다.
\end{remark}

\begin{lemma}
\label{derham-lemma-gysin-differential-hodge}
$X \to S$와 $i : Z \to X$가 보조정리 \ref{derham-lemma-gysin-global}에서와
같다고 하자. $X \to S$가 매끄럽고 $Z \to X$가 코쥘 정칙이라고
가정하자. Gysin 사상들 $\gamma^{p, q}$는 $\Omega^\bullet_{X/S}$와
$\Omega^\bullet_{Z/S}$ 위의 드람 미분과 양립한다.
\end{lemma}

\begin{proof}
이는 보조정리 \ref{derham-lemma-gysin-differential-global}에서 곧바로 따른다.
\end{proof}

\begin{lemma}
\label{derham-lemma-gysin-projection-global}
$X \to S$, $i : Z \to X$, $c \geq 0$이 보조정리
\ref{derham-lemma-gysin-global}에서와 같다고 하자. $X \to S$가 매끄럽고
$Z \to X$가 코쥘 정칙이라고 가정하자.
$\alpha \in H^q(X, \Omega^p_{X/S})$가 주어지면
$\gamma^{p, q}(\alpha|_Z) = \alpha \cup \gamma^{0, 0}(1)$가
$H^{q + c}(X, \Omega^{p + c}_{X/S})$에서 성립한다. 여기서
$\gamma^{a, b}$는 주의 \ref{derham-remark-how-to-use}에서와 같다.
\end{lemma}

\begin{proof}
이 보조정리는 보조정리 \ref{derham-lemma-gysin-projection}과 층 코호몰로지의
보조정리 \ref{cohomology-lemma-support-cup-product}에서 따른다.
독자에게 아래의 더 자세한 논의를 건너뛰기를 권한다.

\medskip\noindent
이후 따로 말하지 않고
$R\mathcal{H}_Z(\Omega^j_{X/S}) = \mathcal{H}^c_Z(\Omega^j_{X/S})[-c]$
$j$ 모두에 대해 성립함을 사용할 것이다. 이는 주의
\ref{derham-remark-how-to-use}에서 지적하였다. 또한 다음 동일시들을 명시하지
않고 사용할 것이다.
$H^{q + c}_Z(X, \Omega^j_{X/S}) = H^{q + c}(Z, R\mathcal{H}_Z(\Omega^j_{X/S}) =
H^q(Z, \mathcal{H}^c_Z(\Omega^j_{X/S}))$.
첫 번째 동일시에 대해서는 층 코호몰로지의 보조정리
\ref{cohomology-lemma-local-to-global-sections-with-support}를 보라.
이 동일시들 아래에서 다음이 성립한다.
\begin{enumerate}
\item $\gamma^0(1) \in H^c_Z(X, \Omega^c_{X/S})$은
$\gamma^{0, 0}(1)$인 $H^c(X, \Omega^c_{X/S})$의 원소로 보내진다.
\item 보조정리 \ref{derham-lemma-gysin-projection}의 등식 우변
$i^{-1}\alpha \wedge \gamma^0(1)$은 층 코호몰로지의 주의
\ref{cohomology-remark-support-cup-product-global}에 나오는 원소
$\alpha$와 $\gamma^0(1)$의 컵곱을 쐐기곱으로 보낸 상이다. 다시 말해
보조정리 \ref{derham-lemma-gysin-projection}의 증명에서의 구성과 층 코호몰로지의
주의 \ref{cohomology-remark-support-cup-product-global}에서의 구성이
일치한다.
\item 층 코호몰로지의 보조정리
\ref{cohomology-lemma-support-cup-product}에 의해 이는
$\alpha \cup \gamma^{0, 0}(1)$인
$H^{q + c}(X, \Omega^p_{X/S} \otimes \Omega^c_{X/S})$의 원소로 보내진다.
\item 보조정리 \ref{derham-lemma-gysin-projection}의 등식 좌변
$\gamma^p(\alpha|_Z)$은 $\gamma^{p, q}(\alpha|_Z)$로 보내진다.
\end{enumerate}
이로써 증명이 끝난다.
\end{proof}

\begin{lemma}
\label{derham-lemma-gysin-transverse-global}
$c \geq 0$이라 하고
$$
\xymatrix{
Z' \ar[d]_h \ar[r] & X' \ar[d]_g \ar[r] & S' \ar[d] \\
Z \ar[r] & X \ar[r] & S
}
$$
이 보조정리 \ref{derham-lemma-gysin-transverse}의 가정을 만족한다고 하자. 또한
$X \to S$와 $X' \to S'$가 매끄럽고 $Z \to X$와 $Z' \to X'$가 코쥘
정칙 몰입이라고 가정하자. 그러면 그림
$$
\xymatrix{
H^q(Z, \Omega^p_{Z/S}) \ar[rr]_-{\gamma^{p, q}} \ar[d] & &
H^{q + c}(X, \Omega^{p + c}_{X/S}) \ar[d] \\
H^q(Z', \Omega^p_{Z'/S'}) \ar[rr]^{\gamma^{p, q}} & &
H^{q + c}(X', \Omega^{p + c}_{X'/S'})
}
$$
은 가환이다. 여기서 $\gamma^{p, q}$는 주의
\ref{derham-remark-how-to-use}에서와 같다.
\end{lemma}

\begin{proof}
이는 보조정리 \ref{derham-lemma-gysin-transverse}와 층 코호몰로지의 보조정리
\ref{cohomology-lemma-support-functorial}을 결합하면 따른다.
\end{proof}

\begin{lemma}
\label{derham-lemma-class-of-a-point}
$k$를 체라 하자. $X$를 $k$ 위의 차원 $d$인 기약 매끄러운 고유
스킴이라 하자. $Z \subset X$를 하나의 $k$-유리점 $x$로 이루어진
축약 닫힌 부분스킴이라 하자. 그러면
$1 \in k = H^0(Z, \mathcal{O}_Z) = H^0(Z, \Omega^0_{Z/k})$
의 주의 \ref{derham-remark-how-to-use}의 사상
$H^0(Z, \Omega^0_{Z/k}) \to H^d(X, \Omega^d_{X/k})$에 의한 상은
영이 아니다.
\end{lemma}

\begin{proof}
사상 $\gamma^0 : \mathcal{O}_Z \to
\mathcal{H}^d_Z(\Omega^d_{X/k}) = R\mathcal{H}_Z(\Omega^d_{X/k})[d]$는
사상
$$
g^0 : i_*\mathcal{O}_Z \longrightarrow \Omega^d_{X/k}[d]
$$
의 $D(\mathcal{O}_X)$에서의 수반이다. $\Omega^d_{X/k} = \omega_X$가
$X/k$의 쌍대화 층임을 상기하자. 스킴의 쌍대성의 보조정리
\ref{duality-lemma-duality-proper-over-field}를 보라. 따라서 보조정리의
명제에 나오는 사상의 $k$-선형 쌍대는 사상
$$
H^0(X, \mathcal{O}_X) \to \Ext^d_X(i_*\mathcal{O}_Z, \omega_X)
$$
이다. 이는 $1$을 $g^0$로 보낸다. 따라서 $g^0$가 영이 아님을 보이면
충분하다. 이는 임의의 근방 $U$에서 확인할 수 있는데, 이 근방은 점
$x$의 근방이다. $U$를
선택하여 $f_1, \ldots, f_d \in \mathcal{O}_X(U)$가 $x$에서만 소멸하고
극대 아이디얼 $\mathfrak m_x \subset \mathcal{O}_{X, x}$을 생성하게
하자. 아핀인 $U = \Spec(R)$라고 가정할 수 있다고 가정할 수 있다.
$\gamma^0$의 구성을 살펴보면 우리의 확장은 다음과 같이 주어진다.
$$
k \to
(R \to \bigoplus\nolimits_{i_0} R_{f_{i_0}} \to
\bigoplus\nolimits_{i_0 < i_1} R_{f_{i_0}f_{i_1}} \to
\ldots \to R_{f_1\ldots f_r})[d] \to R[d]
$$
여기서 $1$은 첫 번째 사상에 의해 $1/f_1 \ldots f_c$로 보내진다.
이 경우 $1/f_1 \ldots f_c$는 국소 코호몰로지 군
$H^d_{(f_1, \ldots, f_d)}(R)$의 영이 아닌 원소이므로, 이는 영이 아니다.
\end{proof}






\section{상대 푸앵카레 쌍대성}
\label{derham-section-relative-poincare-duality}

\noindent
이 절에서는 바탕 위 고유 매끄러운 스킴의 상대 드람 코호몰로지에 대한
푸앵카레 쌍대성을 증명한다. 독자에게 먼저 절
\ref{derham-section-poincare-duality}을 읽기를 강하게 권한다.

\begin{situation}
\label{derham-situation-relative-duality}
여기서 $S$는 준콤팩트 준분리 스킴이고, $f : X \to S$는 모든 올이
공집합이 아니며 차원 $n$인 등차원 스킴이 되는 고유 매끄러운
스킴 사상이다.
\end{situation}

\begin{lemma}
\label{derham-lemma-relative-bottom-part-degenerates}
상황 \ref{derham-situation-relative-duality}에서 직상 $f_*\mathcal{O}_X$는
유한 에탈 $\mathcal{O}_S$-대수이고, $S$ 위에서 국소적으로
$Rf_*\mathcal{O}_X = f_*\mathcal{O}_X \oplus P$가
$D(\mathcal{O}_S)$ 안에서 성립한다. 여기서
$P$는 Tor 진폭이 $[1, \infty)$에 놓이는 완전 대상이다. 사상
$\text{d} : f_*\mathcal{O}_X \to f_*\Omega_{X/S}$는 영이다.
\end{lemma}

\begin{proof}
명제의 첫 부분은 스킴의 유도 범주의 보조정리
\ref{perfect-lemma-proper-flat-geom-red}에서 따른다.
$S' = \underline{\Spec}_S(f_*\mathcal{O}_X)$로 두면 분해
$X \to S' \to S$를 얻는다. 이는 스타인 분해이지만 여기서는 그 사실이
필요하지 않다. 사상 심화의 절
\ref{more-morphisms-section-stein-factorization}을 보라. 예를 들어
스킴의 사상의 보조정리 \ref{morphisms-lemma-triangle-differentials}와
\ref{morphisms-lemma-etale-at-point}에 의해
$\Omega_{X/S} = \Omega_{X/S'}$임을 알 수 있다. 물론 이는
$\text{d} : f_*\mathcal{O}_X \to f_*\Omega_{X/S}$가 영임을 뜻한다.
\end{proof}

\begin{lemma}
\label{derham-lemma-relative-duality-hodge}
상황 \ref{derham-situation-relative-duality}에서 $\mathcal{O}_S$-가군 사상
$$
t : Rf_*\Omega^n_{X/S}[n] \longrightarrow \mathcal{O}_S
$$
이 존재한다. 이는 $H^0(X, \mathcal{O}_X)$의 가역원을 곱하는 사상을
앞에 합성하는 것까지 유일하며, 다음 성질을 갖는다. 모든 $p$에 대해 쌍
$$
Rf_*\Omega^p_{X/S}
\otimes_{\mathcal{O}_S}^\mathbf{L}
Rf_*\Omega^{n - p}_{X/S}[n]
\longrightarrow
\mathcal{O}_S
$$
은 상대 컵곱과 $t$의 합성으로 주어지며, $S$ 위 완전 복합체들의
완전한 쌍이다.
\end{lemma}

\begin{proof}
먼저 $S$가 뇌터 스킴이라고 가정하자. $\omega^\bullet_{X/S}$를
스킴의 쌍대성의 주의
\ref{duality-remark-relative-dualizing-complex}에서처럼 $X$의 $S$ 위
상대 쌍대화 복합체라 하고,
$Rf_*\omega_{X/S}^\bullet \to \mathcal{O}_S$를 그 대각합 사상이라 하자.
스킴의 쌍대성의 보조정리 \ref{duality-lemma-smooth-proper}에 의해
(여기서 $S$가 뇌터라는 사실을 사용한다) 동형
$\omega^\bullet_{X/S} \cong \Omega^n_{X/S}[n]$이 존재하며, 이 동형을
사용하여 $t$를 얻는다. 복합체 $Rf_*\Omega^p_{X/S}$는 보조정리
\ref{derham-lemma-proper-smooth-de-Rham}에 의해 완전하다.
$\Omega^p_{X/S}$가 국소 자유이고,
$\Omega^p_{X/S} \otimes_{\mathcal{O}_X} \Omega^{n - p}_{X/S} \to
\Omega^n_{X/S}$가 동형
$\Omega^p_{X/S} \cong
\SheafHom_{\mathcal{O}_X}(\Omega^{n - p}_{X/S}, \Omega^n_{X/S})$을
나타내므로, 상대 컵곱으로 유도되는 쌍은 스킴의 쌍대성의 주의
\ref{duality-remark-relative-dualizing-complex-relative-cup-product}에 의해
완전하다.

\medskip\noindent
일반적인 경우의 존재성 증명. 절대 뇌터 근사에 의해 데카르트 그림
$$
\xymatrix{
X \ar[r] \ar[d]^f & X_0 \ar[d]^{f_0} \\
S \ar[r]^g & S_0
}
$$
을 찾을 수 있다. 여기서 $S_0$는 뇌터이고, $f_0$는 모든 올이 공집합이
아니며 차원 $n$인 등차원 스킴이 되는 고유 매끄러운 사상이다.
스킴의 극한의 보조정리들 \ref{limits-lemma-descend-finite-presentation},
\ref{limits-lemma-descend-smooth}, \ref{limits-lemma-descend-dimension-d},
\ref{limits-lemma-descend-surjective}, 그리고
\ref{limits-lemma-eventually-proper}를 보라. 사상
$$
t_0 : Rf_{0, *}\Omega^n_{X_0/S_0}[n] \longrightarrow \mathcal{O}_{S_0}
$$
을 보조정리의 명제에서처럼 선택하자. 코호몰로지와 바탕변환에 의해
복합체 $Rf_*(\Omega^p_{X/S})$는
$Lg^*(Rf_{0, *}\Omega^p_{X_0/S_0})$와 같으며, 이는 모든
$p \geq 0$에 대해 성립한다. 보조정리
\ref{derham-lemma-proper-smooth-de-Rham}을 보라. 특히 $t_0$를 $g$로 당기면
보조정리의 명제에서와 같은 사상 $t$를 얻는다. 실제로 층 코호몰로지의
보조정리 \ref{cohomology-lemma-base-change-cup-product}에 의해 컵곱 사상
$$
Rf_*\Omega^p_{X/S}
\otimes_{\mathcal{O}_S}^\mathbf{L}
Rf_*\Omega^{n - p}_{X/S}[n]
\longrightarrow
Rf_*\Omega^n_{X/S}
$$
은 $g$로 당긴 $f_0$에 대한 대응하는 컵곱 사상이다. 따라서
$t_0$가 완전한 쌍들을 준다는 사실을 당기면 $t$에 대해서도 $S$ 위에서
같은 성질을 얻는다.

\medskip\noindent
이제 $t$의 유일성을 보이자. 특별 삼각형
$f_*\mathcal{O}_X \to Rf_*\mathcal{O}_X \to P \to f_*\mathcal{O}_X[1]$.
을 선택하자. 보조정리 \ref{derham-lemma-relative-bottom-part-degenerates}에 의해
대상 $P$는 Tor 진폭이 $[1, \infty)$에 놓이는 완전 대상이고, 이
삼각형은 $S$ 위에서 국소적으로 분할된다. 따라서
$R\SheafHom_{\mathcal{O}_X}(P, \mathcal{O}_X)$는 Tor 진폭이
$(-\infty, -1]$에 놓이는 완전 대상이다. 그러므로 위 쌍대성에 의해
$S$ 위에서 국소적으로
$$
Rf_*\Omega^n_{X/S}[n] \cong
R\SheafHom_{\mathcal{O}_S}(f_*\mathcal{O}_X, \mathcal{O}_S)
\oplus R\SheafHom_{\mathcal{O}_X}(P, \mathcal{O}_X)
$$
가 성립한다. 이로부터 $R^nf_*\Omega^n_{X/S}$가 유한 국소 자유이고,
완전한 $\mathcal{O}_S$-쌍선형 쌍
$$
f_*\mathcal{O}_X \times R^nf_*\Omega^n_{X/S} \longrightarrow \mathcal{O}_S
$$
을 $t$를 사용하여 얻는다. 따라서 임의의 $\mathcal{O}_S$-선형 사상
$t' : R^nf_*\Omega^n_{X/S} \to \mathcal{O}_S$는
$t' = t \circ g$ 꼴이며, 여기서 어떤
$g \in \Gamma(S, f_*\mathcal{O}_X) = \Gamma(X, \mathcal{O}_X)$이다.
$t'$가 여전히 완전한 쌍을 정하려면 $g$가
가역원이어야 한다. 이로써 증명이 끝난다.
\end{proof}

\begin{lemma}
\label{derham-lemma-relative-top-part-degenerates}
상황 \ref{derham-situation-relative-duality}에서 사상
$\text{d} : R^nf_*\Omega^{n - 1}_{X/S} \to R^nf_*\Omega^n_{X/S}$
는 영이다.
\end{lemma}

\noindent
보조정리 \ref{derham-lemma-top-part-degenerates}의 증명에서 언급했듯이, 이
보조정리는 보조정리들 \ref{derham-lemma-relative-duality-hodge}와
\ref{derham-lemma-relative-bottom-part-degenerates}.
의 쉬운 귀결이 아니다.

\begin{proof}[$S$가 축약인 경우의 증명]
$S$가 축약이라고 가정하자. 다음 사상
$\text{d} : R^nf_*\Omega^{n - 1}_{X/S} \to R^nf_*\Omega^n_{X/S}$
은 (준연접) $\mathcal{O}_S$-가군의 $\mathcal{O}_S$-선형 사상이다.
$\mathcal{O}_S$-가군 $R^nf_*\Omega^n_{X/S}$는 유한 국소 자유이다.
실제로 보조정리들
\ref{derham-lemma-relative-duality-hodge}와
\ref{derham-lemma-relative-bottom-part-degenerates}에 의해 유한 국소 자유
$\mathcal{O}_S$-가군 $f_*\mathcal{O}_X$의 쌍대이다.
$S$가 축약이므로 $\text{d}$의 줄기가 모든 일반점 $\eta \in S$에서
영임을 보이면 충분하다. 이는 아핀 열린집합 위의 단면들을 살펴보고,
$\text{d}$의 공역이 국소 자유라는 사실과 가환대수의 보조정리
\ref{algebra-lemma-reduced-ring-sub-product-fields}의 (2)를 사용하면
따른다. $S$가 축약이므로
$\mathcal{O}_{S, \eta} = \kappa(\eta)$이다. 가환대수의 보조정리
\ref{algebra-lemma-minimal-prime-reduced-ring}을 보라. 따라서
$\text{d}_\eta$는 사상
$$
\text{d} :
H^n(X_\eta, \Omega^{n - 1}_{X_\eta/\kappa(\eta)})
\longrightarrow
H^n(X_\eta, \Omega^n_{X_\eta/\kappa(\eta)})
$$
과 동일시되며, 이는 보조정리 \ref{derham-lemma-top-part-degenerates}에 의해
영이다.
\end{proof}

\begin{proof}[일반적인 경우의 증명]
이 문제는 $S$ 위 평탄 국소적이다. 즉, $S' \to S$가 스킴의 전사 평탄
사상이고 사상을 $S'$로 당긴 뒤 영이면 원래 사상도 영이다. 또한 평탄
바탕변환에 의해 이 사상의 구성은 평탄 사상에 의한 바탕변환과
가환한다(스킴의 코호몰로지의 보조정리
\ref{coherent-lemma-flat-base-change-cohomology}).

\medskip\noindent
사상 심화의 정리
\ref{more-morphisms-theorem-stein-factorization-general}에서처럼 스타인
분해 $X \to S' \to S$를 생각하자. 보조정리
\ref{derham-lemma-relative-bottom-part-degenerates}에 의해 사상
$\pi : S' \to S$는 유한 에탈이다. 사상 $f : X \to S'$는 이 정리에
의해 고유이고, 사상 심화의 보조정리
\ref{more-morphisms-lemma-smooth-etale-permanence}에 의해 매끄러우며,
스타인 분해 정리에 의해 기하적으로 연결된 올들을 갖는다. 보조정리
\ref{derham-lemma-relative-bottom-part-degenerates}의 증명에서
$\Omega_{X/S} = \Omega_{X/S'}$임을 보았는데, 이는 $S' \to S$가
에탈이기 때문이다. 따라서
$\Omega^\bullet_{X/S} = \Omega^\bullet_{X/S'}$이다. 또한
$$
R^qf_*\Omega^p_{X/S} = \pi_*R^qf'_*\Omega^p_{X/S'}
$$
가 모든 $p, q$에 대해 성립한다. 이는 르레 스펙트럼 열(층 코호몰로지의
보조정리 \ref{cohomology-lemma-relative-Leray}), $\pi$가 유한이어서
아핀이라는 사실, 그리고 스킴의 코호몰로지의 보조정리
\ref{coherent-lemma-relative-affine-vanishing}에서 따른다. 물론
$R^qf'_*\Omega^p_{X'/S}$가 준연접이라는 사실도 사용한다. 따라서
보조정리의 사상은 $\pi_*$를
$\text{d} : R^nf'_*\Omega^{n - 1}_{X/S'} \to R^nf'_*\Omega^n_{X/S'}$에
적용한 것이다. 다시 말해 보조정리를 증명하기 위해
$f : X \to S$를 $f' : X \to S'$로 바꾸어 다음 문단에서 논할
경우로 환원할 수 있다.

\medskip\noindent
$f$가 기하적으로 연결된 올들을 갖고
$f_*\mathcal{O}_X = \mathcal{O}_S$라고 가정하자. 모든 $s \in S$에
대하여 에탈 근방 $(S', s') \to (S, s)$를 선택하여 바탕변환
$X' \to S'$가 $S$의 바탕변환으로서 단면을 갖게 할 수 있다.
사상 심화의 보조정리
\ref{more-morphisms-lemma-etale-nbhd-dominates-smooth}를 보라. 증명 첫머리의
언급에 의해 다음 문단에서 논할 경우로 환원된다.

\medskip\noindent
$f$가 기하적으로 연결된 올들을 갖고
$f_*\mathcal{O}_X = \mathcal{O}_S$이며, $s : S \to X$가 $f$의
단면이라고 가정하자. $S = \Spec(A)$가 아핀이라고
가정할 수 있으며 그렇게 가정한다. 사상
$s^* : R\Gamma(X, \mathcal{O}_X) \to R\Gamma(S, \mathcal{O}_S) = A$
은 사상 $A \to R\Gamma(X, \mathcal{O}_X)$의 분할이다. 따라서
$$
R\Gamma(X, \mathcal{O}_X) = A \oplus P
$$
라고 쓸 수 있다. 여기서 $P$는 $s^*$의 “핵”이다.
보조정리 \ref{derham-lemma-relative-bottom-part-degenerates}에 의해 $P$는
$D(A)$의 Tor 진폭이 $[1, n]$인 완전 대상이다. 보조정리
\ref{derham-lemma-relative-duality-hodge}의 증명에서와 같이
$H^n(X, \Omega^n_{X/S})$는 국소 자유 $A$-가군이며 계수는 $1$임을 알 수
있다. 실제로 이는 $A$의 쌍대이므로 계수 $1$인 자유 가군이다. 곧
생성원을 선택하겠지만 같은 생성원인지 확인하고 싶지 않으며, 그럴
필요도 없다.

\medskip\noindent
$Z \subset X$를 $s$의 상이라 하자. 이는 스킴의 보조정리
\ref{schemes-lemma-section-immersion}에 의해 $X$의 닫힌 부분스킴이다.
$Z \to X$는 인자의 보조정리
\ref{divisors-lemma-regular-quasi-regular-immersion}에 의해 특히 코쥘
정칙이고, 인자의 보조정리
\ref{divisors-lemma-section-smooth-regular-immersion}에 의해 정칙인
닫힌 몰입이다. 물론 $Z \to X$의 여차원은 $n$이다. 따라서 주의
\ref{derham-remark-how-to-use}에 의해 사상
$$
\gamma^{0, 0} : H^0(Z, \Omega^0_{Z/S}) \longrightarrow H^n(X, \Omega^n_{X/S})
$$
을 생각할 수 있고,
$\xi = \gamma^{0, 0}(1) \in H^n(X, \Omega^n_{X/S})$로 둔다.

\medskip\noindent
$\xi$가 기저 원소라고 주장한다. 실제로 최고차수에서 바탕변환이
성립하므로(예를 들어 스킴의 극한의 보조정리
\ref{limits-lemma-higher-direct-images-zero-above-dimension-fibre}를 보라)
$H^n(X, \Omega^n_{X/S}) \otimes_A k = H^n(X_k, \Omega^n_{X_k/k})$
를 임의의 환 사상 $A \to k$에 대해 얻는다. $\xi$의 구성이
바탕변환과 양립하므로(보조정리 \ref{derham-lemma-gysin-transverse-global}을
보라), $\xi$의 $H^n(X_k, \Omega^n_{X_k/k})$에서의 상은 $k$가 체이면
보조정리 \ref{derham-lemma-class-of-a-point}에 의해 영이 아니다. 따라서
$\xi$는 가역 가군의 어디서도 소멸하지 않는 단면이므로 생성원이다.

\medskip\noindent
$\theta \in H^n(X, \Omega^{n - 1}_{X/S})$라 하자.
$\text{d}(\theta)$가 $H^n(X, \Omega^n_{X/S})$에서 영임을 보여야 한다.
$\text{d}(\theta) = a \xi$라고 어떤 $a \in A$에 대해 쓸 수 있는데,
이는 $\xi$가 기저 원소이기 때문이다. 이제 $a = 0$임을 보이면
된다.

\medskip\noindent
닫힌 몰입
$$
\Delta : X \to X \times_S X
$$
을 생각하자. 이 사상은 또한 매끄러운 사상, 즉 어느 한 사영의
단면이므로 역시 여차원 $n$인 정칙 코쥘 몰입이다. 따라서 주의
\ref{derham-remark-how-to-use}의 사상들
$$
\gamma^{p, q} :
H^q(X, \Omega^p_{X/S})
\longrightarrow
H^{q + n}(X \times_S X, \Omega^{p + n}_{X \times_S X/S})
$$
을 생각할 수 있다. 다음 상을 생각하자.
$$
\gamma^{n - 1, n}(\theta) \in
H^{2n}(X \times_S X, \Omega^{2n - 1}_{X \times_S X})
$$
보조정리 \ref{derham-lemma-de-rham-complex-product}에 의해
$$
\Omega^{2n - 1}_{X \times_S X} =
\Omega^{n - 1}_{X/S} \boxtimes \Omega^n_{X/S} \oplus
\Omega^n_{X/S} \boxtimes \Omega^{n - 1}_{X/S}
$$
이다. 퀴네트 공식, 즉 스킴의 유도 범주의 보조정리
\ref{perfect-lemma-kunneth} 또는 스킴의 유도 범주의
보조정리 \ref{perfect-lemma-kunneth-single-sheaf}에 의해
$$
H^{2n}(X \times_S X, \Omega^{n - 1}_{X/S} \boxtimes \Omega^n_{X/S}) =
H^n(X, \Omega^{n - 1}_{X/S}) \otimes_A H^n(X, \Omega^n_{X/S})
$$
이고
$$
H^{2n}(X \times_S X, \Omega^n_{X/S} \boxtimes \Omega^{n - 1}_{X/S}) =
H^n(X, \Omega^n_{X/S}) \otimes_A H^n(X, \Omega^{n - 1}_{X/S})
$$
이다. 실제로 최고차수를 보고 있으므로 개입하는 더 높은 Tor 군이 없다.
$\xi$가 생성원이라는 사실과 결합하면
$$
\gamma^{n - 1, n}(\theta) = \theta_1 \otimes \xi + \xi \otimes \theta_2
$$
라고 쓸 수 있다. 여기서
$\theta_1, \theta_2 \in H^n(X, \Omega^{n - 1}_{X/S})$이다.
정확히 같은 방식으로 논하면
$$
\gamma^{n, n}(\xi) = b \xi \otimes \xi
$$
라고 쓸 수 있다. 이는
$H^{2n}(X \times_S X, \Omega^{2n}_{X \times_S X/S}) =
H^n(X, \Omega^n_{X/S}) \otimes_A H^n(X, \Omega^n_{X/S})$
안에서 성립하며, 여기서 어떤 $b \in H^0(S, \mathcal{O}_S)$가
존재한다.

\medskip\noindent
{\bf 주장:} $\theta_1 = \theta$, $\theta_2 = \theta$, 그리고 $b = 1$이다.
이 주장이 원하는 결과 $a = 0$을 함의함을 보이자. 실제로 보조정리
\ref{derham-lemma-gysin-differential-hodge}에 의해
$$
\gamma^{n, n}(\text{d}(\theta)) = \text{d}(\gamma^{n - 1, n}(\theta))
$$
이다. 위의 선택들에 의해
$$
a \xi \otimes \xi =
\gamma^{n, n}(a\xi) =
\text{d}(\theta \otimes \xi + \xi \otimes \theta) =
a \xi \otimes \xi + (-1)^n a \xi \otimes \xi
$$
이다. 맨 오른쪽 등식은 사상
$\text{d} : \Omega^{2n - 1}_{X \otimes_S X/S} \to \Omega^{2n}_{X \times_S X/S}$
이 보조정리 \ref{derham-lemma-de-rham-complex-product}에 의해 미분
$\text{d} \boxtimes 1 : \Omega^{n - 1}_{X/S} \boxtimes \Omega^n_{X/S}
\to \Omega^n_{X/S} \boxtimes \Omega^n_{X/S}$
과 미분
$(-1)^n 1 \boxtimes \text{d} : \Omega^n_{X/S} \boxtimes \Omega^{n - 1}_{X/S}
\to \Omega^n_{X/S} \boxtimes \Omega^n_{X/S}$의 합이라는 사실에서
나온다. 자세한 내용은 절 \ref{derham-section-kunneth}과 스킴의 유도 범주의 절 \ref{perfect-section-kunneth-complexes}의 논의를 보라.
$\xi \otimes \xi$는 계수 $1$인 자유 $A$-가군
$H^n(X, \Omega^n_{X/S}) \otimes_A H^n(X, \Omega^n_{X/S})$
의 기저이므로
$$
a = a + (-1)^n a \Rightarrow a = 0
$$
를 얻으며, 이것이 원하는 결론이다.

\medskip\noindent
증명의 나머지에서는 위 주장을 증명한다. 다음과 같이 쓰자.
$\eta = \gamma^{0, 0}(1) \in H^n(X \times_S X, \Omega^n_{X \times_S X/S})$.
다음 등식이 성립하므로
$\Omega^n_{X \times_S X/S} =
\bigoplus_{p + p' = n} \Omega^p_{X/S} \boxtimes \Omega^{p'}_{X/S}$
$$
\eta = \eta_0 + \eta_1 + \ldots + \eta_n
$$
라고 쓸 수 있다. 여기서 $\eta_p$는
$H^n(X \times_S X, \Omega^p_{X/S} \boxtimes \Omega^{n - p}_{X/S})$.
의 원소이다. $p = 0$일 때
\begin{align*}
H^n(X \times_S X, \mathcal{O}_X \boxtimes \Omega^n_{X/S})
& =
H^n(R\Gamma(X, \mathcal{O}_X) \otimes_A^\mathbf{L}
R\Gamma(X, \Omega^n_{X/S})) \\
& =
A \otimes_A H^n(X, \Omega^n_{X/S}) \oplus
H^n(P \otimes_A^\mathbf{L} R\Gamma(X, \Omega^n_{X/S}))
\end{align*}
라고 쓸 수 있다. 여기서는 앞서 주어진 분해
$R\Gamma(X, \mathcal{O}_X) = A \oplus P$를 사용했다.
사상 $(s, \text{id}) : X \to X \times_S X$를 생각하자. 그러면
스킴론적으로 $(s, \text{id})^{-1}(\Delta) = Z$이다. 따라서 보조정리
\ref{derham-lemma-gysin-transverse-global}에 의해
$(s, \text{id})^*\eta = \xi$임을 알 수 있다. 이는
$$
\xi = (s, \text{id})^*\eta = (s^* \otimes \text{id})(\eta_0)
$$
를 뜻한다. 정확히 말해 위 직합 분해에서 $\eta_0$의 첫 번째 성분은
$\xi$이다. 다시 말해
$$
\eta_0 = 1 \otimes \xi + \eta'_0
$$
라고 쓸 수 있으며, 여기서
$\eta'_0 \in H^n(P \otimes_A^\mathbf{L} R\Gamma(X, \Omega^n_{X/S}))$이다.
정확히 같은 방식으로 $p = n$일 때
\begin{align*}
H^n(X \times_S X, \Omega^n_{X/S} \boxtimes \mathcal{O}_X)
& =
H^n(R\Gamma(X, \Omega^n_{X/S}) \otimes_A^\mathbf{L}
R\Gamma(X, \mathcal{O}_X)) \\
& =
H^n(X, \Omega^n_{X/S}) \otimes_A A \oplus
H^n(R\Gamma(X, \Omega^n_{X/S}) \otimes_A^\mathbf{L} P)
\end{align*}
라고 쓸 수 있고, 다시
$$
\eta_n = \xi \otimes 1 + \eta'_n
$$
라고 쓸 수 있다. 여기서
$\eta'_n \in H^n(R\Gamma(X, \Omega^n_{X/S}) \otimes_A^\mathbf{L} P)$이다.

\medskip\noindent
$\text{pr}_1^*\theta = \theta \otimes 1$과
$\text{pr}_2^*\theta = 1 \otimes \theta$는 $X \times_S X$ 위 호지
코호몰로지류이고, 당기면 $\theta$가 되며 그 당김은 $\Delta$에
의함을 관찰하자.
따라서 보조정리 \ref{derham-lemma-gysin-projection-global}에 의해
$$
\theta_1 \otimes \xi + \xi \otimes \theta_2 =
\gamma^{n - 1, n}(\theta) =
(\theta \otimes 1) \cup \eta =
(1 \otimes \theta) \cup \eta
$$
가 $X \times_S X$의 $S$ 위 호지 코호몰로지 환에서 성립한다.
$X \times_S X/S$의 미분 가군 위 직합 분해로 표현하면
$$
\theta_1 \otimes \xi =
(\theta \otimes 1) \cup \eta_0
\quad\text{이고}\quad
\xi \otimes \theta_2 =
(1 \otimes \theta) \cup \eta_n
$$
를 얻는다. 위에서 얻은 공식
$\eta_0 = 1 \otimes \xi + \eta'_0$을 살펴보면
$\theta_1 = \theta$를 보이기 위해
$$
(\theta \otimes 1) \cup \eta'_0 = 0
$$
를 증명하면 충분하다. 이를 위해 $\theta \otimes 1$과의 컵곱은 다음
사상이 코호몰로지에 작용하여 주어짐을 관찰하자.
$$
(P \otimes_A^\mathbf{L} R\Gamma(X, \Omega^n_{X/S}))[-n]
\xrightarrow{\theta \otimes 1}
R\Gamma(X, \Omega^{n - 1}_{X/S}) \otimes_A^\mathbf{L}
R\Gamma(X, \Omega^n_{X/S})
$$
이는 유도 범주 안의 사상이며, 층 코호몰로지의 주의
\ref{cohomology-remark-cup-with-element-map-total-cohomology}를 보라.
이 사상은 다음 두 사상의 유도 텐서곱이다.
$$
\theta : P[-n] \to R\Gamma(X, \Omega^{n - 1}_{X/S})
\quad\text{이고}\quad
1 : R\Gamma(X, \Omega^n_{X/S}) \to R\Gamma(X, \Omega^n_{X/S})
$$
스킴의 유도 범주의 주의
\ref{perfect-remark-annoying-compatibility}를 보라. 그러나 그중 첫 번째는
$D(A)$에서 영이다. 그 이유는 Tor 진폭이 $[n + 1, 2n]$인 완전
복합체에서 코호몰로지가 차수 $0, 1, \ldots, n$에만 놓이는 복합체로
가는 사상이기 때문이다. 대수학 심화의 보조정리
\ref{more-algebra-lemma-splitting-unique}를 보라. 비슷한 논증으로
$(1 \otimes \theta) \cup \eta'_n$의 소멸도 보일 수 있다. 마지막으로
정확히 같은 방식으로
$$
b \xi \otimes \xi = \gamma^{n, n}(\xi) = (\xi \otimes 1) \cup \eta_0
$$
를 얻으며, 위와 같은 방식으로
$(\xi \otimes 1) \cup \eta'_0 = 0$을 보여 앞에서와 같이 결론을 얻는다.
이로써 증명이 끝난다.
\end{proof}

\begin{proposition}
\label{derham-proposition-relative-poincare-duality}
$S$를 준콤팩트 준분리 스킴이라 하자. $f : X \to S$를 모든 올이
공집합이 아니며 차원 $n$인 등차원 스킴이 되는 고유 매끄러운 스킴
사상이라 하자. $\mathcal{O}_S$-가군 사상
$$
t : R^{2n}f_*\Omega^\bullet_{X/S} \longrightarrow \mathcal{O}_S
$$
이 존재한다. 이는 $H^0(X, \mathcal{O}_X)$의 가역원을 곱하는 사상을
앞에 합성하는 것까지 유일하며, 다음 성질을 갖는다. 쌍
$$
Rf_*\Omega^\bullet_{X/S}
\otimes_{\mathcal{O}_S}^\mathbf{L}
Rf_*\Omega^\bullet_{X/S}[2n]
\longrightarrow
\mathcal{O}_S, \quad
(\xi, \xi') \longmapsto t(\xi \cup \xi')
$$
은 $S$ 위 완전 복합체들의 완전한 쌍이다.
\end{proposition}

\begin{proof}
증명은 명제 \ref{derham-proposition-poincare-duality}의 증명과 정확히 같다.

\medskip\noindent
상대 호지-드람 스펙트럼 열
$$
E_1^{p, q} = R^qf_*\Omega^p_{X/S} \Rightarrow R^{p + q}f_*\Omega^\bullet_{X/S}
$$
(절 \ref{derham-section-hodge-to-de-rham}), $\Omega^i_{X/S}$가 $i > n$일 때
소멸한다는 사실, 예를 들어 스킴의 극한의 보조정리
\ref{limits-lemma-higher-direct-images-zero-above-dimension-fibre}에
나오는 소멸, 그리고 보조정리들
\ref{derham-lemma-relative-bottom-part-degenerates}와
\ref{derham-lemma-relative-top-part-degenerates}의 결과들에 의해
$R^0f_*\Omega_{X/S} = R^0f_*\mathcal{O}_X$이고
$R^nf_*\Omega^n_{X/S} = R^{2n}f_*\Omega^\bullet_{X/S}$임을 알 수 있다.
더 정확히 말해 이 동일시들은 복합체들의 사상
$$
\Omega^\bullet_{X/S} \to \mathcal{O}_X[0]
\quad\text{이고}\quad
\Omega^n_{X/S}[-n] \to \Omega^\bullet_{X/S}
$$
에서 나온다. $t : R^{2n}f_*\Omega_{X/S} \to \mathcal{O}_S$를 선택하여,
이 동일시 아래에서 보조정리 \ref{derham-lemma-relative-duality-hodge}에서와
같은 $t$에 대응하게 하자.

\medskip\noindent
$\Omega^\bullet = \Omega^\bullet_{X/S}$로 줄여 쓰자. 우리의 상황에서는
사상 (\ref{derham-equation-wedge})이 다음과 같이 된다.
$$
\wedge :
\text{Tot}(\Omega^\bullet \otimes_{f^{-1}\mathcal{O}_S} \Omega^\bullet)
\longrightarrow
\Omega^\bullet
$$
모든 정수 $p = 0, 1, \ldots, n$에 대해 이 사상은 부분복합체
$\text{Tot}(\sigma_{> p} \Omega^\bullet \otimes_{f^{-1}\mathcal{O}_S}
\sigma_{\geq n - p} \Omega^\bullet)$를 차수상의 이유로 영으로 보낸다.
따라서 $\wedge$를 부분복합체
$\text{Tot}(\Omega^\bullet \otimes_{f^{-1}\mathcal{O}_S}
\geq_{n - p}\Omega^\bullet)$에 제한하면 복합체들의 사상
$$
\gamma_p :
\text{Tot}(\sigma_{\leq p} \Omega^\bullet \otimes_{f^{-1}\mathcal{O}_S}
\sigma_{\geq n - p} \Omega^\bullet)
\longrightarrow
\Omega^\bullet
$$
을 거쳐 인수분해된다. 절 \ref{derham-section-cup-product}에서와 같은 절차를
사용하면 상대 컵곱들
$$
Rf_*\sigma_{\leq p} \Omega^\bullet
\otimes_{\mathcal{O}_S}^\mathbf{L}
Rf_*\sigma_{\geq n - p}\Omega^\bullet
\longrightarrow
Rf_*\Omega^\bullet
$$
을 얻는다. $p$에 대한 귀납법으로, 이 컵곱들이 $t$를 거쳐
$Rf_*\sigma_{\leq p} \Omega^\bullet$와
$Rf_*\sigma_{\geq n - p}\Omega^\bullet[2n]$ 사이의 완전한 쌍을
유도함을 증명하겠다. $p = n$일 때 이것이
명제의 주장이다.

\medskip\noindent
기초 단계는 $p = 0$이다. 이 경우
$$
Rf_*\sigma_{\leq p}\Omega^\bullet = Rf_*\mathcal{O}_X
\quad\text{이고}\quad
Rf_*\sigma_{\geq n - p}\Omega^\bullet[2n] = Rf_*(\Omega^n[-n])[2n] =
Rf_*\Omega^n[n]
$$
이다. 이 경우 보조정리 \ref{derham-lemma-relative-duality-hodge}의
$Rf_*\mathcal{O}_X$와 $Rf_*\Omega^n[n]$ 사이의 쌍을 그대로 얻으므로
결론이 성립한다.

\medskip\noindent
귀납 단계. 결과가 $p$에 대해 성립함을 안다고 하자. 그러면 특별 삼각형
$$
\Omega^{p + 1}[-p - 1] \to
\sigma_{\leq p + 1}\Omega^\bullet \to
\sigma_{\leq p}\Omega^\bullet \to
\Omega^{p + 1}[-p]
$$
과 특별 삼각형
$$
\sigma_{\geq n - p}\Omega^\bullet \to
\sigma_{\geq n - p - 1}\Omega^\bullet \to
\Omega^{n - p - 1}[-n + p + 1] \to
(\sigma_{\geq n - p}\Omega^\bullet)[1]
$$
을 생각하자. 둘 다 $f^{-1}\mathcal{O}_S$-가군 층의 복합체들의 호모토피
범주에서 특별 삼각형임을 관찰하자. 특히 사상
$\sigma_{\leq p}\Omega^\bullet \to \Omega^{p + 1}[-p]$와
$\Omega^{n - p - 1}[-d + p + 1] \to
(\sigma_{\geq n - p}\Omega^\bullet)[1]$은 실제 복합체 사상들로 주어진다.
구체적으로 미분 $\Omega^p \to \Omega^{p + 1}$과 미분
$\Omega^{n - p - 1} \to \Omega^{n - p}$을 사용한다. 이 특별 삼각형들에
$Rf_*$를 적용하여 연관되어 얻은 특별 삼각형들을 생각하자.
$$
\xymatrix{
Rf_*\sigma_{\leq p}\Omega^\bullet \ar[d]_a \\
Rf_*\Omega^{p + 1}[-p - 1] \ar[d]_b \\
Rf_*\sigma_{\leq p + 1}\Omega^\bullet \ar[d]_c \\
Rf_*\sigma_{\leq p}\Omega^\bullet \ar[d]_d \\
Rf_*\Omega^{p + 1}[-p - 1]
}
\quad\quad
\xymatrix{
Rf_*\sigma_{\geq n - p}\Omega^\bullet \\
Rf_*\Omega^{n - p - 1}[-n + p + 1] \ar[u]_{a'} \\
Rf_*\sigma_{\geq n - p - 1}\Omega^\bullet \ar[u]_{b'} \\
Rf_*\sigma_{\geq n - p}\Omega^\bullet \ar[u]_{c'} \\
Rf_*\Omega^{n - p - 1}[-n + p + 1] \ar[u]_{d'}
}
$$
아래에서 쌍들 $(a, a')$, $(b, b')$, $(c, c')$, 그리고 $(d, d')$가
주어진 쌍들과 양립함을 보이겠다. 이는 왼쪽 특별 삼각형에서 오른쪽
특별 삼각형에 $R\SheafHom(-, \mathcal{O}_S)$를 적용하여 얻은 특별
삼각형으로 가는 사상을 얻는다는 뜻이다. 귀납 가정과 보조정리
\ref{derham-lemma-duality-hodge}에 의해 첫째, 둘째, 넷째, 다섯째 행의 복합체들
사이에 위에서 구성한 쌍들은 완전하다. 즉 쌍대를 취하면 동형을 정한다.
유도 범주의 보조정리
\ref{derived-lemma-third-isomorphism-triangle}에 의해 가운데 행의 복합체들
사이의 쌍도 원하는 대로 완전하다는 결론을 얻는다.

\medskip\noindent
$e : K \to K'$와 $e' : M' \to M$를 $D(\mathcal{O}_S)$의 대상들 사이의
사상이라 하고,
$K \otimes_{\mathcal{O}_S}^\mathbf{L} M \to \mathcal{O}_S$와
$K' \otimes_{\mathcal{O}_S}^\mathbf{L} M' \to \mathcal{O}_S$
를 쌍들이라고 하자. 다음 그림이 가환이면 이 쌍들이 양립한다고 한다.
$$
\xymatrix{
K' \otimes_{\mathcal{O}_S}^\mathbf{L} M' \ar[d] &
K \otimes_{\mathcal{O}_S}^\mathbf{L} M'
\ar[l]^{e \otimes 1} \ar[d]^{1 \otimes e'} \\
\mathcal{O}_S &
K \otimes_{\mathcal{O}_S}^\mathbf{L} M \ar[l]
}
$$
실제로 이는 그림
$$
\xymatrix{
K \ar[r] \ar[d]_e & R\SheafHom(M, \mathcal{O}_S) \ar[d]^{R\SheafHom(e', -)} \\
K' \ar[r] & R\SheafHom(M', \mathcal{O}_S)
}
$$
이 가환함을 뜻하므로 우리의 목적에 충분하다.

\medskip\noindent
쌍 $(c, c')$에 대해 이를 증명하자. 여기서는 다음 가환 그림이 있음을
관찰하기만 하면 된다.
$$
\xymatrix{
\text{Tot}(\sigma_{\leq p} \Omega^\bullet \otimes_{f^{-1}\mathcal{O}_S}
\sigma_{\geq n - p} \Omega^\bullet) \ar[d]_{\gamma_p} &
\text{Tot}(\sigma_{\leq p + 1} \Omega^\bullet \otimes_{f^{-1}\mathcal{O}_S}
\sigma_{\geq n - p} \Omega^\bullet) \ar[l] \ar[d] \\
\Omega^\bullet &
\text{Tot}(\sigma_{\leq p + 1} \Omega^\bullet \otimes_{f^{-1}\mathcal{O}_S}
\sigma_{\geq n - p - 1} \Omega^\bullet) \ar[l]_-{\gamma_{p + 1}}
}
$$
컵곱의 함자성에 의해 원하는 그림의 가환성을 얻는다.

\medskip\noindent
마찬가지로 쌍 $(b, b')$에 대해서는 다음 가환 그림을 사용한다.
$$
\xymatrix{
\text{Tot}(\sigma_{\leq p + 1} \Omega^\bullet \otimes_{f^{-1}\mathcal{O}_S}
\sigma_{\geq n - p - 1} \Omega^\bullet) \ar[d]_{\gamma_{p + 1}} &
\text{Tot}(\Omega^{p + 1}[-p - 1] \otimes_{f^{-1}\mathcal{O}_S}
\sigma_{\geq n - p - 1} \Omega^\bullet) \ar[l] \ar[d] \\
\Omega^\bullet &
\Omega^{p + 1}[-p - 1]
\otimes_{f^{-1}\mathcal{O}_S}
\Omega^{n - p - 1}[-n + p + 1] \ar[l]_-\wedge
}
$$

\medskip\noindent
쌍들 $(d, d')$와 $(a, a')$에 대해서는 다음 가환 그림을 사용한다.
$$
\xymatrix{
\Omega^{p + 1}[-p] \otimes_{f^{-1}\mathcal{O}_S}
\Omega^{n - p - 1}[-n + p] \ar[d] &
\text{Tot}(\sigma_{\leq p}\Omega^\bullet \otimes_{f^{-1}\mathcal{O}_S}
\Omega^{n - p - 1}[-n + p]) \ar[l] \ar[d] \\
\Omega^\bullet &
\text{Tot}(\sigma_{\leq p}\Omega^\bullet \otimes_{f^{-1}\mathcal{O}_S}
\sigma_{\geq n - p}\Omega^\bullet) \ar[l]
}
$$

\medskip\noindent
$t$가 $H^0(X, \mathcal{O}_X)$의 가역원을 곱하는 사상을 앞에 합성하는
것까지 유일함을 보이는 논증은 생략한다.
\end{proof}
